% =============================================================================
\section{Condiciones de frontera}\label{sec:frontera}
% =============================================================================

Las ecuaciones $\divg\vD=\rhof$ y $\rot\vE=\vcero$ valen en el interior de cada medio, pero en una interfase entre
dos dieléctricos, o entre un dieléctrico y el vacío o un conductor, los campos pueden ser discontinuos y
las derivadas dejan de tener sentido. Sus formas integrales, en cambio, siguen valiendo, y de ellas se
deducen las relaciones entre los campos a uno y otro lado. Repasamos la deducción del Tema~1~\cite{T1},
ahora con $\vD$.

Sea $\Sigma$ una superficie que separa los medios~1 y~2, con normal unitaria $\nn$ dirigida de~1 a~2 y
posiblemente con una densidad de carga libre $\sigmaf$. Para un campo $\vec F$, escribimos $F_n=\vec F\cdot\nn$ y
$\vec F_t=\vec F-F_n\nn$.

% Figura: deducción de las condiciones de frontera (cilindro de Gauss y circuito
% rectangular) e interpretación del salto de la componente normal de P.
\begin{figure}[t]
\centering
\begin{subfigure}[b]{0.33\textwidth}\centering
\begin{tikzpicture}[scale=0.95]
  \fill[dielec] (-1.8,-1.4) rectangle (1.8,0);
  \fill[dielecB] (-1.8,0) rectangle (1.8,1.4);
  \draw[texto, line width=0.6pt] (-1.8,0) -- (1.8,0);
  \node[etiq, anchor=west] at (-1.75,1.12) {$2$};
  \node[etiq, anchor=west] at (-1.75,-1.12) {$1$};
  % cilindro
  \fill[fondo!50!dielecB] (-0.7,-0.42) rectangle (0.7,0.42);
  \draw[gauss] (-0.7,-0.42) -- (-0.7,0.42) (0.7,-0.42) -- (0.7,0.42);
  \draw[gauss] (0,0.42) ellipse (0.7 and 0.16);
  \draw[gauss] (-0.7,-0.42) arc (180:360:0.7 and 0.16);
  \draw[gauss, densely dotted] (0.7,-0.42) arc (0:180:0.7 and 0.16);
  \draw[vecn] (0,0.42) -- (0,1.2) node[right, etiq] {$\nn$};
  \draw[vecn] (0,-0.58) -- (0,-1.3) node[right, etiq] {$-\nn$};
  \node[etiqt] at (-1.12,0.62) {$\Delta S$};
  \draw[cota] (1.0,-0.42) -- (1.0,0.42) node[midway, right, etiq] {$h\!\to\!0$};
\end{tikzpicture}
\caption{Normal de $\vD$}
\end{subfigure}\hfill
\begin{subfigure}[b]{0.33\textwidth}\centering
\begin{tikzpicture}[scale=0.95]
  \fill[dielec] (-1.8,-1.4) rectangle (1.8,0);
  \fill[dielecB] (-1.8,0) rectangle (1.8,1.4);
  \draw[texto, line width=0.6pt] (-1.8,0) -- (1.8,0);
  \node[etiq, anchor=west] at (-1.75,1.12) {$2$};
  \node[etiq, anchor=west] at (-1.75,-1.12) {$1$};
  \draw[acentoB, line width=0.7pt] (-0.8,-0.35) rectangle (0.8,0.35);
  \draw[acentoB, line width=0.7pt, flechamedio] (-0.8,-0.35) -- (0.8,-0.35);
  \draw[acentoB, line width=0.7pt, flechamedio] (0.8,0.35) -- (-0.8,0.35);
  \node[etiqt] at (0,0.62) {$\Delta\ell$};
  \draw[vecn] (-1.3,0) -- (-1.3,0.75) node[right, etiq] {$\nn$};
  \draw[vecn] (-1.3,0) -- (-0.55,0) node[pos=0.55, below=1pt, etiq] {$\tu$};
  \draw[cota] (1.1,-0.35) -- (1.1,0.35) node[midway, right, etiq] {$h\!\to\!0$};
\end{tikzpicture}
\caption{Tangencial de $\vE$}
\end{subfigure}\hfill
\begin{subfigure}[b]{0.32\textwidth}\centering
\begin{tikzpicture}[scale=0.95]
  \fill[dielec] (-1.7,-1.4) rectangle (1.7,0);
  \fill[dielecB] (-1.7,0) rectangle (1.7,1.4);
  \draw[texto, line width=0.6pt] (-1.7,0) -- (1.7,0);
  \node[etiq, anchor=west] at (-1.65,1.12) {$2$};
  \node[etiq, anchor=west] at (-1.65,-1.12) {$1$};
  \draw[vecP] (-0.45,-1.15) -- (-0.45,-0.05) node[pos=0.35, left, etiq] {$\vP_1$};
  \draw[vecP] (0.45,0.05) -- (0.45,0.85) node[pos=0.6, right, etiq] {$\vP_2$};
  \draw[vecn] (1.25,0) -- (1.25,0.65) node[right, etiq] {$\nn$};
  \node[cpos] at (-0.45,0.17) {$+P_{1n}$};
  \node[cneg] at (0.45,-0.2) {$-P_{2n}$};
\end{tikzpicture}
\caption{$\sigmaP=P_{1n}-P_{2n}$}
\end{subfigure}
\caption{Condiciones de frontera entre dos medios; $\nn$ es la normal unitaria del medio 1 al 2.
(a) Cilindro de Gauss de bases $\Delta S$ y altura $h\to0$: el flujo lateral desaparece y queda
$(D_{2n}-D_{1n})\Delta S=\sigmaf\,\Delta S$. (b) Circuito rectangular de lados $\Delta\ell$ y $h\to0$:
la circulación de $\vE$ da $(E_{2t}-E_{1t})\Delta\ell=0$. (c) Cada medio aporta a la interfase la carga
ligada $\vP\cdot\nn_{\mathrm{ext}}$ de su propia cara: $+P_{1n}$ el medio 1 (su normal exterior es $\nn$)
y $-P_{2n}$ el medio 2 (su normal exterior es $-\nn$).}
\label{fig:frontera}
\end{figure}


\subsection{Las dos condiciones fundamentales}\label{sec:condiciones}

\begin{teorema}[Condiciones de frontera]\label{teo:frontera}
En cada punto de $\Sigma$ en que los campos son acotados a ambos lados:
\begin{enumerate}
  \item la componente normal de $\vD$ salta en la densidad de carga libre,
        \begin{equation}\label{eq:frontera-D}
          D_{2n}-D_{1n}=\sigmaf
          \qquad\Longleftrightarrow\qquad
          \nn\cdot(\vD_2-\vD_1)=\sigmaf ;
        \end{equation}
  \item la componente tangencial de $\vE$ es continua,
        \begin{equation}\label{eq:frontera-E}
          \vE_{2t}=\vE_{1t}
          \qquad\Longleftrightarrow\qquad
          \nn\times(\vE_2-\vE_1)=\vcero ;
        \end{equation}
  \item el potencial es continuo, $\phi_1=\phi_2$ sobre $\Sigma$.
\end{enumerate}
\end{teorema}

\begin{proof}
(i) Aplicamos la forma integral de la ley de Gauss para $\vD$ a un cilindro de bases $\Delta S$, paralelas a $\Sigma$ a uno
y otro lado, y altura $h$ (\cref{fig:frontera}\,(a)); las bases son tan pequeñas que $\vD$ es prácticamente uniforme en
cada una. Al hacer $h\to0$, el área lateral tiende a cero y, como $\vD$ es acotado, su flujo desaparece. Solo quedan las
bases, con normales exteriores $\nn$ (en el medio~2) y $-\nn$ (en el medio~1):
\[
  \oiint\vD\cdot\dS\ \xrightarrow{\;h\to0\;}\ \vD_2\cdot\nn\,\Delta S+\vD_1\cdot(-\nn)\,\Delta S=(D_{2n}-D_{1n})\,\Delta S .
\]
En cuanto a la carga libre encerrada, la volúmica es $\rhof\,\Delta S\,h\to0$ y solo sobrevive la superficial, $\sigmaf\,\Delta S$.
Igualando y dividiendo por $\Delta S$ se obtiene \eqref{eq:frontera-D}.

(ii) Como el campo electrostático es conservativo, $\oint_C\vE\cdot\dl=0$ para cualquier curva cerrada. Tomamos un rectángulo
de lados $\Delta\ell$ paralelos a $\Sigma$, según una dirección tangente $\tu$, y lados $h$ perpendiculares (\cref{fig:frontera}\,(b)).
Al hacer $h\to0$ los lados perpendiculares no contribuyen y
\[
  \vE_2\cdot(-\tu)\,\Delta\ell+\vE_1\cdot\tu\,\Delta\ell=0
  \quad\Longrightarrow\quad (E_{2t}-E_{1t})\,\Delta\ell=0 ,
\]
para cualquier dirección tangente $\tu$, lo que equivale a \eqref{eq:frontera-E}.

(iii) Sean $P_1$ y $P_2$ dos puntos a uno y otro lado de $\Sigma$, separados una distancia $h$ según la normal. Entonces
$\phi(P_2)-\phi(P_1)=-\int_{P_1}^{P_2}\vE\cdot\dl$, que en módulo no supera $h\,\sup\abs{\vE}$ y tiende a cero con $h$ porque $\vE$ es acotado.
\end{proof}

\begin{observacion}[Continuidad de $\phi$ y de $\vE_t$ son equivalentes]\label{obs:phi-Et}
Si $\phi$ es continuo, $\phi_1=\phi_2$ en todos los puntos de $\Sigma$: ambos potenciales coinciden como funciones sobre la
superficie, y también sus derivadas a lo largo de ella, $E_{1t}=-\partial\phi_1/\partial t=-\partial\phi_2/\partial t=E_{2t}$. Recíprocamente, si
$\vE_t$ es continuo, para dos puntos $A$ y $B$ de la frontera, recorriéndola justo por encima y justo por debajo,
$\phi_2(B)-\phi_2(A)=-\int_A^BE_{2t}\,\dd\ell=-\int_A^BE_{1t}\,\dd\ell=\phi_1(B)-\phi_1(A)$, de modo que $\phi_2-\phi_1$ es constante a lo largo
de $\Sigma$; el argumento de (iii) muestra que esa constante es nula. Las dos condiciones son la misma
información; la (iii) es la que se usa al resolver la ecuación de Laplace.
\end{observacion}

\subsection{Las otras dos componentes}\label{sec:otras-componentes}

Las condiciones \eqref{eq:frontera-D} y \eqref{eq:frontera-E} no dicen nada directamente de $E_n$ ni de $\vD_t$. Se obtienen
de ellas despejando $\vE=(\vD-\vP)/\eo$.

\begin{corolario}[Saltos de $E_n$ y de $\vD_t$]\label{cor:otras}
En las mismas condiciones del \cref{teo:frontera}:
\begin{enumerate}[start=4]
  \item la componente normal de $\vE$ salta en la carga total, libre y ligada, dividida por $\eo$:
        \begin{equation}\label{eq:salto-En}
          E_{2n}-E_{1n}=\frac{1}{\eo}\big(\sigmaf+\sigmaP\big),
          \qquad \sigmaP=P_{1n}-P_{2n} ;
        \end{equation}
  \item la componente tangencial de $\vD$ salta como la de $\vP$:
        \begin{equation}\label{eq:salto-Dt}
          \vD_{2t}-\vD_{1t}=\vP_{2t}-\vP_{1t} .
        \end{equation}
\end{enumerate}
\end{corolario}

\begin{proof}
(iv) Sustituyendo $D_n=\eo E_n+P_n$ en \eqref{eq:frontera-D},
$(\eo E_{2n}+P_{2n})-(\eo E_{1n}+P_{1n})=\sigmaf$, es decir, $\eo(E_{2n}-E_{1n})=\sigmaf+P_{1n}-P_{2n}$. Que $P_{1n}-P_{2n}$ es la carga ligada neta de la
interfase se sigue del \cref{teo:densidades}: cada medio aporta $\vP\cdot\nn_{\mathrm{ext}}$ en su propia cara, y la normal exterior del
medio~1 es $\nn$ mientras que la del medio~2 es $-\nn$ (\cref{fig:frontera}\,(c)).
(v) Tomando componentes tangenciales en $\vD=\eo\vE+\vP$ y usando \eqref{eq:frontera-E},
$\vD_{2t}-\vD_{1t}=\eo(\vE_{2t}-\vE_{1t})+(\vP_{2t}-\vP_{1t})=\vP_{2t}-\vP_{1t}$.
\end{proof}

\begin{observacion}[Lectura de los saltos]\label{obs:lectura-saltos}
La ecuación \eqref{eq:salto-En} es la condición del vacío con la carga total: el campo $\vE$ «ve» todas las cargas. La
\eqref{eq:salto-Dt} es la versión superficial de $\rot\vD=\rot\vP$: como $\vD$ no es conservativo, su componente tangencial no tiene
por qué conservarse. Dos casos particulares: en la cara de un dieléctrico en el vacío ($\vP_2=\vcero$) se recupera
$\sigmaP=\vP\cdot\nn$; y junto a un conductor en equilibrio ($\vE_1=\vD_1=\vcero$ en el metal), $D_{2n}=\sigmaf$ es la carga libre del conductor y
$\eo E_{2n}=\sigmaf+\sigmaP$ incluye la carga ligada del dieléctrico adyacente.
\end{observacion}

\subsection{Medios lineales: ley de refracción}\label{sec:refraccion}

Si ambos medios son \LHI{}, $\vD_1=\varepsilon_1\vE_1$ y $\vD_2=\varepsilon_2\vE_2$, y las condiciones fundamentales toman la forma
\begin{equation}\label{eq:frontera-LHI}
  \varepsilon_2E_{2n}-\varepsilon_1E_{1n}=\sigmaf,\qquad
  \frac{\vD_{2t}}{\varepsilon_2}-\frac{\vD_{1t}}{\varepsilon_1}=\vcero .
\end{equation}
Si $\varepsilon_1\neq\varepsilon_2$, ni $E_n$ ni $\vD_t$ son continuos, aunque no haya carga libre en la interfase. Esta discontinuidad del
campo se traduce en una refracción de sus líneas.

\begin{proposicion}[Ley de refracción de las líneas de campo]\label{prop:refraccion}
En una interfase sin carga libre entre dos medios \LHI{}, si $\theta_1$ y $\theta_2$ son los ángulos que forman $\vE_1$ y $\vE_2$ con la
normal,
\begin{equation}\label{eq:refraccion}
  \boxed{\;\frac{\tan\theta_2}{\varepsilon_2}=\frac{\tan\theta_1}{\varepsilon_1}\; .}
\end{equation}
Las líneas de $\vD$ se refractan con los mismos ángulos, pues $\vD\parallel\vE$ en cada medio.
\end{proposicion}

\begin{proof}
Con $\sigmaf=0$, $E_{2t}=E_{1t}$ y $\varepsilon_2E_{2n}=\varepsilon_1E_{1n}$. Como $\tan\theta_i=E_{it}/E_{in}$ (\cref{fig:refraccion}\,(a)),
\[
  \frac{\tan\theta_2}{\tan\theta_1}=\frac{E_{2t}}{E_{1t}}\,\frac{E_{1n}}{E_{2n}}=1\cdot\frac{\varepsilon_2}{\varepsilon_1}. \qedhere
\]
\end{proof}

Al pasar a un medio de mayor permitividad, $\theta_2>\theta_1$: la línea se aleja de la normal y tiende a «tumbarse» a lo largo
de la interfase (\cref{fig:refraccion}\,(b)). En el límite $\varepsilon_2\to\infty$, las líneas del medio~1 llegan perpendiculares a la
interfase, como a la superficie de un conductor~\cite{T3}.

% Figura: ley de refracción de las líneas de campo y ángulo refractado.
\begin{figure}[t]
\centering
\begin{subfigure}[b]{0.46\textwidth}\centering
\begin{tikzpicture}[scale=1.0]
  % medio 1 abajo (eps1), medio 2 arriba (eps2 > eps1); tan t2/eps2 = tan t1/eps1
  \fill[dielec] (-2.4,-1.9) rectangle (2.4,0);
  \fill[dielecB] (-2.4,0) rectangle (2.4,1.9);
  \draw[texto, line width=0.6pt] (-2.4,0) -- (2.4,0);
  \node[etiq, anchor=west] at (-2.35,1.6) {$\varepsilon_2$};
  \node[etiq, anchor=west] at (-2.35,-1.6) {$\varepsilon_1$};
  \draw[texto!50!fondo, dashed, line width=0.35pt] (0,-1.9) -- (0,1.9);
  % E1: ángulo 30 grados con la normal; E2: atan(2.5 tan 30) = 55.3 grados
  \coordinate (E1) at ({-1.55*sin(30)},{-1.55*cos(30)});
  \coordinate (E2) at ({1.75*sin(55.3)},{1.75*cos(55.3)});
  \draw[vecE] (E1) -- (0,0);
  \draw[vecE] (0,0) -- (E2);
  \node[etiq, anchor=north east] at ($(E1)+(0.02,-0.02)$) {$\vE_1$};
  \node[etiq, anchor=west] at ($(E2)+(0.05,0.05)$) {$\vE_2$};
  % componentes
  \draw[campoE!60!fondo, dashed, line width=0.4pt] (E2) -- ({1.75*sin(55.3)},0)
     node[pos=0.5, right, etiq] {$E_{2n}$};
  \draw[campoE!60!fondo, dashed, line width=0.4pt] (E2) -- (0,{1.75*cos(55.3)})
     node[pos=0.5, above, etiq] {$E_{2t}$};
  \draw[campoE!60!fondo, dashed, line width=0.4pt] (E1) -- (0,{-1.55*cos(30)})
     node[pos=0.5, below, etiq] {$E_{1t}$};
  \draw[campoE!60!fondo, dashed, line width=0.4pt] (E1) -- ({-1.55*sin(30)},0)
     node[pos=0.72, left, etiq] {$E_{1n}$};
  % ángulos
  \draw[tenue, line width=0.4pt] (0,0.55) arc (90:34.7:0.55);
  \node[etiqt] at (62:0.78) {$\theta_2$};
  \draw[tenue, line width=0.4pt] (0,-0.55) arc (-90:-120:0.55);
  \node[etiqt] at (-105:0.82) {$\theta_1$};
  \draw[vecn] (-1.75,0) -- (-1.75,0.7) node[right, etiq] {$\nn$};
\end{tikzpicture}
\caption{Refracción de $\vE$ ($\varepsilon_2>\varepsilon_1$, $\sigmaf=0$)}
\end{subfigure}\hfill
\begin{subfigure}[b]{0.5\textwidth}\centering
\begin{tikzpicture}
\begin{axis}[width=6.9cm, height=5.6cm, axis lines=left, xmin=0, xmax=90, ymin=0, ymax=92,
  xlabel={$\theta_1$ (grados)}, ylabel={$\theta_2$ (grados)}, xtick={0,30,60,90}, ytick={0,30,60,90},
  grid=major, clip=false]
  \foreach \k/\c in {0.25/campoE!45!fondo, 0.5/campoE!70!fondo, 2/campoD!70!fondo, 4/campoD} {
    \edef\temp{\noexpand\addplot[draw=\c, line width=0.9pt, domain=0:89.5, samples=120]
      {atan(\k*tan(x))};}\temp
  }
  \addplot[draw=tenue, line width=0.6pt, dashed, domain=0:90] {x};
  \node[etiqf, font=\scriptsize, anchor=west] at (axis cs:5,82) {$\varepsilon_2/\varepsilon_1=4$};
  \node[etiqf, font=\scriptsize] at (axis cs:38,57.4) {$2$};
  \node[etiqf, font=\scriptsize] at (axis cs:50,50) {$1$};
  \node[etiqf, font=\scriptsize] at (axis cs:62,43.2) {$\tfrac12$};
  \node[etiqf, font=\scriptsize] at (axis cs:76,44.7) {$\tfrac14$};
\end{axis}
\end{tikzpicture}
\caption{$\theta_2=\arctan\!\big[(\varepsilon_2/\varepsilon_1)\tan\theta_1\big]$}
\end{subfigure}
\caption{Ley de refracción de las líneas de campo en una interfase sin carga libre. (a) La
componente tangencial de $\vE$ se conserva y la normal se reduce en el factor $\varepsilon_1/\varepsilon_2$: al
pasar al medio de mayor permitividad la línea se aleja de la normal. (b) Ángulo refractado para
varios cocientes de permitividades; la diagonal ($\varepsilon_1=\varepsilon_2$) no tiene refracción.}
\label{fig:refraccion}
\end{figure}


\begin{ejemplo}[Esfera dieléctrica en un campo uniforme]\label{ej:esfera-dielectrica}
Una esfera de radio $a$ y permitividad $\varepsilon=\er\eo$ se coloca en un campo uniforme $\vE_0=E_0\,\uv z$. Busquemos una solución en
la que la esfera se polariza \emph{uniformemente}. Por el \cref{prob:esfera-polarizada}, una esfera con polarización uniforme
$\vP$ crea en su interior el campo uniforme $-\vP/3\eo$, y fuera, el de un dipolo $\vp=\tfrac43\pi a^3\vP$. El campo interior total es
$\vE_{\mathrm{int}}=\vE_0-\vP/3\eo$, y la relación constitutiva exige $\vP=\eo\chi\vE_{\mathrm{int}}$. Resolviendo este sistema autoconsistente,
\[
  \vE_{\mathrm{int}}=\vE_0-\frac{\chi}{3}\vE_{\mathrm{int}}
  \quad\Longrightarrow\quad
  \vE_{\mathrm{int}}=\frac{3}{\er+2}\,\vE_0 ,\qquad
  \vP=3\eo\,\frac{\er-1}{\er+2}\,\vE_0 .
\]
En coordenadas esféricas, el potencial correspondiente es
\[
  \phi_{\mathrm{int}}=-\frac{3}{\er+2}\,E_0\,r\cos\theta,
  \qquad
  \phi_{\mathrm{ext}}=-E_0\,r\cos\theta+\frac{\er-1}{\er+2}\,\frac{E_0a^3\cos\theta}{r^2}.
\]
\emph{Comprobación.} Ambos son armónicos ($r\cos\theta=z$ y $\cos\theta/r^2$ es el potencial de un dipolo); $\phi_{\mathrm{ext}}\to-E_0z$ en el
infinito; en $r=a$ los dos valen $-3E_0a\cos\theta/(\er+2)$, y las componentes radiales de $\vD$ coinciden:
\[
  -\varepsilon\,\frac{\partial\phi_{\mathrm{int}}}{\partial r}=\frac{3\er\eo E_0\cos\theta}{\er+2}
  =-\eo\,\frac{\partial\phi_{\mathrm{ext}}}{\partial r}\bigg|_{r=a}=\eo E_0\cos\theta\Big(1+2\,\frac{\er-1}{\er+2}\Big).
\]
Por el teorema de unicidad~\cite{T3,Jackson}, que se extiende sin cambios a medios \LHI{} a trozos, esta es la solución.
El campo interior es uniforme y más débil que el aplicado; las líneas de $\vD$, continuas, se concentran en la esfera
(\cref{fig:esfera-dielectrica}). En el límite $\er\to\infty$, $\vE_{\mathrm{int}}\to\vcero$ y el dipolo inducido tiende a $4\pi\eo a^3\vE_0$: se
recupera la esfera conductora del Tema~3~\cite{T3}.
\end{ejemplo}

% Figura (simulación): esfera dieléctrica eps_r = 4 en un campo uniforme.
% Datos: scripts/lineas_campo.py -> datos/esfera-dielectrica-{E,D}.tex
\begin{figure}[t]
\centering
\pgfplotsset{esfera/.style={marco, width=7.6cm, xmin=-3, xmax=3, ymin=-2.2, ymax=2.2}}
\begin{subfigure}[b]{0.49\textwidth}\centering
\begin{tikzpicture}
\begin{axis}[esfera]
  \fill[dielec] (axis cs:0,0) circle[radius=1];
  %% Generado por scripts/lineas_campo.py -- no editar a mano
\addplot[equipot] coordinates {(2.781,-2.200) (2.782,-2.165) (2.782,-2.130) (2.783,-2.094) (2.784,-2.059) (2.784,-2.024) (2.785,-1.989) (2.785,-1.954) (2.786,-1.918) (2.787,-1.883) (2.787,-1.848) (2.788,-1.813) (2.789,-1.778) (2.789,-1.742) (2.790,-1.707) (2.791,-1.672) (2.791,-1.637) (2.792,-1.602) (2.793,-1.566) (2.793,-1.531) (2.794,-1.496) (2.795,-1.461) (2.795,-1.426) (2.796,-1.390) (2.797,-1.355) (2.797,-1.320) (2.798,-1.285) (2.799,-1.250) (2.799,-1.214) (2.800,-1.179) (2.801,-1.144) (2.801,-1.109) (2.802,-1.074) (2.802,-1.038) (2.803,-1.003) (2.804,-0.968) (2.804,-0.933) (2.805,-0.898) (2.805,-0.862) (2.806,-0.827) (2.807,-0.792) (2.807,-0.757) (2.808,-0.722) (2.808,-0.686) (2.809,-0.651) (2.809,-0.616) (2.809,-0.581) (2.810,-0.546) (2.810,-0.510) (2.811,-0.475) (2.811,-0.440) (2.811,-0.405) (2.812,-0.370) (2.812,-0.334) (2.812,-0.299) (2.812,-0.264) (2.813,-0.229) (2.813,-0.194) (2.813,-0.158) (2.813,-0.123) (2.813,-0.088) (2.813,-0.053) (2.813,-0.018) (2.813,0.018) (2.813,0.053) (2.813,0.088) (2.813,0.123) (2.813,0.158) (2.813,0.194) (2.813,0.229) (2.812,0.264) (2.812,0.299) (2.812,0.334) (2.812,0.370) (2.811,0.405) (2.811,0.440) (2.811,0.475) (2.810,0.510) (2.810,0.546) (2.809,0.581) (2.809,0.616) (2.809,0.651) (2.808,0.686) (2.808,0.722) (2.807,0.757) (2.807,0.792) (2.806,0.827) (2.805,0.862) (2.805,0.898) (2.804,0.933) (2.804,0.968) (2.803,1.003) (2.802,1.038) (2.802,1.074) (2.801,1.109) (2.801,1.144) (2.800,1.179) (2.799,1.214) (2.799,1.250) (2.798,1.285) (2.797,1.320) (2.797,1.355) (2.796,1.390) (2.795,1.426) (2.795,1.461) (2.794,1.496) (2.793,1.531) (2.793,1.566) (2.792,1.602) (2.791,1.637) (2.791,1.672) (2.790,1.707) (2.789,1.742) (2.789,1.778) (2.788,1.813) (2.787,1.848) (2.787,1.883) (2.786,1.918) (2.785,1.954) (2.785,1.989) (2.784,2.024) (2.784,2.059) (2.783,2.094) (2.782,2.130) (2.782,2.165) (2.781,2.200)};
\addplot[equipot] coordinates {(2.286,-2.200) (2.287,-2.165) (2.287,-2.130) (2.288,-2.094) (2.289,-2.059) (2.290,-2.024) (2.291,-1.989) (2.292,-1.954) (2.293,-1.918) (2.294,-1.883) (2.295,-1.848) (2.296,-1.813) (2.297,-1.778) (2.298,-1.742) (2.299,-1.707) (2.300,-1.672) (2.301,-1.637) (2.302,-1.602) (2.303,-1.566) (2.304,-1.531) (2.306,-1.496) (2.307,-1.461) (2.308,-1.426) (2.309,-1.390) (2.310,-1.355) (2.311,-1.320) (2.312,-1.285) (2.314,-1.250) (2.315,-1.214) (2.316,-1.179) (2.317,-1.144) (2.318,-1.109) (2.319,-1.074) (2.321,-1.039) (2.322,-1.003) (2.323,-0.968) (2.324,-0.933) (2.325,-0.898) (2.326,-0.863) (2.327,-0.827) (2.328,-0.792) (2.329,-0.757) (2.330,-0.722) (2.331,-0.687) (2.332,-0.651) (2.333,-0.616) (2.334,-0.581) (2.335,-0.546) (2.335,-0.511) (2.336,-0.475) (2.337,-0.440) (2.338,-0.405) (2.338,-0.370) (2.339,-0.335) (2.339,-0.299) (2.340,-0.264) (2.340,-0.229) (2.340,-0.194) (2.341,-0.158) (2.341,-0.123) (2.341,-0.088) (2.341,-0.053) (2.341,-0.018) (2.341,0.018) (2.341,0.053) (2.341,0.088) (2.341,0.123) (2.341,0.158) (2.340,0.194) (2.340,0.229) (2.340,0.264) (2.339,0.299) (2.339,0.335) (2.338,0.370) (2.338,0.405) (2.337,0.440) (2.336,0.475) (2.335,0.511) (2.335,0.546) (2.334,0.581) (2.333,0.616) (2.332,0.651) (2.331,0.687) (2.330,0.722) (2.329,0.757) (2.328,0.792) (2.327,0.827) (2.326,0.863) (2.325,0.898) (2.324,0.933) (2.323,0.968) (2.322,1.003) (2.321,1.039) (2.319,1.074) (2.318,1.109) (2.317,1.144) (2.316,1.179) (2.315,1.214) (2.314,1.250) (2.312,1.285) (2.311,1.320) (2.310,1.355) (2.309,1.390) (2.308,1.426) (2.307,1.461) (2.306,1.496) (2.304,1.531) (2.303,1.566) (2.302,1.602) (2.301,1.637) (2.300,1.672) (2.299,1.707) (2.298,1.742) (2.297,1.778) (2.296,1.813) (2.295,1.848) (2.294,1.883) (2.293,1.918) (2.292,1.954) (2.291,1.989) (2.290,2.024) (2.289,2.059) (2.288,2.094) (2.287,2.130) (2.287,2.165) (2.286,2.200)};
\addplot[equipot] coordinates {(1.789,-2.200) (1.790,-2.165) (1.792,-2.130) (1.793,-2.094) (1.794,-2.059) (1.795,-2.024) (1.797,-1.989) (1.798,-1.953) (1.799,-1.918) (1.801,-1.883) (1.802,-1.848) (1.804,-1.813) (1.806,-1.777) (1.807,-1.742) (1.809,-1.707) (1.810,-1.672) (1.812,-1.637) (1.814,-1.601) (1.816,-1.566) (1.818,-1.531) (1.820,-1.496) (1.822,-1.461) (1.824,-1.425) (1.826,-1.390) (1.828,-1.355) (1.830,-1.320) (1.832,-1.285) (1.834,-1.250) (1.836,-1.214) (1.838,-1.179) (1.840,-1.144) (1.843,-1.109) (1.845,-1.074) (1.847,-1.039) (1.849,-1.003) (1.851,-0.968) (1.854,-0.933) (1.856,-0.898) (1.858,-0.863) (1.860,-0.827) (1.862,-0.792) (1.864,-0.757) (1.866,-0.722) (1.868,-0.687) (1.870,-0.652) (1.872,-0.616) (1.874,-0.581) (1.876,-0.546) (1.877,-0.511) (1.879,-0.476) (1.881,-0.440) (1.882,-0.405) (1.883,-0.370) (1.884,-0.335) (1.885,-0.300) (1.886,-0.264) (1.887,-0.229) (1.888,-0.194) (1.889,-0.159) (1.889,-0.123) (1.890,-0.088) (1.890,-0.053) (1.890,-0.018) (1.890,0.018) (1.890,0.053) (1.890,0.088) (1.889,0.123) (1.889,0.159) (1.888,0.194) (1.887,0.229) (1.886,0.264) (1.885,0.300) (1.884,0.335) (1.883,0.370) (1.882,0.405) (1.881,0.440) (1.879,0.476) (1.877,0.511) (1.876,0.546) (1.874,0.581) (1.872,0.616) (1.870,0.652) (1.868,0.687) (1.866,0.722) (1.864,0.757) (1.862,0.792) (1.860,0.827) (1.858,0.863) (1.856,0.898) (1.854,0.933) (1.851,0.968) (1.849,1.003) (1.847,1.039) (1.845,1.074) (1.843,1.109) (1.840,1.144) (1.838,1.179) (1.836,1.214) (1.834,1.250) (1.832,1.285) (1.830,1.320) (1.828,1.355) (1.826,1.390) (1.824,1.425) (1.822,1.461) (1.820,1.496) (1.818,1.531) (1.816,1.566) (1.814,1.601) (1.812,1.637) (1.810,1.672) (1.809,1.707) (1.807,1.742) (1.806,1.777) (1.804,1.813) (1.802,1.848) (1.801,1.883) (1.799,1.918) (1.798,1.953) (1.797,1.989) (1.795,2.024) (1.794,2.059) (1.793,2.094) (1.792,2.130) (1.790,2.165) (1.789,2.200)};
\addplot[equipot] coordinates {(1.289,-2.200) (1.290,-2.165) (1.292,-2.130) (1.293,-2.095) (1.295,-2.060) (1.297,-2.025) (1.298,-1.990) (1.300,-1.955) (1.302,-1.920) (1.304,-1.885) (1.306,-1.850) (1.308,-1.815) (1.311,-1.780) (1.313,-1.745) (1.316,-1.710) (1.318,-1.675) (1.321,-1.640) (1.324,-1.605) (1.326,-1.570) (1.329,-1.535) (1.332,-1.500) (1.336,-1.465) (1.339,-1.430) (1.342,-1.395) (1.346,-1.360) (1.350,-1.326) (1.353,-1.291) (1.357,-1.256) (1.361,-1.221) (1.365,-1.186) (1.370,-1.151) (1.374,-1.116) (1.378,-1.082) (1.383,-1.047) (1.387,-1.012) (1.392,-0.977) (1.396,-0.943) (1.401,-0.908) (1.405,-0.873) (1.410,-0.838) (1.414,-0.803) (1.419,-0.769) (1.423,-0.734) (1.428,-0.699) (1.432,-0.664) (1.436,-0.629) (1.440,-0.595) (1.444,-0.560) (1.448,-0.525) (1.452,-0.490) (1.455,-0.455) (1.459,-0.420) (1.462,-0.385) (1.464,-0.350) (1.467,-0.315) (1.469,-0.280) (1.472,-0.245) (1.473,-0.210) (1.475,-0.175) (1.476,-0.140) (1.477,-0.105) (1.478,-0.070) (1.479,-0.035) (1.479,0.000) (1.479,0.035) (1.478,0.070) (1.477,0.105) (1.476,0.140) (1.475,0.175) (1.473,0.210) (1.472,0.245) (1.469,0.280) (1.467,0.315) (1.464,0.350) (1.462,0.385) (1.459,0.420) (1.455,0.455) (1.452,0.490) (1.448,0.525) (1.444,0.560) (1.440,0.595) (1.436,0.629) (1.432,0.664) (1.428,0.699) (1.423,0.734) (1.419,0.769) (1.414,0.803) (1.410,0.838) (1.405,0.873) (1.401,0.908) (1.396,0.943) (1.392,0.977) (1.387,1.012) (1.383,1.047) (1.378,1.082) (1.374,1.116) (1.370,1.151) (1.365,1.186) (1.361,1.221) (1.357,1.256) (1.353,1.291) (1.350,1.326) (1.346,1.360) (1.342,1.395) (1.339,1.430) (1.336,1.465) (1.332,1.500) (1.329,1.535) (1.326,1.570) (1.324,1.605) (1.321,1.640) (1.318,1.675) (1.316,1.710) (1.313,1.745) (1.311,1.780) (1.308,1.815) (1.306,1.850) (1.304,1.885) (1.302,1.920) (1.300,1.955) (1.298,1.990) (1.297,2.025) (1.295,2.060) (1.293,2.095) (1.292,2.130) (1.290,2.165) (1.289,2.200)};
\addplot[equipot] coordinates {(0.781,-2.200) (0.782,-2.165) (0.784,-2.130) (0.785,-2.094) (0.787,-2.059) (0.788,-2.024) (0.790,-1.989) (0.792,-1.954) (0.794,-1.918) (0.797,-1.883) (0.799,-1.848) (0.801,-1.813) (0.804,-1.778) (0.807,-1.743) (0.810,-1.708) (0.813,-1.673) (0.817,-1.638) (0.820,-1.602) (0.824,-1.567) (0.828,-1.532) (0.833,-1.498) (0.837,-1.463) (0.842,-1.428) (0.848,-1.393) (0.853,-1.358) (0.859,-1.323) (0.866,-1.289) (0.872,-1.254) (0.879,-1.220) (0.887,-1.185) (0.894,-1.151) (0.903,-1.117) (0.911,-1.082) (0.920,-1.048) (0.929,-1.014) (0.938,-0.980) (0.947,-0.946) (0.957,-0.912) (0.967,-0.878) (0.977,-0.845) (0.987,-0.811) (0.997,-0.777) (1.007,-0.743) (1.017,-0.710) (1.027,-0.676) (1.036,-0.642) (1.046,-0.608) (1.055,-0.574) (1.063,-0.540) (1.072,-0.505) (1.080,-0.471) (1.088,-0.437) (1.095,-0.402) (1.102,-0.368) (1.108,-0.333) (1.113,-0.298) (1.119,-0.263) (1.123,-0.228) (1.127,-0.193) (1.130,-0.158) (1.133,-0.123) (1.135,-0.088) (1.136,-0.053) (1.137,-0.018) (1.137,0.018) (1.136,0.053) (1.135,0.088) (1.133,0.123) (1.130,0.158) (1.127,0.193) (1.123,0.228) (1.119,0.263) (1.113,0.298) (1.108,0.333) (1.102,0.368) (1.095,0.402) (1.088,0.437) (1.080,0.471) (1.072,0.505) (1.063,0.540) (1.055,0.574) (1.046,0.608) (1.036,0.642) (1.027,0.676) (1.017,0.710) (1.007,0.743) (0.997,0.777) (0.987,0.811) (0.977,0.845) (0.967,0.878) (0.957,0.912) (0.947,0.946) (0.938,0.980) (0.929,1.014) (0.920,1.048) (0.911,1.082) (0.903,1.117) (0.894,1.151) (0.887,1.185) (0.879,1.220) (0.872,1.254) (0.866,1.289) (0.859,1.323) (0.853,1.358) (0.848,1.393) (0.842,1.428) (0.837,1.463) (0.833,1.498) (0.828,1.532) (0.824,1.567) (0.820,1.602) (0.817,1.638) (0.813,1.673) (0.810,1.708) (0.807,1.743) (0.804,1.778) (0.801,1.813) (0.799,1.848) (0.797,1.883) (0.794,1.918) (0.792,1.954) (0.790,1.989) (0.788,2.024) (0.787,2.059) (0.785,2.094) (0.784,2.130) (0.782,2.165) (0.781,2.200)};
\addplot[equipot] coordinates {(0.262,-2.200) (0.263,-2.165) (0.263,-2.130) (0.264,-2.095) (0.265,-2.060) (0.266,-2.025) (0.266,-1.990) (0.267,-1.955) (0.268,-1.919) (0.270,-1.884) (0.271,-1.849) (0.272,-1.814) (0.273,-1.779) (0.275,-1.744) (0.277,-1.709) (0.278,-1.674) (0.280,-1.639) (0.283,-1.604) (0.285,-1.569) (0.288,-1.534) (0.291,-1.499) (0.294,-1.464) (0.298,-1.429) (0.302,-1.395) (0.307,-1.360) (0.312,-1.325) (0.318,-1.291) (0.324,-1.256) (0.332,-1.222) (0.340,-1.188) (0.350,-1.154) (0.360,-1.121) (0.373,-1.088) (0.386,-1.055) (0.401,-1.024) (0.417,-0.993) (0.435,-0.962) (0.453,-0.933) (0.473,-0.904) (0.493,-0.875) (0.500,-0.842) (0.500,-0.807) (0.500,-0.772) (0.500,-0.737) (0.500,-0.701) (0.500,-0.666) (0.500,-0.631) (0.500,-0.596) (0.500,-0.561) (0.500,-0.526) (0.500,-0.491) (0.500,-0.456) (0.500,-0.421) (0.500,-0.386) (0.500,-0.351) (0.500,-0.316) (0.500,-0.281) (0.500,-0.246) (0.500,-0.210) (0.500,-0.175) (0.500,-0.140) (0.500,-0.105) (0.500,-0.070) (0.500,-0.035) (0.500,0.000) (0.500,0.035) (0.500,0.070) (0.500,0.105) (0.500,0.140) (0.500,0.175) (0.500,0.210) (0.500,0.246) (0.500,0.281) (0.500,0.316) (0.500,0.351) (0.500,0.386) (0.500,0.421) (0.500,0.456) (0.500,0.491) (0.500,0.526) (0.500,0.561) (0.500,0.596) (0.500,0.631) (0.500,0.666) (0.500,0.701) (0.500,0.737) (0.500,0.772) (0.500,0.807) (0.500,0.842) (0.493,0.875) (0.473,0.904) (0.453,0.933) (0.435,0.962) (0.417,0.993) (0.401,1.024) (0.386,1.055) (0.373,1.088) (0.360,1.121) (0.350,1.154) (0.340,1.188) (0.332,1.222) (0.324,1.256) (0.318,1.291) (0.312,1.325) (0.307,1.360) (0.302,1.395) (0.298,1.429) (0.294,1.464) (0.291,1.499) (0.288,1.534) (0.285,1.569) (0.283,1.604) (0.280,1.639) (0.278,1.674) (0.277,1.709) (0.275,1.744) (0.273,1.779) (0.272,1.814) (0.271,1.849) (0.270,1.884) (0.268,1.919) (0.267,1.955) (0.266,1.990) (0.266,2.025) (0.265,2.060) (0.264,2.095) (0.263,2.130) (0.263,2.165) (0.262,2.200)};
\addplot[equipot] coordinates {(-0.262,-2.200) (-0.263,-2.165) (-0.263,-2.130) (-0.264,-2.095) (-0.265,-2.060) (-0.266,-2.025) (-0.266,-1.990) (-0.267,-1.955) (-0.268,-1.919) (-0.270,-1.884) (-0.271,-1.849) (-0.272,-1.814) (-0.273,-1.779) (-0.275,-1.744) (-0.277,-1.709) (-0.278,-1.674) (-0.280,-1.639) (-0.283,-1.604) (-0.285,-1.569) (-0.288,-1.534) (-0.291,-1.499) (-0.294,-1.464) (-0.298,-1.429) (-0.302,-1.395) (-0.307,-1.360) (-0.312,-1.325) (-0.318,-1.291) (-0.324,-1.256) (-0.332,-1.222) (-0.340,-1.188) (-0.350,-1.154) (-0.360,-1.121) (-0.373,-1.088) (-0.386,-1.055) (-0.401,-1.024) (-0.417,-0.993) (-0.435,-0.962) (-0.453,-0.933) (-0.473,-0.904) (-0.493,-0.875) (-0.500,-0.842) (-0.500,-0.807) (-0.500,-0.772) (-0.500,-0.737) (-0.500,-0.701) (-0.500,-0.666) (-0.500,-0.631) (-0.500,-0.596) (-0.500,-0.561) (-0.500,-0.526) (-0.500,-0.491) (-0.500,-0.456) (-0.500,-0.421) (-0.500,-0.386) (-0.500,-0.351) (-0.500,-0.316) (-0.500,-0.281) (-0.500,-0.246) (-0.500,-0.210) (-0.500,-0.175) (-0.500,-0.140) (-0.500,-0.105) (-0.500,-0.070) (-0.500,-0.035) (-0.500,-0.000) (-0.500,0.035) (-0.500,0.070) (-0.500,0.105) (-0.500,0.140) (-0.500,0.175) (-0.500,0.210) (-0.500,0.246) (-0.500,0.281) (-0.500,0.316) (-0.500,0.351) (-0.500,0.386) (-0.500,0.421) (-0.500,0.456) (-0.500,0.491) (-0.500,0.526) (-0.500,0.561) (-0.500,0.596) (-0.500,0.631) (-0.500,0.666) (-0.500,0.701) (-0.500,0.737) (-0.500,0.772) (-0.500,0.807) (-0.500,0.842) (-0.493,0.875) (-0.473,0.904) (-0.453,0.933) (-0.435,0.962) (-0.417,0.993) (-0.401,1.024) (-0.386,1.055) (-0.373,1.088) (-0.360,1.121) (-0.350,1.154) (-0.340,1.188) (-0.332,1.222) (-0.324,1.256) (-0.318,1.291) (-0.312,1.325) (-0.307,1.360) (-0.302,1.395) (-0.298,1.429) (-0.294,1.464) (-0.291,1.499) (-0.288,1.534) (-0.285,1.569) (-0.283,1.604) (-0.280,1.639) (-0.278,1.674) (-0.277,1.709) (-0.275,1.744) (-0.273,1.779) (-0.272,1.814) (-0.271,1.849) (-0.270,1.884) (-0.268,1.919) (-0.267,1.955) (-0.266,1.990) (-0.266,2.025) (-0.265,2.060) (-0.264,2.095) (-0.263,2.130) (-0.263,2.165) (-0.262,2.200)};
\addplot[equipot] coordinates {(-0.781,-2.200) (-0.782,-2.165) (-0.784,-2.130) (-0.785,-2.094) (-0.787,-2.059) (-0.788,-2.024) (-0.790,-1.989) (-0.792,-1.954) (-0.794,-1.918) (-0.797,-1.883) (-0.799,-1.848) (-0.801,-1.813) (-0.804,-1.778) (-0.807,-1.743) (-0.810,-1.708) (-0.813,-1.673) (-0.817,-1.638) (-0.820,-1.602) (-0.824,-1.567) (-0.828,-1.532) (-0.833,-1.498) (-0.837,-1.463) (-0.842,-1.428) (-0.848,-1.393) (-0.853,-1.358) (-0.859,-1.323) (-0.866,-1.289) (-0.872,-1.254) (-0.879,-1.220) (-0.887,-1.185) (-0.894,-1.151) (-0.903,-1.117) (-0.911,-1.082) (-0.920,-1.048) (-0.929,-1.014) (-0.938,-0.980) (-0.947,-0.946) (-0.957,-0.912) (-0.967,-0.878) (-0.977,-0.845) (-0.987,-0.811) (-0.997,-0.777) (-1.007,-0.743) (-1.017,-0.710) (-1.027,-0.676) (-1.036,-0.642) (-1.046,-0.608) (-1.055,-0.574) (-1.063,-0.540) (-1.072,-0.505) (-1.080,-0.471) (-1.088,-0.437) (-1.095,-0.402) (-1.102,-0.368) (-1.108,-0.333) (-1.113,-0.298) (-1.119,-0.263) (-1.123,-0.228) (-1.127,-0.193) (-1.130,-0.158) (-1.133,-0.123) (-1.135,-0.088) (-1.136,-0.053) (-1.137,-0.018) (-1.137,0.018) (-1.136,0.053) (-1.135,0.088) (-1.133,0.123) (-1.130,0.158) (-1.127,0.193) (-1.123,0.228) (-1.119,0.263) (-1.113,0.298) (-1.108,0.333) (-1.102,0.368) (-1.095,0.402) (-1.088,0.437) (-1.080,0.471) (-1.072,0.505) (-1.063,0.540) (-1.055,0.574) (-1.046,0.608) (-1.036,0.642) (-1.027,0.676) (-1.017,0.710) (-1.007,0.743) (-0.997,0.777) (-0.987,0.811) (-0.977,0.845) (-0.967,0.878) (-0.957,0.912) (-0.947,0.946) (-0.938,0.980) (-0.929,1.014) (-0.920,1.048) (-0.911,1.082) (-0.903,1.117) (-0.894,1.151) (-0.887,1.185) (-0.879,1.220) (-0.872,1.254) (-0.866,1.289) (-0.859,1.323) (-0.853,1.358) (-0.848,1.393) (-0.842,1.428) (-0.837,1.463) (-0.833,1.498) (-0.828,1.532) (-0.824,1.567) (-0.820,1.602) (-0.817,1.638) (-0.813,1.673) (-0.810,1.708) (-0.807,1.743) (-0.804,1.778) (-0.801,1.813) (-0.799,1.848) (-0.797,1.883) (-0.794,1.918) (-0.792,1.954) (-0.790,1.989) (-0.788,2.024) (-0.787,2.059) (-0.785,2.094) (-0.784,2.130) (-0.782,2.165) (-0.781,2.200)};
\addplot[equipot] coordinates {(-1.289,-2.200) (-1.290,-2.165) (-1.292,-2.130) (-1.293,-2.095) (-1.295,-2.060) (-1.297,-2.025) (-1.298,-1.990) (-1.300,-1.955) (-1.302,-1.920) (-1.304,-1.885) (-1.306,-1.850) (-1.308,-1.815) (-1.311,-1.780) (-1.313,-1.745) (-1.316,-1.710) (-1.318,-1.675) (-1.321,-1.640) (-1.324,-1.605) (-1.326,-1.570) (-1.329,-1.535) (-1.332,-1.500) (-1.336,-1.465) (-1.339,-1.430) (-1.342,-1.395) (-1.346,-1.360) (-1.350,-1.326) (-1.353,-1.291) (-1.357,-1.256) (-1.361,-1.221) (-1.365,-1.186) (-1.370,-1.151) (-1.374,-1.116) (-1.378,-1.082) (-1.383,-1.047) (-1.387,-1.012) (-1.392,-0.977) (-1.396,-0.943) (-1.401,-0.908) (-1.405,-0.873) (-1.410,-0.838) (-1.414,-0.803) (-1.419,-0.769) (-1.423,-0.734) (-1.428,-0.699) (-1.432,-0.664) (-1.436,-0.629) (-1.440,-0.595) (-1.444,-0.560) (-1.448,-0.525) (-1.452,-0.490) (-1.455,-0.455) (-1.459,-0.420) (-1.462,-0.385) (-1.464,-0.350) (-1.467,-0.315) (-1.469,-0.280) (-1.472,-0.245) (-1.473,-0.210) (-1.475,-0.175) (-1.476,-0.140) (-1.477,-0.105) (-1.478,-0.070) (-1.479,-0.035) (-1.479,0.000) (-1.479,0.035) (-1.478,0.070) (-1.477,0.105) (-1.476,0.140) (-1.475,0.175) (-1.473,0.210) (-1.472,0.245) (-1.469,0.280) (-1.467,0.315) (-1.464,0.350) (-1.462,0.385) (-1.459,0.420) (-1.455,0.455) (-1.452,0.490) (-1.448,0.525) (-1.444,0.560) (-1.440,0.595) (-1.436,0.629) (-1.432,0.664) (-1.428,0.699) (-1.423,0.734) (-1.419,0.769) (-1.414,0.803) (-1.410,0.838) (-1.405,0.873) (-1.401,0.908) (-1.396,0.943) (-1.392,0.977) (-1.387,1.012) (-1.383,1.047) (-1.378,1.082) (-1.374,1.116) (-1.370,1.151) (-1.365,1.186) (-1.361,1.221) (-1.357,1.256) (-1.353,1.291) (-1.350,1.326) (-1.346,1.360) (-1.342,1.395) (-1.339,1.430) (-1.336,1.465) (-1.332,1.500) (-1.329,1.535) (-1.326,1.570) (-1.324,1.605) (-1.321,1.640) (-1.318,1.675) (-1.316,1.710) (-1.313,1.745) (-1.311,1.780) (-1.308,1.815) (-1.306,1.850) (-1.304,1.885) (-1.302,1.920) (-1.300,1.955) (-1.298,1.990) (-1.297,2.025) (-1.295,2.060) (-1.293,2.095) (-1.292,2.130) (-1.290,2.165) (-1.289,2.200)};
\addplot[equipot] coordinates {(-1.789,-2.200) (-1.790,-2.165) (-1.792,-2.130) (-1.793,-2.094) (-1.794,-2.059) (-1.795,-2.024) (-1.797,-1.989) (-1.798,-1.953) (-1.799,-1.918) (-1.801,-1.883) (-1.802,-1.848) (-1.804,-1.813) (-1.806,-1.777) (-1.807,-1.742) (-1.809,-1.707) (-1.810,-1.672) (-1.812,-1.637) (-1.814,-1.601) (-1.816,-1.566) (-1.818,-1.531) (-1.820,-1.496) (-1.822,-1.461) (-1.824,-1.425) (-1.826,-1.390) (-1.828,-1.355) (-1.830,-1.320) (-1.832,-1.285) (-1.834,-1.250) (-1.836,-1.214) (-1.838,-1.179) (-1.840,-1.144) (-1.843,-1.109) (-1.845,-1.074) (-1.847,-1.039) (-1.849,-1.003) (-1.851,-0.968) (-1.854,-0.933) (-1.856,-0.898) (-1.858,-0.863) (-1.860,-0.827) (-1.862,-0.792) (-1.864,-0.757) (-1.866,-0.722) (-1.868,-0.687) (-1.870,-0.652) (-1.872,-0.616) (-1.874,-0.581) (-1.876,-0.546) (-1.877,-0.511) (-1.879,-0.476) (-1.881,-0.440) (-1.882,-0.405) (-1.883,-0.370) (-1.884,-0.335) (-1.885,-0.300) (-1.886,-0.264) (-1.887,-0.229) (-1.888,-0.194) (-1.889,-0.159) (-1.889,-0.123) (-1.890,-0.088) (-1.890,-0.053) (-1.890,-0.018) (-1.890,0.018) (-1.890,0.053) (-1.890,0.088) (-1.889,0.123) (-1.889,0.159) (-1.888,0.194) (-1.887,0.229) (-1.886,0.264) (-1.885,0.300) (-1.884,0.335) (-1.883,0.370) (-1.882,0.405) (-1.881,0.440) (-1.879,0.476) (-1.877,0.511) (-1.876,0.546) (-1.874,0.581) (-1.872,0.616) (-1.870,0.652) (-1.868,0.687) (-1.866,0.722) (-1.864,0.757) (-1.862,0.792) (-1.860,0.827) (-1.858,0.863) (-1.856,0.898) (-1.854,0.933) (-1.851,0.968) (-1.849,1.003) (-1.847,1.039) (-1.845,1.074) (-1.843,1.109) (-1.840,1.144) (-1.838,1.179) (-1.836,1.214) (-1.834,1.250) (-1.832,1.285) (-1.830,1.320) (-1.828,1.355) (-1.826,1.390) (-1.824,1.425) (-1.822,1.461) (-1.820,1.496) (-1.818,1.531) (-1.816,1.566) (-1.814,1.601) (-1.812,1.637) (-1.810,1.672) (-1.809,1.707) (-1.807,1.742) (-1.806,1.777) (-1.804,1.813) (-1.802,1.848) (-1.801,1.883) (-1.799,1.918) (-1.798,1.953) (-1.797,1.989) (-1.795,2.024) (-1.794,2.059) (-1.793,2.094) (-1.792,2.130) (-1.790,2.165) (-1.789,2.200)};
\addplot[equipot] coordinates {(-2.286,-2.200) (-2.287,-2.165) (-2.287,-2.130) (-2.288,-2.094) (-2.289,-2.059) (-2.290,-2.024) (-2.291,-1.989) (-2.292,-1.954) (-2.293,-1.918) (-2.294,-1.883) (-2.295,-1.848) (-2.296,-1.813) (-2.297,-1.778) (-2.298,-1.742) (-2.299,-1.707) (-2.300,-1.672) (-2.301,-1.637) (-2.302,-1.602) (-2.303,-1.566) (-2.304,-1.531) (-2.306,-1.496) (-2.307,-1.461) (-2.308,-1.426) (-2.309,-1.390) (-2.310,-1.355) (-2.311,-1.320) (-2.312,-1.285) (-2.314,-1.250) (-2.315,-1.214) (-2.316,-1.179) (-2.317,-1.144) (-2.318,-1.109) (-2.319,-1.074) (-2.321,-1.039) (-2.322,-1.003) (-2.323,-0.968) (-2.324,-0.933) (-2.325,-0.898) (-2.326,-0.863) (-2.327,-0.827) (-2.328,-0.792) (-2.329,-0.757) (-2.330,-0.722) (-2.331,-0.687) (-2.332,-0.651) (-2.333,-0.616) (-2.334,-0.581) (-2.335,-0.546) (-2.335,-0.511) (-2.336,-0.475) (-2.337,-0.440) (-2.338,-0.405) (-2.338,-0.370) (-2.339,-0.335) (-2.339,-0.299) (-2.340,-0.264) (-2.340,-0.229) (-2.340,-0.194) (-2.341,-0.158) (-2.341,-0.123) (-2.341,-0.088) (-2.341,-0.053) (-2.341,-0.018) (-2.341,0.018) (-2.341,0.053) (-2.341,0.088) (-2.341,0.123) (-2.341,0.158) (-2.340,0.194) (-2.340,0.229) (-2.340,0.264) (-2.339,0.299) (-2.339,0.335) (-2.338,0.370) (-2.338,0.405) (-2.337,0.440) (-2.336,0.475) (-2.335,0.511) (-2.335,0.546) (-2.334,0.581) (-2.333,0.616) (-2.332,0.651) (-2.331,0.687) (-2.330,0.722) (-2.329,0.757) (-2.328,0.792) (-2.327,0.827) (-2.326,0.863) (-2.325,0.898) (-2.324,0.933) (-2.323,0.968) (-2.322,1.003) (-2.321,1.039) (-2.319,1.074) (-2.318,1.109) (-2.317,1.144) (-2.316,1.179) (-2.315,1.214) (-2.314,1.250) (-2.312,1.285) (-2.311,1.320) (-2.310,1.355) (-2.309,1.390) (-2.308,1.426) (-2.307,1.461) (-2.306,1.496) (-2.304,1.531) (-2.303,1.566) (-2.302,1.602) (-2.301,1.637) (-2.300,1.672) (-2.299,1.707) (-2.298,1.742) (-2.297,1.778) (-2.296,1.813) (-2.295,1.848) (-2.294,1.883) (-2.293,1.918) (-2.292,1.954) (-2.291,1.989) (-2.290,2.024) (-2.289,2.059) (-2.288,2.094) (-2.287,2.130) (-2.287,2.165) (-2.286,2.200)};
\addplot[equipot] coordinates {(-2.781,-2.200) (-2.782,-2.165) (-2.782,-2.130) (-2.783,-2.094) (-2.784,-2.059) (-2.784,-2.024) (-2.785,-1.989) (-2.785,-1.954) (-2.786,-1.918) (-2.787,-1.883) (-2.787,-1.848) (-2.788,-1.813) (-2.789,-1.778) (-2.789,-1.742) (-2.790,-1.707) (-2.791,-1.672) (-2.791,-1.637) (-2.792,-1.602) (-2.793,-1.566) (-2.793,-1.531) (-2.794,-1.496) (-2.795,-1.461) (-2.795,-1.426) (-2.796,-1.390) (-2.797,-1.355) (-2.797,-1.320) (-2.798,-1.285) (-2.799,-1.250) (-2.799,-1.214) (-2.800,-1.179) (-2.801,-1.144) (-2.801,-1.109) (-2.802,-1.074) (-2.802,-1.038) (-2.803,-1.003) (-2.804,-0.968) (-2.804,-0.933) (-2.805,-0.898) (-2.805,-0.862) (-2.806,-0.827) (-2.807,-0.792) (-2.807,-0.757) (-2.808,-0.722) (-2.808,-0.686) (-2.809,-0.651) (-2.809,-0.616) (-2.809,-0.581) (-2.810,-0.546) (-2.810,-0.510) (-2.811,-0.475) (-2.811,-0.440) (-2.811,-0.405) (-2.812,-0.370) (-2.812,-0.334) (-2.812,-0.299) (-2.812,-0.264) (-2.813,-0.229) (-2.813,-0.194) (-2.813,-0.158) (-2.813,-0.123) (-2.813,-0.088) (-2.813,-0.053) (-2.813,-0.018) (-2.813,0.018) (-2.813,0.053) (-2.813,0.088) (-2.813,0.123) (-2.813,0.158) (-2.813,0.194) (-2.813,0.229) (-2.812,0.264) (-2.812,0.299) (-2.812,0.334) (-2.812,0.370) (-2.811,0.405) (-2.811,0.440) (-2.811,0.475) (-2.810,0.510) (-2.810,0.546) (-2.809,0.581) (-2.809,0.616) (-2.809,0.651) (-2.808,0.686) (-2.808,0.722) (-2.807,0.757) (-2.807,0.792) (-2.806,0.827) (-2.805,0.862) (-2.805,0.898) (-2.804,0.933) (-2.804,0.968) (-2.803,1.003) (-2.802,1.038) (-2.802,1.074) (-2.801,1.109) (-2.801,1.144) (-2.800,1.179) (-2.799,1.214) (-2.799,1.250) (-2.798,1.285) (-2.797,1.320) (-2.797,1.355) (-2.796,1.390) (-2.795,1.426) (-2.795,1.461) (-2.794,1.496) (-2.793,1.531) (-2.793,1.566) (-2.792,1.602) (-2.791,1.637) (-2.791,1.672) (-2.790,1.707) (-2.789,1.742) (-2.789,1.778) (-2.788,1.813) (-2.787,1.848) (-2.787,1.883) (-2.786,1.918) (-2.785,1.954) (-2.785,1.989) (-2.784,2.024) (-2.784,2.059) (-2.783,2.094) (-2.782,2.130) (-2.782,2.165) (-2.781,2.200)};
\addplot[lineaE, flechaen=0.577] coordinates {(0.990,0.141) (1.025,0.145) (1.060,0.148) (1.095,0.152) (1.130,0.155) (1.166,0.157) (1.201,0.160) (1.236,0.162) (1.271,0.165) (1.307,0.167) (1.342,0.169) (1.377,0.171) (1.412,0.172) (1.447,0.174) (1.483,0.175) (1.518,0.177) (1.553,0.178) (1.589,0.179) (1.624,0.180) (1.659,0.181) (1.694,0.182) (1.730,0.183) (1.765,0.184) (1.800,0.185) (1.835,0.186) (1.871,0.186) (1.906,0.187) (1.941,0.188) (1.977,0.188) (2.012,0.189) (2.047,0.189) (2.082,0.190) (2.118,0.190) (2.153,0.191) (2.188,0.191) (2.224,0.192) (2.259,0.192) (2.294,0.192) (2.330,0.193) (2.365,0.193) (2.400,0.193) (2.435,0.193) (2.471,0.194) (2.506,0.194) (2.541,0.194) (2.577,0.194) (2.612,0.195) (2.647,0.195) (2.682,0.195) (2.718,0.195) (2.753,0.195) (2.788,0.196) (2.824,0.196) (2.859,0.196) (2.894,0.196) (2.929,0.196) (2.965,0.196) (3.000,0.196)};
\addplot[lineaE, flechaen=0.423] coordinates {(-3.000,0.196) (-2.965,0.196) (-2.929,0.196) (-2.894,0.196) (-2.859,0.196) (-2.824,0.196) (-2.788,0.196) (-2.753,0.195) (-2.718,0.195) (-2.682,0.195) (-2.647,0.195) (-2.612,0.195) (-2.577,0.194) (-2.541,0.194) (-2.506,0.194) (-2.471,0.194) (-2.435,0.193) (-2.400,0.193) (-2.365,0.193) (-2.330,0.193) (-2.294,0.192) (-2.259,0.192) (-2.224,0.192) (-2.188,0.191) (-2.153,0.191) (-2.118,0.190) (-2.082,0.190) (-2.047,0.189) (-2.012,0.189) (-1.977,0.188) (-1.941,0.188) (-1.906,0.187) (-1.871,0.186) (-1.835,0.186) (-1.800,0.185) (-1.765,0.184) (-1.730,0.183) (-1.694,0.182) (-1.659,0.181) (-1.624,0.180) (-1.589,0.179) (-1.553,0.178) (-1.518,0.177) (-1.483,0.175) (-1.447,0.174) (-1.412,0.172) (-1.377,0.171) (-1.342,0.169) (-1.307,0.167) (-1.271,0.165) (-1.236,0.162) (-1.201,0.160) (-1.166,0.157) (-1.130,0.155) (-1.095,0.152) (-1.060,0.148) (-1.025,0.145) (-0.990,0.141)};
\addplot[lineaE, flechaen=0.585] coordinates {(0.959,0.282) (0.994,0.290) (1.028,0.297) (1.063,0.303) (1.098,0.309) (1.133,0.315) (1.168,0.320) (1.203,0.325) (1.238,0.329) (1.273,0.333) (1.308,0.337) (1.343,0.341) (1.378,0.344) (1.413,0.348) (1.448,0.351) (1.483,0.353) (1.519,0.356) (1.554,0.358) (1.589,0.360) (1.624,0.363) (1.659,0.365) (1.695,0.366) (1.730,0.368) (1.765,0.370) (1.800,0.371) (1.836,0.373) (1.871,0.374) (1.906,0.375) (1.942,0.376) (1.977,0.378) (2.012,0.379) (2.047,0.380) (2.083,0.380) (2.118,0.381) (2.153,0.382) (2.188,0.383) (2.224,0.384) (2.259,0.384) (2.294,0.385) (2.330,0.386) (2.365,0.386) (2.400,0.387) (2.435,0.387) (2.471,0.388) (2.506,0.388) (2.541,0.389) (2.577,0.389) (2.612,0.390) (2.647,0.390) (2.682,0.390) (2.718,0.391) (2.753,0.391) (2.788,0.391) (2.824,0.392) (2.859,0.392) (2.894,0.392) (2.929,0.392) (2.965,0.393) (3.000,0.393)};
\addplot[lineaE, flechaen=0.415] coordinates {(-3.000,0.393) (-2.965,0.393) (-2.929,0.392) (-2.894,0.392) (-2.859,0.392) (-2.824,0.392) (-2.788,0.391) (-2.753,0.391) (-2.718,0.391) (-2.682,0.390) (-2.647,0.390) (-2.612,0.390) (-2.577,0.389) (-2.541,0.389) (-2.506,0.388) (-2.471,0.388) (-2.435,0.387) (-2.400,0.387) (-2.365,0.386) (-2.330,0.386) (-2.294,0.385) (-2.259,0.384) (-2.224,0.384) (-2.188,0.383) (-2.153,0.382) (-2.118,0.381) (-2.083,0.380) (-2.047,0.380) (-2.012,0.379) (-1.977,0.378) (-1.942,0.376) (-1.906,0.375) (-1.871,0.374) (-1.836,0.373) (-1.800,0.371) (-1.765,0.370) (-1.730,0.368) (-1.695,0.366) (-1.659,0.365) (-1.624,0.363) (-1.589,0.360) (-1.554,0.358) (-1.519,0.356) (-1.483,0.353) (-1.448,0.351) (-1.413,0.348) (-1.378,0.344) (-1.343,0.341) (-1.308,0.337) (-1.273,0.333) (-1.238,0.329) (-1.203,0.325) (-1.168,0.320) (-1.133,0.315) (-1.098,0.309) (-1.063,0.303) (-1.028,0.297) (-0.994,0.290) (-0.959,0.282)};
\addplot[lineaE, flechaen=0.597] coordinates {(0.906,0.424) (0.939,0.435) (0.973,0.445) (1.006,0.455) (1.040,0.463) (1.075,0.472) (1.109,0.479) (1.143,0.486) (1.178,0.493) (1.212,0.499) (1.247,0.505) (1.282,0.511) (1.316,0.516) (1.351,0.520) (1.386,0.525) (1.421,0.529) (1.456,0.533) (1.491,0.536) (1.526,0.540) (1.561,0.543) (1.596,0.546) (1.631,0.549) (1.666,0.551) (1.701,0.554) (1.736,0.556) (1.771,0.558) (1.806,0.560) (1.841,0.562) (1.876,0.564) (1.911,0.565) (1.946,0.567) (1.982,0.568) (2.017,0.570) (2.052,0.571) (2.087,0.572) (2.122,0.574) (2.157,0.575) (2.192,0.576) (2.227,0.577) (2.262,0.578) (2.298,0.579) (2.333,0.579) (2.368,0.580) (2.403,0.581) (2.438,0.582) (2.473,0.582) (2.508,0.583) (2.543,0.584) (2.579,0.584) (2.614,0.585) (2.649,0.586) (2.684,0.586) (2.719,0.587) (2.754,0.587) (2.789,0.587) (2.824,0.588) (2.859,0.588) (2.895,0.589) (2.930,0.589) (2.965,0.589) (3.000,0.590)};
\addplot[lineaE, flechaen=0.403] coordinates {(-3.000,0.590) (-2.965,0.589) (-2.930,0.589) (-2.895,0.589) (-2.859,0.588) (-2.824,0.588) (-2.789,0.587) (-2.754,0.587) (-2.719,0.587) (-2.684,0.586) (-2.649,0.586) (-2.614,0.585) (-2.579,0.584) (-2.543,0.584) (-2.508,0.583) (-2.473,0.582) (-2.438,0.582) (-2.403,0.581) (-2.368,0.580) (-2.333,0.579) (-2.298,0.579) (-2.262,0.578) (-2.227,0.577) (-2.192,0.576) (-2.157,0.575) (-2.122,0.574) (-2.087,0.572) (-2.052,0.571) (-2.017,0.570) (-1.982,0.568) (-1.946,0.567) (-1.911,0.565) (-1.876,0.564) (-1.841,0.562) (-1.806,0.560) (-1.771,0.558) (-1.736,0.556) (-1.701,0.554) (-1.666,0.551) (-1.631,0.549) (-1.596,0.546) (-1.561,0.543) (-1.526,0.540) (-1.491,0.536) (-1.456,0.533) (-1.421,0.529) (-1.386,0.525) (-1.351,0.520) (-1.316,0.516) (-1.282,0.511) (-1.247,0.505) (-1.212,0.499) (-1.178,0.493) (-1.143,0.486) (-1.109,0.479) (-1.075,0.472) (-1.040,0.463) (-1.006,0.455) (-0.973,0.445) (-0.939,0.435) (-0.906,0.424)};
\addplot[lineaE, flechaen=0.613] coordinates {(0.825,0.565) (0.857,0.580) (0.890,0.593) (0.923,0.606) (0.957,0.618) (0.990,0.629) (1.024,0.639) (1.058,0.648) (1.093,0.657) (1.127,0.665) (1.162,0.673) (1.197,0.680) (1.231,0.687) (1.266,0.693) (1.301,0.699) (1.336,0.704) (1.371,0.709) (1.407,0.714) (1.442,0.719) (1.477,0.723) (1.512,0.727) (1.548,0.730) (1.583,0.734) (1.618,0.737) (1.653,0.740) (1.689,0.743) (1.724,0.746) (1.760,0.748) (1.795,0.751) (1.830,0.753) (1.866,0.755) (1.901,0.757) (1.937,0.759) (1.972,0.761) (2.007,0.762) (2.043,0.764) (2.078,0.766) (2.114,0.767) (2.149,0.768) (2.185,0.770) (2.220,0.771) (2.255,0.772) (2.291,0.773) (2.326,0.774) (2.362,0.775) (2.397,0.776) (2.433,0.777) (2.468,0.778) (2.504,0.779) (2.539,0.779) (2.575,0.780) (2.610,0.781) (2.645,0.782) (2.681,0.782) (2.716,0.783) (2.752,0.783) (2.787,0.784) (2.823,0.785) (2.858,0.785) (2.894,0.786) (2.929,0.786) (2.965,0.786) (3.000,0.787)};
\addplot[lineaE, flechaen=0.387] coordinates {(-3.000,0.787) (-2.965,0.786) (-2.929,0.786) (-2.894,0.786) (-2.858,0.785) (-2.823,0.785) (-2.787,0.784) (-2.752,0.783) (-2.716,0.783) (-2.681,0.782) (-2.645,0.782) (-2.610,0.781) (-2.575,0.780) (-2.539,0.779) (-2.504,0.779) (-2.468,0.778) (-2.433,0.777) (-2.397,0.776) (-2.362,0.775) (-2.326,0.774) (-2.291,0.773) (-2.255,0.772) (-2.220,0.771) (-2.185,0.770) (-2.149,0.768) (-2.114,0.767) (-2.078,0.766) (-2.043,0.764) (-2.007,0.762) (-1.972,0.761) (-1.937,0.759) (-1.901,0.757) (-1.866,0.755) (-1.830,0.753) (-1.795,0.751) (-1.760,0.748) (-1.724,0.746) (-1.689,0.743) (-1.653,0.740) (-1.618,0.737) (-1.583,0.734) (-1.548,0.730) (-1.512,0.727) (-1.477,0.723) (-1.442,0.719) (-1.407,0.714) (-1.371,0.709) (-1.336,0.704) (-1.301,0.699) (-1.266,0.693) (-1.231,0.687) (-1.197,0.680) (-1.162,0.673) (-1.127,0.665) (-1.093,0.657) (-1.058,0.648) (-1.024,0.639) (-0.990,0.629) (-0.957,0.618) (-0.923,0.606) (-0.890,0.593) (-0.857,0.580) (-0.825,0.565)};
\addplot[lineaE, flechaen=0.635] coordinates {(0.708,0.707) (0.738,0.725) (0.769,0.741) (0.801,0.756) (0.833,0.771) (0.866,0.784) (0.899,0.796) (0.932,0.808) (0.966,0.818) (1.000,0.828) (1.034,0.838) (1.068,0.847) (1.102,0.855) (1.137,0.862) (1.171,0.870) (1.206,0.876) (1.240,0.883) (1.275,0.889) (1.310,0.894) (1.345,0.899) (1.380,0.904) (1.415,0.909) (1.450,0.913) (1.485,0.917) (1.520,0.921) (1.555,0.925) (1.590,0.928) (1.625,0.932) (1.661,0.935) (1.696,0.938) (1.731,0.940) (1.766,0.943) (1.801,0.945) (1.836,0.948) (1.872,0.950) (1.907,0.952) (1.942,0.954) (1.977,0.956) (2.013,0.958) (2.048,0.959) (2.083,0.961) (2.118,0.963) (2.154,0.964) (2.189,0.965) (2.224,0.967) (2.259,0.968) (2.295,0.969) (2.330,0.970) (2.365,0.971) (2.400,0.972) (2.436,0.973) (2.471,0.974) (2.506,0.975) (2.541,0.976) (2.577,0.977) (2.612,0.978) (2.647,0.978) (2.683,0.979) (2.718,0.980) (2.753,0.981) (2.788,0.981) (2.824,0.982) (2.859,0.982) (2.894,0.983) (2.929,0.983) (2.965,0.984) (3.000,0.984)};
\addplot[lineaE, flechaen=0.365] coordinates {(-3.000,0.984) (-2.965,0.984) (-2.929,0.983) (-2.894,0.983) (-2.859,0.982) (-2.824,0.982) (-2.788,0.981) (-2.753,0.981) (-2.718,0.980) (-2.683,0.979) (-2.647,0.978) (-2.612,0.978) (-2.577,0.977) (-2.541,0.976) (-2.506,0.975) (-2.471,0.974) (-2.436,0.973) (-2.400,0.972) (-2.365,0.971) (-2.330,0.970) (-2.295,0.969) (-2.259,0.968) (-2.224,0.967) (-2.189,0.965) (-2.154,0.964) (-2.118,0.963) (-2.083,0.961) (-2.048,0.959) (-2.013,0.958) (-1.977,0.956) (-1.942,0.954) (-1.907,0.952) (-1.872,0.950) (-1.836,0.948) (-1.801,0.945) (-1.766,0.943) (-1.731,0.940) (-1.696,0.938) (-1.661,0.935) (-1.625,0.932) (-1.590,0.928) (-1.555,0.925) (-1.520,0.921) (-1.485,0.917) (-1.450,0.913) (-1.415,0.909) (-1.380,0.904) (-1.345,0.899) (-1.310,0.894) (-1.275,0.889) (-1.240,0.883) (-1.206,0.876) (-1.171,0.870) (-1.137,0.862) (-1.102,0.855) (-1.068,0.847) (-1.034,0.838) (-1.000,0.828) (-0.966,0.818) (-0.932,0.808) (-0.899,0.796) (-0.866,0.784) (-0.833,0.771) (-0.801,0.756) (-0.769,0.741) (-0.738,0.725) (-0.708,0.707)};
\addplot[lineaE, flechaen=0.662] coordinates {(0.531,0.848) (0.559,0.869) (0.589,0.888) (0.620,0.905) (0.651,0.921) (0.683,0.936) (0.716,0.951) (0.749,0.964) (0.782,0.976) (0.816,0.988) (0.849,0.999) (0.883,1.009) (0.917,1.018) (0.952,1.027) (0.986,1.036) (1.021,1.044) (1.055,1.051) (1.090,1.059) (1.125,1.065) (1.160,1.072) (1.195,1.078) (1.230,1.083) (1.265,1.089) (1.300,1.094) (1.335,1.098) (1.371,1.103) (1.406,1.107) (1.441,1.111) (1.476,1.115) (1.512,1.119) (1.547,1.122) (1.582,1.126) (1.618,1.129) (1.653,1.132) (1.688,1.134) (1.724,1.137) (1.759,1.140) (1.794,1.142) (1.830,1.144) (1.865,1.147) (1.901,1.149) (1.936,1.151) (1.972,1.153) (2.007,1.154) (2.042,1.156) (2.078,1.158) (2.113,1.159) (2.149,1.161) (2.184,1.162) (2.220,1.164) (2.255,1.165) (2.291,1.166) (2.326,1.167) (2.361,1.168) (2.397,1.170) (2.432,1.171) (2.468,1.172) (2.503,1.173) (2.539,1.173) (2.574,1.174) (2.610,1.175) (2.645,1.176) (2.681,1.177) (2.716,1.178) (2.752,1.178) (2.787,1.179) (2.823,1.180) (2.858,1.180) (2.894,1.181) (2.929,1.181) (2.965,1.182) (3.000,1.182)};
\addplot[lineaE, flechaen=0.338] coordinates {(-3.000,1.182) (-2.965,1.182) (-2.929,1.181) (-2.894,1.181) (-2.858,1.180) (-2.823,1.180) (-2.787,1.179) (-2.752,1.178) (-2.716,1.178) (-2.681,1.177) (-2.645,1.176) (-2.610,1.175) (-2.574,1.174) (-2.539,1.173) (-2.503,1.173) (-2.468,1.172) (-2.432,1.171) (-2.397,1.170) (-2.361,1.168) (-2.326,1.167) (-2.291,1.166) (-2.255,1.165) (-2.220,1.164) (-2.184,1.162) (-2.149,1.161) (-2.113,1.159) (-2.078,1.158) (-2.042,1.156) (-2.007,1.154) (-1.972,1.153) (-1.936,1.151) (-1.901,1.149) (-1.865,1.147) (-1.830,1.144) (-1.794,1.142) (-1.759,1.140) (-1.724,1.137) (-1.688,1.134) (-1.653,1.132) (-1.618,1.129) (-1.582,1.126) (-1.547,1.122) (-1.512,1.119) (-1.476,1.115) (-1.441,1.111) (-1.406,1.107) (-1.371,1.103) (-1.335,1.098) (-1.300,1.094) (-1.265,1.089) (-1.230,1.083) (-1.195,1.078) (-1.160,1.072) (-1.125,1.065) (-1.090,1.059) (-1.055,1.051) (-1.021,1.044) (-0.986,1.036) (-0.952,1.027) (-0.917,1.018) (-0.883,1.009) (-0.849,0.999) (-0.816,0.988) (-0.782,0.976) (-0.749,0.964) (-0.716,0.951) (-0.683,0.936) (-0.651,0.921) (-0.620,0.905) (-0.589,0.888) (-0.559,0.869) (-0.531,0.848)};
\addplot[lineaE, flechaen=0.707] coordinates {(0.148,0.989) (0.178,1.005) (0.210,1.019) (0.243,1.033) (0.275,1.047) (0.308,1.061) (0.340,1.074) (0.373,1.087) (0.406,1.100) (0.440,1.112) (0.473,1.124) (0.507,1.135) (0.540,1.146) (0.574,1.157) (0.608,1.167) (0.642,1.177) (0.676,1.186) (0.710,1.195) (0.745,1.204) (0.779,1.212) (0.814,1.220) (0.848,1.227) (0.883,1.235) (0.918,1.241) (0.953,1.248) (0.987,1.254) (1.022,1.260) (1.057,1.266) (1.092,1.271) (1.127,1.277) (1.162,1.282) (1.197,1.287) (1.232,1.291) (1.268,1.295) (1.303,1.300) (1.338,1.304) (1.373,1.307) (1.408,1.311) (1.444,1.314) (1.479,1.318) (1.514,1.321) (1.549,1.324) (1.585,1.327) (1.620,1.330) (1.655,1.332) (1.691,1.335) (1.726,1.337) (1.761,1.340) (1.797,1.342) (1.832,1.344) (1.867,1.346) (1.903,1.348) (1.938,1.350) (1.974,1.351) (2.009,1.353) (2.044,1.355) (2.080,1.356) (2.115,1.358) (2.150,1.359) (2.186,1.361) (2.221,1.362) (2.257,1.363) (2.292,1.365) (2.327,1.366) (2.363,1.367) (2.398,1.368) (2.434,1.369) (2.469,1.370) (2.504,1.371) (2.540,1.372) (2.575,1.373) (2.611,1.373) (2.646,1.374) (2.681,1.375) (2.717,1.376) (2.752,1.377) (2.788,1.377) (2.823,1.378) (2.858,1.379) (2.894,1.379) (2.929,1.380) (2.965,1.380) (3.000,1.381)};
\addplot[lineaE, flechaen=0.293] coordinates {(-3.000,1.381) (-2.965,1.380) (-2.929,1.380) (-2.894,1.379) (-2.858,1.379) (-2.823,1.378) (-2.788,1.377) (-2.752,1.377) (-2.717,1.376) (-2.681,1.375) (-2.646,1.374) (-2.611,1.373) (-2.575,1.373) (-2.540,1.372) (-2.504,1.371) (-2.469,1.370) (-2.434,1.369) (-2.398,1.368) (-2.363,1.367) (-2.327,1.366) (-2.292,1.365) (-2.257,1.363) (-2.221,1.362) (-2.186,1.361) (-2.150,1.359) (-2.115,1.358) (-2.080,1.356) (-2.044,1.355) (-2.009,1.353) (-1.974,1.351) (-1.938,1.350) (-1.903,1.348) (-1.867,1.346) (-1.832,1.344) (-1.797,1.342) (-1.761,1.340) (-1.726,1.337) (-1.691,1.335) (-1.655,1.332) (-1.620,1.330) (-1.585,1.327) (-1.549,1.324) (-1.514,1.321) (-1.479,1.318) (-1.444,1.314) (-1.408,1.311) (-1.373,1.307) (-1.338,1.304) (-1.303,1.300) (-1.268,1.295) (-1.232,1.291) (-1.197,1.287) (-1.162,1.282) (-1.127,1.277) (-1.092,1.271) (-1.057,1.266) (-1.022,1.260) (-0.987,1.254) (-0.953,1.248) (-0.918,1.241) (-0.883,1.235) (-0.848,1.227) (-0.814,1.220) (-0.779,1.212) (-0.745,1.204) (-0.710,1.195) (-0.676,1.186) (-0.642,1.177) (-0.608,1.167) (-0.574,1.157) (-0.540,1.146) (-0.507,1.135) (-0.473,1.124) (-0.440,1.112) (-0.406,1.100) (-0.373,1.087) (-0.340,1.074) (-0.308,1.061) (-0.275,1.047) (-0.243,1.033) (-0.210,1.019) (-0.178,1.005) (-0.148,0.989)};
\addplot[lineaE, flechaen=0.141, flechaen=0.859] coordinates {(-3.000,1.580) (-2.965,1.579) (-2.930,1.579) (-2.895,1.578) (-2.860,1.578) (-2.825,1.577) (-2.790,1.576) (-2.755,1.576) (-2.720,1.575) (-2.685,1.574) (-2.650,1.573) (-2.615,1.573) (-2.580,1.572) (-2.545,1.571) (-2.510,1.570) (-2.475,1.569) (-2.440,1.568) (-2.405,1.567) (-2.370,1.566) (-2.335,1.565) (-2.300,1.564) (-2.265,1.563) (-2.230,1.562) (-2.195,1.561) (-2.160,1.559) (-2.125,1.558) (-2.090,1.557) (-2.055,1.555) (-2.020,1.554) (-1.985,1.552) (-1.950,1.551) (-1.915,1.549) (-1.880,1.547) (-1.845,1.546) (-1.810,1.544) (-1.775,1.542) (-1.740,1.540) (-1.705,1.538) (-1.670,1.536) (-1.635,1.533) (-1.600,1.531) (-1.565,1.529) (-1.530,1.526) (-1.495,1.524) (-1.460,1.521) (-1.425,1.518) (-1.390,1.515) (-1.355,1.512) (-1.320,1.509) (-1.286,1.506) (-1.251,1.503) (-1.216,1.499) (-1.181,1.496) (-1.146,1.492) (-1.111,1.488) (-1.077,1.484) (-1.042,1.480) (-1.007,1.476) (-0.972,1.471) (-0.937,1.467) (-0.903,1.462) (-0.868,1.458) (-0.833,1.453) (-0.799,1.448) (-0.764,1.443) (-0.729,1.438) (-0.695,1.432) (-0.660,1.427) (-0.626,1.421) (-0.591,1.416) (-0.556,1.410) (-0.522,1.405) (-0.487,1.399) (-0.453,1.394) (-0.418,1.388) (-0.383,1.383) (-0.349,1.378) (-0.314,1.373) (-0.279,1.368) (-0.245,1.364) (-0.210,1.360) (-0.175,1.357) (-0.140,1.354) (-0.105,1.352) (-0.070,1.350) (-0.035,1.349) (0.000,1.348) (0.035,1.349) (0.070,1.350) (0.105,1.352) (0.140,1.354) (0.175,1.357) (0.210,1.360) (0.245,1.364) (0.279,1.368) (0.314,1.373) (0.349,1.378) (0.383,1.383) (0.418,1.388) (0.453,1.394) (0.487,1.399) (0.522,1.405) (0.556,1.410) (0.591,1.416) (0.626,1.421) (0.660,1.427) (0.695,1.432) (0.729,1.438) (0.764,1.443) (0.799,1.448) (0.833,1.453) (0.868,1.458) (0.903,1.462) (0.937,1.467) (0.972,1.471) (1.007,1.476) (1.042,1.480) (1.077,1.484) (1.111,1.488) (1.146,1.492) (1.181,1.496) (1.216,1.499) (1.251,1.503) (1.286,1.506) (1.320,1.509) (1.355,1.512) (1.390,1.515) (1.425,1.518) (1.460,1.521) (1.495,1.524) (1.530,1.526) (1.565,1.529) (1.600,1.531) (1.635,1.533) (1.670,1.536) (1.705,1.538) (1.740,1.540) (1.775,1.542) (1.810,1.544) (1.845,1.546) (1.880,1.547) (1.915,1.549) (1.950,1.551) (1.985,1.552) (2.020,1.554) (2.055,1.555) (2.090,1.557) (2.125,1.558) (2.160,1.559) (2.195,1.561) (2.230,1.562) (2.265,1.563) (2.300,1.564) (2.335,1.565) (2.370,1.566) (2.405,1.567) (2.440,1.568) (2.475,1.569) (2.510,1.570) (2.545,1.571) (2.580,1.572) (2.615,1.573) (2.650,1.573) (2.685,1.574) (2.720,1.575) (2.755,1.576) (2.790,1.576) (2.825,1.577) (2.860,1.578) (2.895,1.578) (2.930,1.579) (2.965,1.579) (3.000,1.580)};
\addplot[lineaE, flechaen=0.141, flechaen=0.859] coordinates {(-3.000,1.779) (-2.965,1.779) (-2.930,1.778) (-2.895,1.777) (-2.859,1.777) (-2.824,1.776) (-2.789,1.776) (-2.754,1.775) (-2.719,1.774) (-2.684,1.774) (-2.649,1.773) (-2.613,1.772) (-2.578,1.771) (-2.543,1.770) (-2.508,1.770) (-2.473,1.769) (-2.438,1.768) (-2.403,1.767) (-2.367,1.766) (-2.332,1.765) (-2.297,1.764) (-2.262,1.763) (-2.227,1.762) (-2.192,1.761) (-2.157,1.760) (-2.121,1.758) (-2.086,1.757) (-2.051,1.756) (-2.016,1.755) (-1.981,1.753) (-1.946,1.752) (-1.911,1.750) (-1.876,1.749) (-1.841,1.747) (-1.805,1.746) (-1.770,1.744) (-1.735,1.742) (-1.700,1.741) (-1.665,1.739) (-1.630,1.737) (-1.595,1.735) (-1.560,1.733) (-1.525,1.731) (-1.489,1.729) (-1.454,1.727) (-1.419,1.724) (-1.384,1.722) (-1.349,1.720) (-1.314,1.717) (-1.279,1.715) (-1.244,1.712) (-1.209,1.709) (-1.174,1.707) (-1.139,1.704) (-1.104,1.701) (-1.069,1.698) (-1.034,1.695) (-0.999,1.692) (-0.964,1.689) (-0.929,1.686) (-0.894,1.682) (-0.859,1.679) (-0.824,1.676) (-0.789,1.673) (-0.754,1.669) (-0.719,1.666) (-0.684,1.663) (-0.649,1.659) (-0.614,1.656) (-0.579,1.653) (-0.544,1.649) (-0.509,1.646) (-0.474,1.643) (-0.439,1.640) (-0.404,1.637) (-0.369,1.635) (-0.334,1.632) (-0.299,1.630) (-0.263,1.628) (-0.228,1.626) (-0.193,1.624) (-0.158,1.622) (-0.123,1.621) (-0.088,1.620) (-0.053,1.620) (-0.018,1.619) (0.018,1.619) (0.053,1.620) (0.088,1.620) (0.123,1.621) (0.158,1.622) (0.193,1.624) (0.228,1.626) (0.263,1.628) (0.299,1.630) (0.334,1.632) (0.369,1.635) (0.404,1.637) (0.439,1.640) (0.474,1.643) (0.509,1.646) (0.544,1.649) (0.579,1.653) (0.614,1.656) (0.649,1.659) (0.684,1.663) (0.719,1.666) (0.754,1.669) (0.789,1.673) (0.824,1.676) (0.859,1.679) (0.894,1.682) (0.929,1.686) (0.964,1.689) (0.999,1.692) (1.034,1.695) (1.069,1.698) (1.104,1.701) (1.139,1.704) (1.174,1.707) (1.209,1.709) (1.244,1.712) (1.279,1.715) (1.314,1.717) (1.349,1.720) (1.384,1.722) (1.419,1.724) (1.454,1.727) (1.489,1.729) (1.525,1.731) (1.560,1.733) (1.595,1.735) (1.630,1.737) (1.665,1.739) (1.700,1.741) (1.735,1.742) (1.770,1.744) (1.805,1.746) (1.841,1.747) (1.876,1.749) (1.911,1.750) (1.946,1.752) (1.981,1.753) (2.016,1.755) (2.051,1.756) (2.086,1.757) (2.121,1.758) (2.157,1.760) (2.192,1.761) (2.227,1.762) (2.262,1.763) (2.297,1.764) (2.332,1.765) (2.367,1.766) (2.403,1.767) (2.438,1.768) (2.473,1.769) (2.508,1.770) (2.543,1.770) (2.578,1.771) (2.613,1.772) (2.649,1.773) (2.684,1.774) (2.719,1.774) (2.754,1.775) (2.789,1.776) (2.824,1.776) (2.859,1.777) (2.895,1.777) (2.930,1.778) (2.965,1.779) (3.000,1.779)};
\addplot[lineaE, flechaen=0.142, flechaen=0.858] coordinates {(-3.000,1.979) (-2.965,1.978) (-2.930,1.978) (-2.895,1.977) (-2.860,1.977) (-2.824,1.976) (-2.789,1.975) (-2.754,1.975) (-2.719,1.974) (-2.684,1.973) (-2.649,1.973) (-2.614,1.972) (-2.579,1.971) (-2.544,1.971) (-2.508,1.970) (-2.473,1.969) (-2.438,1.968) (-2.403,1.967) (-2.368,1.967) (-2.333,1.966) (-2.298,1.965) (-2.263,1.964) (-2.228,1.963) (-2.192,1.962) (-2.157,1.961) (-2.122,1.960) (-2.087,1.959) (-2.052,1.958) (-2.017,1.956) (-1.982,1.955) (-1.947,1.954) (-1.912,1.953) (-1.877,1.951) (-1.841,1.950) (-1.806,1.949) (-1.771,1.947) (-1.736,1.946) (-1.701,1.944) (-1.666,1.943) (-1.631,1.941) (-1.596,1.940) (-1.561,1.938) (-1.526,1.936) (-1.491,1.935) (-1.455,1.933) (-1.420,1.931) (-1.385,1.929) (-1.350,1.927) (-1.315,1.926) (-1.280,1.924) (-1.245,1.922) (-1.210,1.920) (-1.175,1.917) (-1.140,1.915) (-1.105,1.913) (-1.070,1.911) (-1.035,1.909) (-1.000,1.907) (-0.965,1.905) (-0.930,1.902) (-0.895,1.900) (-0.859,1.898) (-0.824,1.896) (-0.789,1.893) (-0.754,1.891) (-0.719,1.889) (-0.684,1.887) (-0.649,1.885) (-0.614,1.882) (-0.579,1.880) (-0.544,1.878) (-0.509,1.876) (-0.474,1.875) (-0.439,1.873) (-0.404,1.871) (-0.369,1.869) (-0.334,1.868) (-0.298,1.867) (-0.263,1.865) (-0.228,1.864) (-0.193,1.863) (-0.158,1.862) (-0.123,1.862) (-0.088,1.861) (-0.053,1.861) (-0.018,1.861) (0.018,1.861) (0.053,1.861) (0.088,1.861) (0.123,1.862) (0.158,1.862) (0.193,1.863) (0.228,1.864) (0.263,1.865) (0.298,1.867) (0.334,1.868) (0.369,1.869) (0.404,1.871) (0.439,1.873) (0.474,1.875) (0.509,1.876) (0.544,1.878) (0.579,1.880) (0.614,1.882) (0.649,1.885) (0.684,1.887) (0.719,1.889) (0.754,1.891) (0.789,1.893) (0.824,1.896) (0.859,1.898) (0.895,1.900) (0.930,1.902) (0.965,1.905) (1.000,1.907) (1.035,1.909) (1.070,1.911) (1.105,1.913) (1.140,1.915) (1.175,1.917) (1.210,1.920) (1.245,1.922) (1.280,1.924) (1.315,1.926) (1.350,1.927) (1.385,1.929) (1.420,1.931) (1.455,1.933) (1.491,1.935) (1.526,1.936) (1.561,1.938) (1.596,1.940) (1.631,1.941) (1.666,1.943) (1.701,1.944) (1.736,1.946) (1.771,1.947) (1.806,1.949) (1.841,1.950) (1.877,1.951) (1.912,1.953) (1.947,1.954) (1.982,1.955) (2.017,1.956) (2.052,1.958) (2.087,1.959) (2.122,1.960) (2.157,1.961) (2.192,1.962) (2.228,1.963) (2.263,1.964) (2.298,1.965) (2.333,1.966) (2.368,1.967) (2.403,1.967) (2.438,1.968) (2.473,1.969) (2.508,1.970) (2.544,1.971) (2.579,1.971) (2.614,1.972) (2.649,1.973) (2.684,1.973) (2.719,1.974) (2.754,1.975) (2.789,1.975) (2.824,1.976) (2.860,1.977) (2.895,1.977) (2.930,1.978) (2.965,1.978) (3.000,1.979)};
\addplot[lineaE, flechaen=0.142, flechaen=0.858] coordinates {(-3.000,2.179) (-2.965,2.178) (-2.930,2.178) (-2.895,2.177) (-2.860,2.177) (-2.824,2.176) (-2.789,2.176) (-2.754,2.175) (-2.719,2.174) (-2.684,2.174) (-2.649,2.173) (-2.614,2.173) (-2.579,2.172) (-2.544,2.171) (-2.509,2.170) (-2.473,2.170) (-2.438,2.169) (-2.403,2.168) (-2.368,2.167) (-2.333,2.167) (-2.298,2.166) (-2.263,2.165) (-2.228,2.164) (-2.193,2.163) (-2.158,2.162) (-2.123,2.161) (-2.087,2.161) (-2.052,2.160) (-2.017,2.159) (-1.982,2.158) (-1.947,2.156) (-1.912,2.155) (-1.877,2.154) (-1.842,2.153) (-1.807,2.152) (-1.772,2.151) (-1.737,2.150) (-1.701,2.148) (-1.666,2.147) (-1.631,2.146) (-1.596,2.145) (-1.561,2.143) (-1.526,2.142) (-1.491,2.141) (-1.456,2.139) (-1.421,2.138) (-1.386,2.136) (-1.351,2.135) (-1.316,2.133) (-1.281,2.132) (-1.245,2.130) (-1.210,2.129) (-1.175,2.127) (-1.140,2.126) (-1.105,2.124) (-1.070,2.122) (-1.035,2.121) (-1.000,2.119) (-0.965,2.118) (-0.930,2.116) (-0.895,2.114) (-0.860,2.113) (-0.825,2.111) (-0.790,2.110) (-0.754,2.108) (-0.719,2.107) (-0.684,2.105) (-0.649,2.104) (-0.614,2.102) (-0.579,2.101) (-0.544,2.100) (-0.509,2.098) (-0.474,2.097) (-0.439,2.096) (-0.404,2.095) (-0.369,2.094) (-0.333,2.093) (-0.298,2.092) (-0.263,2.091) (-0.228,2.090) (-0.193,2.090) (-0.158,2.089) (-0.123,2.089) (-0.088,2.089) (-0.053,2.088) (-0.018,2.088) (0.018,2.088) (0.053,2.088) (0.088,2.089) (0.123,2.089) (0.158,2.089) (0.193,2.090) (0.228,2.090) (0.263,2.091) (0.298,2.092) (0.333,2.093) (0.369,2.094) (0.404,2.095) (0.439,2.096) (0.474,2.097) (0.509,2.098) (0.544,2.100) (0.579,2.101) (0.614,2.102) (0.649,2.104) (0.684,2.105) (0.719,2.107) (0.754,2.108) (0.790,2.110) (0.825,2.111) (0.860,2.113) (0.895,2.114) (0.930,2.116) (0.965,2.118) (1.000,2.119) (1.035,2.121) (1.070,2.122) (1.105,2.124) (1.140,2.126) (1.175,2.127) (1.210,2.129) (1.245,2.130) (1.281,2.132) (1.316,2.133) (1.351,2.135) (1.386,2.136) (1.421,2.138) (1.456,2.139) (1.491,2.141) (1.526,2.142) (1.561,2.143) (1.596,2.145) (1.631,2.146) (1.666,2.147) (1.701,2.148) (1.737,2.150) (1.772,2.151) (1.807,2.152) (1.842,2.153) (1.877,2.154) (1.912,2.155) (1.947,2.156) (1.982,2.158) (2.017,2.159) (2.052,2.160) (2.087,2.161) (2.123,2.161) (2.158,2.162) (2.193,2.163) (2.228,2.164) (2.263,2.165) (2.298,2.166) (2.333,2.167) (2.368,2.167) (2.403,2.168) (2.438,2.169) (2.473,2.170) (2.509,2.170) (2.544,2.171) (2.579,2.172) (2.614,2.173) (2.649,2.173) (2.684,2.174) (2.719,2.174) (2.754,2.175) (2.789,2.176) (2.824,2.176) (2.860,2.177) (2.895,2.177) (2.930,2.178) (2.965,2.178) (3.000,2.179)};
\addplot[lineaE, flechaen=0.239] coordinates {(-0.959,0.283) (-0.924,0.283) (-0.888,0.283) (-0.853,0.283) (-0.817,0.283) (-0.781,0.283) (-0.746,0.283) (-0.710,0.283) (-0.675,0.283) (-0.639,0.283) (-0.604,0.283) (-0.568,0.283) (-0.533,0.283) (-0.497,0.283) (-0.462,0.283) (-0.426,0.283) (-0.391,0.283) (-0.355,0.283) (-0.320,0.283) (-0.284,0.283) (-0.249,0.283) (-0.213,0.283) (-0.178,0.283) (-0.142,0.283) (-0.107,0.283) (-0.071,0.283) (-0.036,0.283) (-0.000,0.283) (0.036,0.283) (0.071,0.283) (0.107,0.283) (0.142,0.283) (0.178,0.283) (0.213,0.283) (0.249,0.283) (0.284,0.283) (0.320,0.283) (0.355,0.283) (0.391,0.283) (0.426,0.283) (0.462,0.283) (0.497,0.283) (0.533,0.283) (0.568,0.283) (0.604,0.283) (0.639,0.283) (0.675,0.283) (0.710,0.283) (0.746,0.283) (0.781,0.283) (0.817,0.283) (0.853,0.283) (0.888,0.283) (0.924,0.283) (0.959,0.283)};
\addplot[lineaE, flechaen=0.804] coordinates {(-0.824,0.567) (-0.789,0.566) (-0.754,0.566) (-0.719,0.566) (-0.684,0.566) (-0.649,0.566) (-0.614,0.566) (-0.579,0.566) (-0.544,0.566) (-0.508,0.566) (-0.473,0.566) (-0.438,0.566) (-0.403,0.566) (-0.368,0.566) (-0.333,0.566) (-0.298,0.566) (-0.263,0.566) (-0.228,0.566) (-0.193,0.566) (-0.158,0.566) (-0.123,0.566) (-0.088,0.566) (-0.053,0.566) (-0.018,0.566) (0.018,0.566) (0.053,0.566) (0.088,0.566) (0.123,0.566) (0.158,0.566) (0.193,0.566) (0.228,0.566) (0.263,0.566) (0.298,0.566) (0.333,0.566) (0.368,0.566) (0.403,0.566) (0.438,0.566) (0.473,0.566) (0.508,0.566) (0.544,0.566) (0.579,0.566) (0.614,0.566) (0.649,0.566) (0.684,0.566) (0.719,0.566) (0.754,0.566) (0.789,0.566) (0.824,0.567)};
\addplot[lineaE, flechaen=0.500] coordinates {(-0.525,0.851) (-0.491,0.849) (-0.456,0.849) (-0.421,0.849) (-0.386,0.849) (-0.351,0.849) (-0.316,0.849) (-0.281,0.849) (-0.246,0.849) (-0.210,0.849) (-0.175,0.849) (-0.140,0.849) (-0.105,0.849) (-0.070,0.849) (-0.035,0.849) (0.000,0.849) (0.035,0.849) (0.070,0.849) (0.105,0.849) (0.140,0.849) (0.175,0.849) (0.210,0.849) (0.246,0.849) (0.281,0.849) (0.316,0.849) (0.351,0.849) (0.386,0.849) (0.421,0.849) (0.456,0.849) (0.491,0.849) (0.525,0.851)};
\addplot[lineaE, flechaen=0.577] coordinates {(0.990,-0.141) (1.025,-0.145) (1.060,-0.148) (1.095,-0.152) (1.130,-0.155) (1.166,-0.157) (1.201,-0.160) (1.236,-0.162) (1.271,-0.165) (1.307,-0.167) (1.342,-0.169) (1.377,-0.171) (1.412,-0.172) (1.447,-0.174) (1.483,-0.175) (1.518,-0.177) (1.553,-0.178) (1.589,-0.179) (1.624,-0.180) (1.659,-0.181) (1.694,-0.182) (1.730,-0.183) (1.765,-0.184) (1.800,-0.185) (1.835,-0.186) (1.871,-0.186) (1.906,-0.187) (1.941,-0.188) (1.977,-0.188) (2.012,-0.189) (2.047,-0.189) (2.082,-0.190) (2.118,-0.190) (2.153,-0.191) (2.188,-0.191) (2.224,-0.192) (2.259,-0.192) (2.294,-0.192) (2.330,-0.193) (2.365,-0.193) (2.400,-0.193) (2.435,-0.193) (2.471,-0.194) (2.506,-0.194) (2.541,-0.194) (2.577,-0.194) (2.612,-0.195) (2.647,-0.195) (2.682,-0.195) (2.718,-0.195) (2.753,-0.195) (2.788,-0.196) (2.824,-0.196) (2.859,-0.196) (2.894,-0.196) (2.929,-0.196) (2.965,-0.196) (3.000,-0.196)};
\addplot[lineaE, flechaen=0.423] coordinates {(-3.000,-0.196) (-2.965,-0.196) (-2.929,-0.196) (-2.894,-0.196) (-2.859,-0.196) (-2.824,-0.196) (-2.788,-0.196) (-2.753,-0.195) (-2.718,-0.195) (-2.682,-0.195) (-2.647,-0.195) (-2.612,-0.195) (-2.577,-0.194) (-2.541,-0.194) (-2.506,-0.194) (-2.471,-0.194) (-2.435,-0.193) (-2.400,-0.193) (-2.365,-0.193) (-2.330,-0.193) (-2.294,-0.192) (-2.259,-0.192) (-2.224,-0.192) (-2.188,-0.191) (-2.153,-0.191) (-2.118,-0.190) (-2.082,-0.190) (-2.047,-0.189) (-2.012,-0.189) (-1.977,-0.188) (-1.941,-0.188) (-1.906,-0.187) (-1.871,-0.186) (-1.835,-0.186) (-1.800,-0.185) (-1.765,-0.184) (-1.730,-0.183) (-1.694,-0.182) (-1.659,-0.181) (-1.624,-0.180) (-1.589,-0.179) (-1.553,-0.178) (-1.518,-0.177) (-1.483,-0.175) (-1.447,-0.174) (-1.412,-0.172) (-1.377,-0.171) (-1.342,-0.169) (-1.307,-0.167) (-1.271,-0.165) (-1.236,-0.162) (-1.201,-0.160) (-1.166,-0.157) (-1.130,-0.155) (-1.095,-0.152) (-1.060,-0.148) (-1.025,-0.145) (-0.990,-0.141)};
\addplot[lineaE, flechaen=0.585] coordinates {(0.959,-0.282) (0.994,-0.290) (1.028,-0.297) (1.063,-0.303) (1.098,-0.309) (1.133,-0.315) (1.168,-0.320) (1.203,-0.325) (1.238,-0.329) (1.273,-0.333) (1.308,-0.337) (1.343,-0.341) (1.378,-0.344) (1.413,-0.348) (1.448,-0.351) (1.483,-0.353) (1.519,-0.356) (1.554,-0.358) (1.589,-0.360) (1.624,-0.363) (1.659,-0.365) (1.695,-0.366) (1.730,-0.368) (1.765,-0.370) (1.800,-0.371) (1.836,-0.373) (1.871,-0.374) (1.906,-0.375) (1.942,-0.376) (1.977,-0.378) (2.012,-0.379) (2.047,-0.380) (2.083,-0.380) (2.118,-0.381) (2.153,-0.382) (2.188,-0.383) (2.224,-0.384) (2.259,-0.384) (2.294,-0.385) (2.330,-0.386) (2.365,-0.386) (2.400,-0.387) (2.435,-0.387) (2.471,-0.388) (2.506,-0.388) (2.541,-0.389) (2.577,-0.389) (2.612,-0.390) (2.647,-0.390) (2.682,-0.390) (2.718,-0.391) (2.753,-0.391) (2.788,-0.391) (2.824,-0.392) (2.859,-0.392) (2.894,-0.392) (2.929,-0.392) (2.965,-0.393) (3.000,-0.393)};
\addplot[lineaE, flechaen=0.415] coordinates {(-3.000,-0.393) (-2.965,-0.393) (-2.929,-0.392) (-2.894,-0.392) (-2.859,-0.392) (-2.824,-0.392) (-2.788,-0.391) (-2.753,-0.391) (-2.718,-0.391) (-2.682,-0.390) (-2.647,-0.390) (-2.612,-0.390) (-2.577,-0.389) (-2.541,-0.389) (-2.506,-0.388) (-2.471,-0.388) (-2.435,-0.387) (-2.400,-0.387) (-2.365,-0.386) (-2.330,-0.386) (-2.294,-0.385) (-2.259,-0.384) (-2.224,-0.384) (-2.188,-0.383) (-2.153,-0.382) (-2.118,-0.381) (-2.083,-0.380) (-2.047,-0.380) (-2.012,-0.379) (-1.977,-0.378) (-1.942,-0.376) (-1.906,-0.375) (-1.871,-0.374) (-1.836,-0.373) (-1.800,-0.371) (-1.765,-0.370) (-1.730,-0.368) (-1.695,-0.366) (-1.659,-0.365) (-1.624,-0.363) (-1.589,-0.360) (-1.554,-0.358) (-1.519,-0.356) (-1.483,-0.353) (-1.448,-0.351) (-1.413,-0.348) (-1.378,-0.344) (-1.343,-0.341) (-1.308,-0.337) (-1.273,-0.333) (-1.238,-0.329) (-1.203,-0.325) (-1.168,-0.320) (-1.133,-0.315) (-1.098,-0.309) (-1.063,-0.303) (-1.028,-0.297) (-0.994,-0.290) (-0.959,-0.282)};
\addplot[lineaE, flechaen=0.597] coordinates {(0.906,-0.424) (0.939,-0.435) (0.973,-0.445) (1.006,-0.455) (1.040,-0.463) (1.075,-0.472) (1.109,-0.479) (1.143,-0.486) (1.178,-0.493) (1.212,-0.499) (1.247,-0.505) (1.282,-0.511) (1.316,-0.516) (1.351,-0.520) (1.386,-0.525) (1.421,-0.529) (1.456,-0.533) (1.491,-0.536) (1.526,-0.540) (1.561,-0.543) (1.596,-0.546) (1.631,-0.549) (1.666,-0.551) (1.701,-0.554) (1.736,-0.556) (1.771,-0.558) (1.806,-0.560) (1.841,-0.562) (1.876,-0.564) (1.911,-0.565) (1.946,-0.567) (1.982,-0.568) (2.017,-0.570) (2.052,-0.571) (2.087,-0.572) (2.122,-0.574) (2.157,-0.575) (2.192,-0.576) (2.227,-0.577) (2.262,-0.578) (2.298,-0.579) (2.333,-0.579) (2.368,-0.580) (2.403,-0.581) (2.438,-0.582) (2.473,-0.582) (2.508,-0.583) (2.543,-0.584) (2.579,-0.584) (2.614,-0.585) (2.649,-0.586) (2.684,-0.586) (2.719,-0.587) (2.754,-0.587) (2.789,-0.587) (2.824,-0.588) (2.859,-0.588) (2.895,-0.589) (2.930,-0.589) (2.965,-0.589) (3.000,-0.590)};
\addplot[lineaE, flechaen=0.403] coordinates {(-3.000,-0.590) (-2.965,-0.589) (-2.930,-0.589) (-2.895,-0.589) (-2.859,-0.588) (-2.824,-0.588) (-2.789,-0.587) (-2.754,-0.587) (-2.719,-0.587) (-2.684,-0.586) (-2.649,-0.586) (-2.614,-0.585) (-2.579,-0.584) (-2.543,-0.584) (-2.508,-0.583) (-2.473,-0.582) (-2.438,-0.582) (-2.403,-0.581) (-2.368,-0.580) (-2.333,-0.579) (-2.298,-0.579) (-2.262,-0.578) (-2.227,-0.577) (-2.192,-0.576) (-2.157,-0.575) (-2.122,-0.574) (-2.087,-0.572) (-2.052,-0.571) (-2.017,-0.570) (-1.982,-0.568) (-1.946,-0.567) (-1.911,-0.565) (-1.876,-0.564) (-1.841,-0.562) (-1.806,-0.560) (-1.771,-0.558) (-1.736,-0.556) (-1.701,-0.554) (-1.666,-0.551) (-1.631,-0.549) (-1.596,-0.546) (-1.561,-0.543) (-1.526,-0.540) (-1.491,-0.536) (-1.456,-0.533) (-1.421,-0.529) (-1.386,-0.525) (-1.351,-0.520) (-1.316,-0.516) (-1.282,-0.511) (-1.247,-0.505) (-1.212,-0.499) (-1.178,-0.493) (-1.143,-0.486) (-1.109,-0.479) (-1.075,-0.472) (-1.040,-0.463) (-1.006,-0.455) (-0.973,-0.445) (-0.939,-0.435) (-0.906,-0.424)};
\addplot[lineaE, flechaen=0.613] coordinates {(0.825,-0.565) (0.857,-0.580) (0.890,-0.593) (0.923,-0.606) (0.957,-0.618) (0.990,-0.629) (1.024,-0.639) (1.058,-0.648) (1.093,-0.657) (1.127,-0.665) (1.162,-0.673) (1.197,-0.680) (1.231,-0.687) (1.266,-0.693) (1.301,-0.699) (1.336,-0.704) (1.371,-0.709) (1.407,-0.714) (1.442,-0.719) (1.477,-0.723) (1.512,-0.727) (1.548,-0.730) (1.583,-0.734) (1.618,-0.737) (1.653,-0.740) (1.689,-0.743) (1.724,-0.746) (1.760,-0.748) (1.795,-0.751) (1.830,-0.753) (1.866,-0.755) (1.901,-0.757) (1.937,-0.759) (1.972,-0.761) (2.007,-0.762) (2.043,-0.764) (2.078,-0.766) (2.114,-0.767) (2.149,-0.768) (2.185,-0.770) (2.220,-0.771) (2.255,-0.772) (2.291,-0.773) (2.326,-0.774) (2.362,-0.775) (2.397,-0.776) (2.433,-0.777) (2.468,-0.778) (2.504,-0.779) (2.539,-0.779) (2.575,-0.780) (2.610,-0.781) (2.645,-0.782) (2.681,-0.782) (2.716,-0.783) (2.752,-0.783) (2.787,-0.784) (2.823,-0.785) (2.858,-0.785) (2.894,-0.786) (2.929,-0.786) (2.965,-0.786) (3.000,-0.787)};
\addplot[lineaE, flechaen=0.387] coordinates {(-3.000,-0.787) (-2.965,-0.786) (-2.929,-0.786) (-2.894,-0.786) (-2.858,-0.785) (-2.823,-0.785) (-2.787,-0.784) (-2.752,-0.783) (-2.716,-0.783) (-2.681,-0.782) (-2.645,-0.782) (-2.610,-0.781) (-2.575,-0.780) (-2.539,-0.779) (-2.504,-0.779) (-2.468,-0.778) (-2.433,-0.777) (-2.397,-0.776) (-2.362,-0.775) (-2.326,-0.774) (-2.291,-0.773) (-2.255,-0.772) (-2.220,-0.771) (-2.185,-0.770) (-2.149,-0.768) (-2.114,-0.767) (-2.078,-0.766) (-2.043,-0.764) (-2.007,-0.762) (-1.972,-0.761) (-1.937,-0.759) (-1.901,-0.757) (-1.866,-0.755) (-1.830,-0.753) (-1.795,-0.751) (-1.760,-0.748) (-1.724,-0.746) (-1.689,-0.743) (-1.653,-0.740) (-1.618,-0.737) (-1.583,-0.734) (-1.548,-0.730) (-1.512,-0.727) (-1.477,-0.723) (-1.442,-0.719) (-1.407,-0.714) (-1.371,-0.709) (-1.336,-0.704) (-1.301,-0.699) (-1.266,-0.693) (-1.231,-0.687) (-1.197,-0.680) (-1.162,-0.673) (-1.127,-0.665) (-1.093,-0.657) (-1.058,-0.648) (-1.024,-0.639) (-0.990,-0.629) (-0.957,-0.618) (-0.923,-0.606) (-0.890,-0.593) (-0.857,-0.580) (-0.825,-0.565)};
\addplot[lineaE, flechaen=0.635] coordinates {(0.708,-0.707) (0.738,-0.725) (0.769,-0.741) (0.801,-0.756) (0.833,-0.771) (0.866,-0.784) (0.899,-0.796) (0.932,-0.808) (0.966,-0.818) (1.000,-0.828) (1.034,-0.838) (1.068,-0.847) (1.102,-0.855) (1.137,-0.862) (1.171,-0.870) (1.206,-0.876) (1.240,-0.883) (1.275,-0.889) (1.310,-0.894) (1.345,-0.899) (1.380,-0.904) (1.415,-0.909) (1.450,-0.913) (1.485,-0.917) (1.520,-0.921) (1.555,-0.925) (1.590,-0.928) (1.625,-0.932) (1.661,-0.935) (1.696,-0.938) (1.731,-0.940) (1.766,-0.943) (1.801,-0.945) (1.836,-0.948) (1.872,-0.950) (1.907,-0.952) (1.942,-0.954) (1.977,-0.956) (2.013,-0.958) (2.048,-0.959) (2.083,-0.961) (2.118,-0.963) (2.154,-0.964) (2.189,-0.965) (2.224,-0.967) (2.259,-0.968) (2.295,-0.969) (2.330,-0.970) (2.365,-0.971) (2.400,-0.972) (2.436,-0.973) (2.471,-0.974) (2.506,-0.975) (2.541,-0.976) (2.577,-0.977) (2.612,-0.978) (2.647,-0.978) (2.683,-0.979) (2.718,-0.980) (2.753,-0.981) (2.788,-0.981) (2.824,-0.982) (2.859,-0.982) (2.894,-0.983) (2.929,-0.983) (2.965,-0.984) (3.000,-0.984)};
\addplot[lineaE, flechaen=0.365] coordinates {(-3.000,-0.984) (-2.965,-0.984) (-2.929,-0.983) (-2.894,-0.983) (-2.859,-0.982) (-2.824,-0.982) (-2.788,-0.981) (-2.753,-0.981) (-2.718,-0.980) (-2.683,-0.979) (-2.647,-0.978) (-2.612,-0.978) (-2.577,-0.977) (-2.541,-0.976) (-2.506,-0.975) (-2.471,-0.974) (-2.436,-0.973) (-2.400,-0.972) (-2.365,-0.971) (-2.330,-0.970) (-2.295,-0.969) (-2.259,-0.968) (-2.224,-0.967) (-2.189,-0.965) (-2.154,-0.964) (-2.118,-0.963) (-2.083,-0.961) (-2.048,-0.959) (-2.013,-0.958) (-1.977,-0.956) (-1.942,-0.954) (-1.907,-0.952) (-1.872,-0.950) (-1.836,-0.948) (-1.801,-0.945) (-1.766,-0.943) (-1.731,-0.940) (-1.696,-0.938) (-1.661,-0.935) (-1.625,-0.932) (-1.590,-0.928) (-1.555,-0.925) (-1.520,-0.921) (-1.485,-0.917) (-1.450,-0.913) (-1.415,-0.909) (-1.380,-0.904) (-1.345,-0.899) (-1.310,-0.894) (-1.275,-0.889) (-1.240,-0.883) (-1.206,-0.876) (-1.171,-0.870) (-1.137,-0.862) (-1.102,-0.855) (-1.068,-0.847) (-1.034,-0.838) (-1.000,-0.828) (-0.966,-0.818) (-0.932,-0.808) (-0.899,-0.796) (-0.866,-0.784) (-0.833,-0.771) (-0.801,-0.756) (-0.769,-0.741) (-0.738,-0.725) (-0.708,-0.707)};
\addplot[lineaE, flechaen=0.662] coordinates {(0.531,-0.848) (0.559,-0.869) (0.589,-0.888) (0.620,-0.905) (0.651,-0.921) (0.683,-0.936) (0.716,-0.951) (0.749,-0.964) (0.782,-0.976) (0.816,-0.988) (0.849,-0.999) (0.883,-1.009) (0.917,-1.018) (0.952,-1.027) (0.986,-1.036) (1.021,-1.044) (1.055,-1.051) (1.090,-1.059) (1.125,-1.065) (1.160,-1.072) (1.195,-1.078) (1.230,-1.083) (1.265,-1.089) (1.300,-1.094) (1.335,-1.098) (1.371,-1.103) (1.406,-1.107) (1.441,-1.111) (1.476,-1.115) (1.512,-1.119) (1.547,-1.122) (1.582,-1.126) (1.618,-1.129) (1.653,-1.132) (1.688,-1.134) (1.724,-1.137) (1.759,-1.140) (1.794,-1.142) (1.830,-1.144) (1.865,-1.147) (1.901,-1.149) (1.936,-1.151) (1.972,-1.153) (2.007,-1.154) (2.042,-1.156) (2.078,-1.158) (2.113,-1.159) (2.149,-1.161) (2.184,-1.162) (2.220,-1.164) (2.255,-1.165) (2.291,-1.166) (2.326,-1.167) (2.361,-1.168) (2.397,-1.170) (2.432,-1.171) (2.468,-1.172) (2.503,-1.173) (2.539,-1.173) (2.574,-1.174) (2.610,-1.175) (2.645,-1.176) (2.681,-1.177) (2.716,-1.178) (2.752,-1.178) (2.787,-1.179) (2.823,-1.180) (2.858,-1.180) (2.894,-1.181) (2.929,-1.181) (2.965,-1.182) (3.000,-1.182)};
\addplot[lineaE, flechaen=0.338] coordinates {(-3.000,-1.182) (-2.965,-1.182) (-2.929,-1.181) (-2.894,-1.181) (-2.858,-1.180) (-2.823,-1.180) (-2.787,-1.179) (-2.752,-1.178) (-2.716,-1.178) (-2.681,-1.177) (-2.645,-1.176) (-2.610,-1.175) (-2.574,-1.174) (-2.539,-1.173) (-2.503,-1.173) (-2.468,-1.172) (-2.432,-1.171) (-2.397,-1.170) (-2.361,-1.168) (-2.326,-1.167) (-2.291,-1.166) (-2.255,-1.165) (-2.220,-1.164) (-2.184,-1.162) (-2.149,-1.161) (-2.113,-1.159) (-2.078,-1.158) (-2.042,-1.156) (-2.007,-1.154) (-1.972,-1.153) (-1.936,-1.151) (-1.901,-1.149) (-1.865,-1.147) (-1.830,-1.144) (-1.794,-1.142) (-1.759,-1.140) (-1.724,-1.137) (-1.688,-1.134) (-1.653,-1.132) (-1.618,-1.129) (-1.582,-1.126) (-1.547,-1.122) (-1.512,-1.119) (-1.476,-1.115) (-1.441,-1.111) (-1.406,-1.107) (-1.371,-1.103) (-1.335,-1.098) (-1.300,-1.094) (-1.265,-1.089) (-1.230,-1.083) (-1.195,-1.078) (-1.160,-1.072) (-1.125,-1.065) (-1.090,-1.059) (-1.055,-1.051) (-1.021,-1.044) (-0.986,-1.036) (-0.952,-1.027) (-0.917,-1.018) (-0.883,-1.009) (-0.849,-0.999) (-0.816,-0.988) (-0.782,-0.976) (-0.749,-0.964) (-0.716,-0.951) (-0.683,-0.936) (-0.651,-0.921) (-0.620,-0.905) (-0.589,-0.888) (-0.559,-0.869) (-0.531,-0.848)};
\addplot[lineaE, flechaen=0.707] coordinates {(0.148,-0.989) (0.178,-1.005) (0.210,-1.019) (0.243,-1.033) (0.275,-1.047) (0.308,-1.061) (0.340,-1.074) (0.373,-1.087) (0.406,-1.100) (0.440,-1.112) (0.473,-1.124) (0.507,-1.135) (0.540,-1.146) (0.574,-1.157) (0.608,-1.167) (0.642,-1.177) (0.676,-1.186) (0.710,-1.195) (0.745,-1.204) (0.779,-1.212) (0.814,-1.220) (0.848,-1.227) (0.883,-1.235) (0.918,-1.241) (0.953,-1.248) (0.987,-1.254) (1.022,-1.260) (1.057,-1.266) (1.092,-1.271) (1.127,-1.277) (1.162,-1.282) (1.197,-1.287) (1.232,-1.291) (1.268,-1.295) (1.303,-1.300) (1.338,-1.304) (1.373,-1.307) (1.408,-1.311) (1.444,-1.314) (1.479,-1.318) (1.514,-1.321) (1.549,-1.324) (1.585,-1.327) (1.620,-1.330) (1.655,-1.332) (1.691,-1.335) (1.726,-1.337) (1.761,-1.340) (1.797,-1.342) (1.832,-1.344) (1.867,-1.346) (1.903,-1.348) (1.938,-1.350) (1.974,-1.351) (2.009,-1.353) (2.044,-1.355) (2.080,-1.356) (2.115,-1.358) (2.150,-1.359) (2.186,-1.361) (2.221,-1.362) (2.257,-1.363) (2.292,-1.365) (2.327,-1.366) (2.363,-1.367) (2.398,-1.368) (2.434,-1.369) (2.469,-1.370) (2.504,-1.371) (2.540,-1.372) (2.575,-1.373) (2.611,-1.373) (2.646,-1.374) (2.681,-1.375) (2.717,-1.376) (2.752,-1.377) (2.788,-1.377) (2.823,-1.378) (2.858,-1.379) (2.894,-1.379) (2.929,-1.380) (2.965,-1.380) (3.000,-1.381)};
\addplot[lineaE, flechaen=0.293] coordinates {(-3.000,-1.381) (-2.965,-1.380) (-2.929,-1.380) (-2.894,-1.379) (-2.858,-1.379) (-2.823,-1.378) (-2.788,-1.377) (-2.752,-1.377) (-2.717,-1.376) (-2.681,-1.375) (-2.646,-1.374) (-2.611,-1.373) (-2.575,-1.373) (-2.540,-1.372) (-2.504,-1.371) (-2.469,-1.370) (-2.434,-1.369) (-2.398,-1.368) (-2.363,-1.367) (-2.327,-1.366) (-2.292,-1.365) (-2.257,-1.363) (-2.221,-1.362) (-2.186,-1.361) (-2.150,-1.359) (-2.115,-1.358) (-2.080,-1.356) (-2.044,-1.355) (-2.009,-1.353) (-1.974,-1.351) (-1.938,-1.350) (-1.903,-1.348) (-1.867,-1.346) (-1.832,-1.344) (-1.797,-1.342) (-1.761,-1.340) (-1.726,-1.337) (-1.691,-1.335) (-1.655,-1.332) (-1.620,-1.330) (-1.585,-1.327) (-1.549,-1.324) (-1.514,-1.321) (-1.479,-1.318) (-1.444,-1.314) (-1.408,-1.311) (-1.373,-1.307) (-1.338,-1.304) (-1.303,-1.300) (-1.268,-1.295) (-1.232,-1.291) (-1.197,-1.287) (-1.162,-1.282) (-1.127,-1.277) (-1.092,-1.271) (-1.057,-1.266) (-1.022,-1.260) (-0.987,-1.254) (-0.953,-1.248) (-0.918,-1.241) (-0.883,-1.235) (-0.848,-1.227) (-0.814,-1.220) (-0.779,-1.212) (-0.745,-1.204) (-0.710,-1.195) (-0.676,-1.186) (-0.642,-1.177) (-0.608,-1.167) (-0.574,-1.157) (-0.540,-1.146) (-0.507,-1.135) (-0.473,-1.124) (-0.440,-1.112) (-0.406,-1.100) (-0.373,-1.087) (-0.340,-1.074) (-0.308,-1.061) (-0.275,-1.047) (-0.243,-1.033) (-0.210,-1.019) (-0.178,-1.005) (-0.148,-0.989)};
\addplot[lineaE, flechaen=0.141, flechaen=0.859] coordinates {(-3.000,-1.580) (-2.965,-1.579) (-2.930,-1.579) (-2.895,-1.578) (-2.860,-1.578) (-2.825,-1.577) (-2.790,-1.576) (-2.755,-1.576) (-2.720,-1.575) (-2.685,-1.574) (-2.650,-1.573) (-2.615,-1.573) (-2.580,-1.572) (-2.545,-1.571) (-2.510,-1.570) (-2.475,-1.569) (-2.440,-1.568) (-2.405,-1.567) (-2.370,-1.566) (-2.335,-1.565) (-2.300,-1.564) (-2.265,-1.563) (-2.230,-1.562) (-2.195,-1.561) (-2.160,-1.559) (-2.125,-1.558) (-2.090,-1.557) (-2.055,-1.555) (-2.020,-1.554) (-1.985,-1.552) (-1.950,-1.551) (-1.915,-1.549) (-1.880,-1.547) (-1.845,-1.546) (-1.810,-1.544) (-1.775,-1.542) (-1.740,-1.540) (-1.705,-1.538) (-1.670,-1.536) (-1.635,-1.533) (-1.600,-1.531) (-1.565,-1.529) (-1.530,-1.526) (-1.495,-1.524) (-1.460,-1.521) (-1.425,-1.518) (-1.390,-1.515) (-1.355,-1.512) (-1.320,-1.509) (-1.286,-1.506) (-1.251,-1.503) (-1.216,-1.499) (-1.181,-1.496) (-1.146,-1.492) (-1.111,-1.488) (-1.077,-1.484) (-1.042,-1.480) (-1.007,-1.476) (-0.972,-1.471) (-0.937,-1.467) (-0.903,-1.462) (-0.868,-1.458) (-0.833,-1.453) (-0.799,-1.448) (-0.764,-1.443) (-0.729,-1.438) (-0.695,-1.432) (-0.660,-1.427) (-0.626,-1.421) (-0.591,-1.416) (-0.556,-1.410) (-0.522,-1.405) (-0.487,-1.399) (-0.453,-1.394) (-0.418,-1.388) (-0.383,-1.383) (-0.349,-1.378) (-0.314,-1.373) (-0.279,-1.368) (-0.245,-1.364) (-0.210,-1.360) (-0.175,-1.357) (-0.140,-1.354) (-0.105,-1.352) (-0.070,-1.350) (-0.035,-1.349) (0.000,-1.348) (0.035,-1.349) (0.070,-1.350) (0.105,-1.352) (0.140,-1.354) (0.175,-1.357) (0.210,-1.360) (0.245,-1.364) (0.279,-1.368) (0.314,-1.373) (0.349,-1.378) (0.383,-1.383) (0.418,-1.388) (0.453,-1.394) (0.487,-1.399) (0.522,-1.405) (0.556,-1.410) (0.591,-1.416) (0.626,-1.421) (0.660,-1.427) (0.695,-1.432) (0.729,-1.438) (0.764,-1.443) (0.799,-1.448) (0.833,-1.453) (0.868,-1.458) (0.903,-1.462) (0.937,-1.467) (0.972,-1.471) (1.007,-1.476) (1.042,-1.480) (1.077,-1.484) (1.111,-1.488) (1.146,-1.492) (1.181,-1.496) (1.216,-1.499) (1.251,-1.503) (1.286,-1.506) (1.320,-1.509) (1.355,-1.512) (1.390,-1.515) (1.425,-1.518) (1.460,-1.521) (1.495,-1.524) (1.530,-1.526) (1.565,-1.529) (1.600,-1.531) (1.635,-1.533) (1.670,-1.536) (1.705,-1.538) (1.740,-1.540) (1.775,-1.542) (1.810,-1.544) (1.845,-1.546) (1.880,-1.547) (1.915,-1.549) (1.950,-1.551) (1.985,-1.552) (2.020,-1.554) (2.055,-1.555) (2.090,-1.557) (2.125,-1.558) (2.160,-1.559) (2.195,-1.561) (2.230,-1.562) (2.265,-1.563) (2.300,-1.564) (2.335,-1.565) (2.370,-1.566) (2.405,-1.567) (2.440,-1.568) (2.475,-1.569) (2.510,-1.570) (2.545,-1.571) (2.580,-1.572) (2.615,-1.573) (2.650,-1.573) (2.685,-1.574) (2.720,-1.575) (2.755,-1.576) (2.790,-1.576) (2.825,-1.577) (2.860,-1.578) (2.895,-1.578) (2.930,-1.579) (2.965,-1.579) (3.000,-1.580)};
\addplot[lineaE, flechaen=0.141, flechaen=0.859] coordinates {(-3.000,-1.779) (-2.965,-1.779) (-2.930,-1.778) (-2.895,-1.777) (-2.859,-1.777) (-2.824,-1.776) (-2.789,-1.776) (-2.754,-1.775) (-2.719,-1.774) (-2.684,-1.774) (-2.649,-1.773) (-2.613,-1.772) (-2.578,-1.771) (-2.543,-1.770) (-2.508,-1.770) (-2.473,-1.769) (-2.438,-1.768) (-2.403,-1.767) (-2.367,-1.766) (-2.332,-1.765) (-2.297,-1.764) (-2.262,-1.763) (-2.227,-1.762) (-2.192,-1.761) (-2.157,-1.760) (-2.121,-1.758) (-2.086,-1.757) (-2.051,-1.756) (-2.016,-1.755) (-1.981,-1.753) (-1.946,-1.752) (-1.911,-1.750) (-1.876,-1.749) (-1.841,-1.747) (-1.805,-1.746) (-1.770,-1.744) (-1.735,-1.742) (-1.700,-1.741) (-1.665,-1.739) (-1.630,-1.737) (-1.595,-1.735) (-1.560,-1.733) (-1.525,-1.731) (-1.489,-1.729) (-1.454,-1.727) (-1.419,-1.724) (-1.384,-1.722) (-1.349,-1.720) (-1.314,-1.717) (-1.279,-1.715) (-1.244,-1.712) (-1.209,-1.709) (-1.174,-1.707) (-1.139,-1.704) (-1.104,-1.701) (-1.069,-1.698) (-1.034,-1.695) (-0.999,-1.692) (-0.964,-1.689) (-0.929,-1.686) (-0.894,-1.682) (-0.859,-1.679) (-0.824,-1.676) (-0.789,-1.673) (-0.754,-1.669) (-0.719,-1.666) (-0.684,-1.663) (-0.649,-1.659) (-0.614,-1.656) (-0.579,-1.653) (-0.544,-1.649) (-0.509,-1.646) (-0.474,-1.643) (-0.439,-1.640) (-0.404,-1.637) (-0.369,-1.635) (-0.334,-1.632) (-0.299,-1.630) (-0.263,-1.628) (-0.228,-1.626) (-0.193,-1.624) (-0.158,-1.622) (-0.123,-1.621) (-0.088,-1.620) (-0.053,-1.620) (-0.018,-1.619) (0.018,-1.619) (0.053,-1.620) (0.088,-1.620) (0.123,-1.621) (0.158,-1.622) (0.193,-1.624) (0.228,-1.626) (0.263,-1.628) (0.299,-1.630) (0.334,-1.632) (0.369,-1.635) (0.404,-1.637) (0.439,-1.640) (0.474,-1.643) (0.509,-1.646) (0.544,-1.649) (0.579,-1.653) (0.614,-1.656) (0.649,-1.659) (0.684,-1.663) (0.719,-1.666) (0.754,-1.669) (0.789,-1.673) (0.824,-1.676) (0.859,-1.679) (0.894,-1.682) (0.929,-1.686) (0.964,-1.689) (0.999,-1.692) (1.034,-1.695) (1.069,-1.698) (1.104,-1.701) (1.139,-1.704) (1.174,-1.707) (1.209,-1.709) (1.244,-1.712) (1.279,-1.715) (1.314,-1.717) (1.349,-1.720) (1.384,-1.722) (1.419,-1.724) (1.454,-1.727) (1.489,-1.729) (1.525,-1.731) (1.560,-1.733) (1.595,-1.735) (1.630,-1.737) (1.665,-1.739) (1.700,-1.741) (1.735,-1.742) (1.770,-1.744) (1.805,-1.746) (1.841,-1.747) (1.876,-1.749) (1.911,-1.750) (1.946,-1.752) (1.981,-1.753) (2.016,-1.755) (2.051,-1.756) (2.086,-1.757) (2.121,-1.758) (2.157,-1.760) (2.192,-1.761) (2.227,-1.762) (2.262,-1.763) (2.297,-1.764) (2.332,-1.765) (2.367,-1.766) (2.403,-1.767) (2.438,-1.768) (2.473,-1.769) (2.508,-1.770) (2.543,-1.770) (2.578,-1.771) (2.613,-1.772) (2.649,-1.773) (2.684,-1.774) (2.719,-1.774) (2.754,-1.775) (2.789,-1.776) (2.824,-1.776) (2.859,-1.777) (2.895,-1.777) (2.930,-1.778) (2.965,-1.779) (3.000,-1.779)};
\addplot[lineaE, flechaen=0.142, flechaen=0.858] coordinates {(-3.000,-1.979) (-2.965,-1.978) (-2.930,-1.978) (-2.895,-1.977) (-2.860,-1.977) (-2.824,-1.976) (-2.789,-1.975) (-2.754,-1.975) (-2.719,-1.974) (-2.684,-1.973) (-2.649,-1.973) (-2.614,-1.972) (-2.579,-1.971) (-2.544,-1.971) (-2.508,-1.970) (-2.473,-1.969) (-2.438,-1.968) (-2.403,-1.967) (-2.368,-1.967) (-2.333,-1.966) (-2.298,-1.965) (-2.263,-1.964) (-2.228,-1.963) (-2.192,-1.962) (-2.157,-1.961) (-2.122,-1.960) (-2.087,-1.959) (-2.052,-1.958) (-2.017,-1.956) (-1.982,-1.955) (-1.947,-1.954) (-1.912,-1.953) (-1.877,-1.951) (-1.841,-1.950) (-1.806,-1.949) (-1.771,-1.947) (-1.736,-1.946) (-1.701,-1.944) (-1.666,-1.943) (-1.631,-1.941) (-1.596,-1.940) (-1.561,-1.938) (-1.526,-1.936) (-1.491,-1.935) (-1.455,-1.933) (-1.420,-1.931) (-1.385,-1.929) (-1.350,-1.927) (-1.315,-1.926) (-1.280,-1.924) (-1.245,-1.922) (-1.210,-1.920) (-1.175,-1.917) (-1.140,-1.915) (-1.105,-1.913) (-1.070,-1.911) (-1.035,-1.909) (-1.000,-1.907) (-0.965,-1.905) (-0.930,-1.902) (-0.895,-1.900) (-0.859,-1.898) (-0.824,-1.896) (-0.789,-1.893) (-0.754,-1.891) (-0.719,-1.889) (-0.684,-1.887) (-0.649,-1.885) (-0.614,-1.882) (-0.579,-1.880) (-0.544,-1.878) (-0.509,-1.876) (-0.474,-1.875) (-0.439,-1.873) (-0.404,-1.871) (-0.369,-1.869) (-0.334,-1.868) (-0.298,-1.867) (-0.263,-1.865) (-0.228,-1.864) (-0.193,-1.863) (-0.158,-1.862) (-0.123,-1.862) (-0.088,-1.861) (-0.053,-1.861) (-0.018,-1.861) (0.018,-1.861) (0.053,-1.861) (0.088,-1.861) (0.123,-1.862) (0.158,-1.862) (0.193,-1.863) (0.228,-1.864) (0.263,-1.865) (0.298,-1.867) (0.334,-1.868) (0.369,-1.869) (0.404,-1.871) (0.439,-1.873) (0.474,-1.875) (0.509,-1.876) (0.544,-1.878) (0.579,-1.880) (0.614,-1.882) (0.649,-1.885) (0.684,-1.887) (0.719,-1.889) (0.754,-1.891) (0.789,-1.893) (0.824,-1.896) (0.859,-1.898) (0.895,-1.900) (0.930,-1.902) (0.965,-1.905) (1.000,-1.907) (1.035,-1.909) (1.070,-1.911) (1.105,-1.913) (1.140,-1.915) (1.175,-1.917) (1.210,-1.920) (1.245,-1.922) (1.280,-1.924) (1.315,-1.926) (1.350,-1.927) (1.385,-1.929) (1.420,-1.931) (1.455,-1.933) (1.491,-1.935) (1.526,-1.936) (1.561,-1.938) (1.596,-1.940) (1.631,-1.941) (1.666,-1.943) (1.701,-1.944) (1.736,-1.946) (1.771,-1.947) (1.806,-1.949) (1.841,-1.950) (1.877,-1.951) (1.912,-1.953) (1.947,-1.954) (1.982,-1.955) (2.017,-1.956) (2.052,-1.958) (2.087,-1.959) (2.122,-1.960) (2.157,-1.961) (2.192,-1.962) (2.228,-1.963) (2.263,-1.964) (2.298,-1.965) (2.333,-1.966) (2.368,-1.967) (2.403,-1.967) (2.438,-1.968) (2.473,-1.969) (2.508,-1.970) (2.544,-1.971) (2.579,-1.971) (2.614,-1.972) (2.649,-1.973) (2.684,-1.973) (2.719,-1.974) (2.754,-1.975) (2.789,-1.975) (2.824,-1.976) (2.860,-1.977) (2.895,-1.977) (2.930,-1.978) (2.965,-1.978) (3.000,-1.979)};
\addplot[lineaE, flechaen=0.142, flechaen=0.858] coordinates {(-3.000,-2.179) (-2.965,-2.178) (-2.930,-2.178) (-2.895,-2.177) (-2.860,-2.177) (-2.824,-2.176) (-2.789,-2.176) (-2.754,-2.175) (-2.719,-2.174) (-2.684,-2.174) (-2.649,-2.173) (-2.614,-2.173) (-2.579,-2.172) (-2.544,-2.171) (-2.509,-2.170) (-2.473,-2.170) (-2.438,-2.169) (-2.403,-2.168) (-2.368,-2.167) (-2.333,-2.167) (-2.298,-2.166) (-2.263,-2.165) (-2.228,-2.164) (-2.193,-2.163) (-2.158,-2.162) (-2.123,-2.161) (-2.087,-2.161) (-2.052,-2.160) (-2.017,-2.159) (-1.982,-2.158) (-1.947,-2.156) (-1.912,-2.155) (-1.877,-2.154) (-1.842,-2.153) (-1.807,-2.152) (-1.772,-2.151) (-1.737,-2.150) (-1.701,-2.148) (-1.666,-2.147) (-1.631,-2.146) (-1.596,-2.145) (-1.561,-2.143) (-1.526,-2.142) (-1.491,-2.141) (-1.456,-2.139) (-1.421,-2.138) (-1.386,-2.136) (-1.351,-2.135) (-1.316,-2.133) (-1.281,-2.132) (-1.245,-2.130) (-1.210,-2.129) (-1.175,-2.127) (-1.140,-2.126) (-1.105,-2.124) (-1.070,-2.122) (-1.035,-2.121) (-1.000,-2.119) (-0.965,-2.118) (-0.930,-2.116) (-0.895,-2.114) (-0.860,-2.113) (-0.825,-2.111) (-0.790,-2.110) (-0.754,-2.108) (-0.719,-2.107) (-0.684,-2.105) (-0.649,-2.104) (-0.614,-2.102) (-0.579,-2.101) (-0.544,-2.100) (-0.509,-2.098) (-0.474,-2.097) (-0.439,-2.096) (-0.404,-2.095) (-0.369,-2.094) (-0.333,-2.093) (-0.298,-2.092) (-0.263,-2.091) (-0.228,-2.090) (-0.193,-2.090) (-0.158,-2.089) (-0.123,-2.089) (-0.088,-2.089) (-0.053,-2.088) (-0.018,-2.088) (0.018,-2.088) (0.053,-2.088) (0.088,-2.089) (0.123,-2.089) (0.158,-2.089) (0.193,-2.090) (0.228,-2.090) (0.263,-2.091) (0.298,-2.092) (0.333,-2.093) (0.369,-2.094) (0.404,-2.095) (0.439,-2.096) (0.474,-2.097) (0.509,-2.098) (0.544,-2.100) (0.579,-2.101) (0.614,-2.102) (0.649,-2.104) (0.684,-2.105) (0.719,-2.107) (0.754,-2.108) (0.790,-2.110) (0.825,-2.111) (0.860,-2.113) (0.895,-2.114) (0.930,-2.116) (0.965,-2.118) (1.000,-2.119) (1.035,-2.121) (1.070,-2.122) (1.105,-2.124) (1.140,-2.126) (1.175,-2.127) (1.210,-2.129) (1.245,-2.130) (1.281,-2.132) (1.316,-2.133) (1.351,-2.135) (1.386,-2.136) (1.421,-2.138) (1.456,-2.139) (1.491,-2.141) (1.526,-2.142) (1.561,-2.143) (1.596,-2.145) (1.631,-2.146) (1.666,-2.147) (1.701,-2.148) (1.737,-2.150) (1.772,-2.151) (1.807,-2.152) (1.842,-2.153) (1.877,-2.154) (1.912,-2.155) (1.947,-2.156) (1.982,-2.158) (2.017,-2.159) (2.052,-2.160) (2.087,-2.161) (2.123,-2.161) (2.158,-2.162) (2.193,-2.163) (2.228,-2.164) (2.263,-2.165) (2.298,-2.166) (2.333,-2.167) (2.368,-2.167) (2.403,-2.168) (2.438,-2.169) (2.473,-2.170) (2.509,-2.170) (2.544,-2.171) (2.579,-2.172) (2.614,-2.173) (2.649,-2.173) (2.684,-2.174) (2.719,-2.174) (2.754,-2.175) (2.789,-2.176) (2.824,-2.176) (2.860,-2.177) (2.895,-2.177) (2.930,-2.178) (2.965,-2.178) (3.000,-2.179)};
\addplot[lineaE, flechaen=0.239] coordinates {(-0.959,-0.283) (-0.924,-0.283) (-0.888,-0.283) (-0.853,-0.283) (-0.817,-0.283) (-0.781,-0.283) (-0.746,-0.283) (-0.710,-0.283) (-0.675,-0.283) (-0.639,-0.283) (-0.604,-0.283) (-0.568,-0.283) (-0.533,-0.283) (-0.497,-0.283) (-0.462,-0.283) (-0.426,-0.283) (-0.391,-0.283) (-0.355,-0.283) (-0.320,-0.283) (-0.284,-0.283) (-0.249,-0.283) (-0.213,-0.283) (-0.178,-0.283) (-0.142,-0.283) (-0.107,-0.283) (-0.071,-0.283) (-0.036,-0.283) (-0.000,-0.283) (0.036,-0.283) (0.071,-0.283) (0.107,-0.283) (0.142,-0.283) (0.178,-0.283) (0.213,-0.283) (0.249,-0.283) (0.284,-0.283) (0.320,-0.283) (0.355,-0.283) (0.391,-0.283) (0.426,-0.283) (0.462,-0.283) (0.497,-0.283) (0.533,-0.283) (0.568,-0.283) (0.604,-0.283) (0.639,-0.283) (0.675,-0.283) (0.710,-0.283) (0.746,-0.283) (0.781,-0.283) (0.817,-0.283) (0.853,-0.283) (0.888,-0.283) (0.924,-0.283) (0.959,-0.283)};
\addplot[lineaE, flechaen=0.804] coordinates {(-0.824,-0.567) (-0.789,-0.566) (-0.754,-0.566) (-0.719,-0.566) (-0.684,-0.566) (-0.649,-0.566) (-0.614,-0.566) (-0.579,-0.566) (-0.544,-0.566) (-0.508,-0.566) (-0.473,-0.566) (-0.438,-0.566) (-0.403,-0.566) (-0.368,-0.566) (-0.333,-0.566) (-0.298,-0.566) (-0.263,-0.566) (-0.228,-0.566) (-0.193,-0.566) (-0.158,-0.566) (-0.123,-0.566) (-0.088,-0.566) (-0.053,-0.566) (-0.018,-0.566) (0.018,-0.566) (0.053,-0.566) (0.088,-0.566) (0.123,-0.566) (0.158,-0.566) (0.193,-0.566) (0.228,-0.566) (0.263,-0.566) (0.298,-0.566) (0.333,-0.566) (0.368,-0.566) (0.403,-0.566) (0.438,-0.566) (0.473,-0.566) (0.508,-0.566) (0.544,-0.566) (0.579,-0.566) (0.614,-0.566) (0.649,-0.566) (0.684,-0.566) (0.719,-0.566) (0.754,-0.566) (0.789,-0.566) (0.824,-0.567)};
\addplot[lineaE, flechaen=0.500] coordinates {(-0.525,-0.851) (-0.491,-0.849) (-0.456,-0.849) (-0.421,-0.849) (-0.386,-0.849) (-0.351,-0.849) (-0.316,-0.849) (-0.281,-0.849) (-0.246,-0.849) (-0.210,-0.849) (-0.175,-0.849) (-0.140,-0.849) (-0.105,-0.849) (-0.070,-0.849) (-0.035,-0.849) (0.000,-0.849) (0.035,-0.849) (0.070,-0.849) (0.105,-0.849) (0.140,-0.849) (0.175,-0.849) (0.210,-0.849) (0.246,-0.849) (0.281,-0.849) (0.316,-0.849) (0.351,-0.849) (0.386,-0.849) (0.421,-0.849) (0.456,-0.849) (0.491,-0.849) (0.525,-0.851)};
\addplot[lineaE, flechaen=0.142, flechaen=0.858] coordinates {(-3.000,0.000) (3.000,0.000)};

  \draw[texto, line width=0.6pt] (axis cs:0,0) circle[radius=1];
  \pgfplotsinvokeforeach{-60,-30,0,30,60}{
    \node[cpos, font=\scriptsize\bfseries, fill=fondo, circle, inner sep=0.2pt] at (axis cs:{1.16*cos(#1)},{1.16*sin(#1)}) {$+$};
    \node[cneg, font=\scriptsize\bfseries, fill=fondo, circle, inner sep=0.2pt] at (axis cs:{-1.16*cos(#1)},{1.16*sin(#1)}) {$-$};
  }
\end{axis}
\end{tikzpicture}
\caption{Líneas de $\vE$}
\end{subfigure}\hfill
\begin{subfigure}[b]{0.49\textwidth}\centering
\begin{tikzpicture}
\begin{axis}[esfera]
  \fill[dielec] (axis cs:0,0) circle[radius=1];
  %% Generado por scripts/lineas_campo.py -- no editar a mano
\addplot[equipot] coordinates {(2.781,-2.200) (2.782,-2.165) (2.782,-2.130) (2.783,-2.094) (2.784,-2.059) (2.784,-2.024) (2.785,-1.989) (2.785,-1.954) (2.786,-1.918) (2.787,-1.883) (2.787,-1.848) (2.788,-1.813) (2.789,-1.778) (2.789,-1.742) (2.790,-1.707) (2.791,-1.672) (2.791,-1.637) (2.792,-1.602) (2.793,-1.566) (2.793,-1.531) (2.794,-1.496) (2.795,-1.461) (2.795,-1.426) (2.796,-1.390) (2.797,-1.355) (2.797,-1.320) (2.798,-1.285) (2.799,-1.250) (2.799,-1.214) (2.800,-1.179) (2.801,-1.144) (2.801,-1.109) (2.802,-1.074) (2.802,-1.038) (2.803,-1.003) (2.804,-0.968) (2.804,-0.933) (2.805,-0.898) (2.805,-0.862) (2.806,-0.827) (2.807,-0.792) (2.807,-0.757) (2.808,-0.722) (2.808,-0.686) (2.809,-0.651) (2.809,-0.616) (2.809,-0.581) (2.810,-0.546) (2.810,-0.510) (2.811,-0.475) (2.811,-0.440) (2.811,-0.405) (2.812,-0.370) (2.812,-0.334) (2.812,-0.299) (2.812,-0.264) (2.813,-0.229) (2.813,-0.194) (2.813,-0.158) (2.813,-0.123) (2.813,-0.088) (2.813,-0.053) (2.813,-0.018) (2.813,0.018) (2.813,0.053) (2.813,0.088) (2.813,0.123) (2.813,0.158) (2.813,0.194) (2.813,0.229) (2.812,0.264) (2.812,0.299) (2.812,0.334) (2.812,0.370) (2.811,0.405) (2.811,0.440) (2.811,0.475) (2.810,0.510) (2.810,0.546) (2.809,0.581) (2.809,0.616) (2.809,0.651) (2.808,0.686) (2.808,0.722) (2.807,0.757) (2.807,0.792) (2.806,0.827) (2.805,0.862) (2.805,0.898) (2.804,0.933) (2.804,0.968) (2.803,1.003) (2.802,1.038) (2.802,1.074) (2.801,1.109) (2.801,1.144) (2.800,1.179) (2.799,1.214) (2.799,1.250) (2.798,1.285) (2.797,1.320) (2.797,1.355) (2.796,1.390) (2.795,1.426) (2.795,1.461) (2.794,1.496) (2.793,1.531) (2.793,1.566) (2.792,1.602) (2.791,1.637) (2.791,1.672) (2.790,1.707) (2.789,1.742) (2.789,1.778) (2.788,1.813) (2.787,1.848) (2.787,1.883) (2.786,1.918) (2.785,1.954) (2.785,1.989) (2.784,2.024) (2.784,2.059) (2.783,2.094) (2.782,2.130) (2.782,2.165) (2.781,2.200)};
\addplot[equipot] coordinates {(2.286,-2.200) (2.287,-2.165) (2.287,-2.130) (2.288,-2.094) (2.289,-2.059) (2.290,-2.024) (2.291,-1.989) (2.292,-1.954) (2.293,-1.918) (2.294,-1.883) (2.295,-1.848) (2.296,-1.813) (2.297,-1.778) (2.298,-1.742) (2.299,-1.707) (2.300,-1.672) (2.301,-1.637) (2.302,-1.602) (2.303,-1.566) (2.304,-1.531) (2.306,-1.496) (2.307,-1.461) (2.308,-1.426) (2.309,-1.390) (2.310,-1.355) (2.311,-1.320) (2.312,-1.285) (2.314,-1.250) (2.315,-1.214) (2.316,-1.179) (2.317,-1.144) (2.318,-1.109) (2.319,-1.074) (2.321,-1.039) (2.322,-1.003) (2.323,-0.968) (2.324,-0.933) (2.325,-0.898) (2.326,-0.863) (2.327,-0.827) (2.328,-0.792) (2.329,-0.757) (2.330,-0.722) (2.331,-0.687) (2.332,-0.651) (2.333,-0.616) (2.334,-0.581) (2.335,-0.546) (2.335,-0.511) (2.336,-0.475) (2.337,-0.440) (2.338,-0.405) (2.338,-0.370) (2.339,-0.335) (2.339,-0.299) (2.340,-0.264) (2.340,-0.229) (2.340,-0.194) (2.341,-0.158) (2.341,-0.123) (2.341,-0.088) (2.341,-0.053) (2.341,-0.018) (2.341,0.018) (2.341,0.053) (2.341,0.088) (2.341,0.123) (2.341,0.158) (2.340,0.194) (2.340,0.229) (2.340,0.264) (2.339,0.299) (2.339,0.335) (2.338,0.370) (2.338,0.405) (2.337,0.440) (2.336,0.475) (2.335,0.511) (2.335,0.546) (2.334,0.581) (2.333,0.616) (2.332,0.651) (2.331,0.687) (2.330,0.722) (2.329,0.757) (2.328,0.792) (2.327,0.827) (2.326,0.863) (2.325,0.898) (2.324,0.933) (2.323,0.968) (2.322,1.003) (2.321,1.039) (2.319,1.074) (2.318,1.109) (2.317,1.144) (2.316,1.179) (2.315,1.214) (2.314,1.250) (2.312,1.285) (2.311,1.320) (2.310,1.355) (2.309,1.390) (2.308,1.426) (2.307,1.461) (2.306,1.496) (2.304,1.531) (2.303,1.566) (2.302,1.602) (2.301,1.637) (2.300,1.672) (2.299,1.707) (2.298,1.742) (2.297,1.778) (2.296,1.813) (2.295,1.848) (2.294,1.883) (2.293,1.918) (2.292,1.954) (2.291,1.989) (2.290,2.024) (2.289,2.059) (2.288,2.094) (2.287,2.130) (2.287,2.165) (2.286,2.200)};
\addplot[equipot] coordinates {(1.789,-2.200) (1.790,-2.165) (1.792,-2.130) (1.793,-2.094) (1.794,-2.059) (1.795,-2.024) (1.797,-1.989) (1.798,-1.953) (1.799,-1.918) (1.801,-1.883) (1.802,-1.848) (1.804,-1.813) (1.806,-1.777) (1.807,-1.742) (1.809,-1.707) (1.810,-1.672) (1.812,-1.637) (1.814,-1.601) (1.816,-1.566) (1.818,-1.531) (1.820,-1.496) (1.822,-1.461) (1.824,-1.425) (1.826,-1.390) (1.828,-1.355) (1.830,-1.320) (1.832,-1.285) (1.834,-1.250) (1.836,-1.214) (1.838,-1.179) (1.840,-1.144) (1.843,-1.109) (1.845,-1.074) (1.847,-1.039) (1.849,-1.003) (1.851,-0.968) (1.854,-0.933) (1.856,-0.898) (1.858,-0.863) (1.860,-0.827) (1.862,-0.792) (1.864,-0.757) (1.866,-0.722) (1.868,-0.687) (1.870,-0.652) (1.872,-0.616) (1.874,-0.581) (1.876,-0.546) (1.877,-0.511) (1.879,-0.476) (1.881,-0.440) (1.882,-0.405) (1.883,-0.370) (1.884,-0.335) (1.885,-0.300) (1.886,-0.264) (1.887,-0.229) (1.888,-0.194) (1.889,-0.159) (1.889,-0.123) (1.890,-0.088) (1.890,-0.053) (1.890,-0.018) (1.890,0.018) (1.890,0.053) (1.890,0.088) (1.889,0.123) (1.889,0.159) (1.888,0.194) (1.887,0.229) (1.886,0.264) (1.885,0.300) (1.884,0.335) (1.883,0.370) (1.882,0.405) (1.881,0.440) (1.879,0.476) (1.877,0.511) (1.876,0.546) (1.874,0.581) (1.872,0.616) (1.870,0.652) (1.868,0.687) (1.866,0.722) (1.864,0.757) (1.862,0.792) (1.860,0.827) (1.858,0.863) (1.856,0.898) (1.854,0.933) (1.851,0.968) (1.849,1.003) (1.847,1.039) (1.845,1.074) (1.843,1.109) (1.840,1.144) (1.838,1.179) (1.836,1.214) (1.834,1.250) (1.832,1.285) (1.830,1.320) (1.828,1.355) (1.826,1.390) (1.824,1.425) (1.822,1.461) (1.820,1.496) (1.818,1.531) (1.816,1.566) (1.814,1.601) (1.812,1.637) (1.810,1.672) (1.809,1.707) (1.807,1.742) (1.806,1.777) (1.804,1.813) (1.802,1.848) (1.801,1.883) (1.799,1.918) (1.798,1.953) (1.797,1.989) (1.795,2.024) (1.794,2.059) (1.793,2.094) (1.792,2.130) (1.790,2.165) (1.789,2.200)};
\addplot[equipot] coordinates {(1.289,-2.200) (1.290,-2.165) (1.292,-2.130) (1.293,-2.095) (1.295,-2.060) (1.297,-2.025) (1.298,-1.990) (1.300,-1.955) (1.302,-1.920) (1.304,-1.885) (1.306,-1.850) (1.308,-1.815) (1.311,-1.780) (1.313,-1.745) (1.316,-1.710) (1.318,-1.675) (1.321,-1.640) (1.324,-1.605) (1.326,-1.570) (1.329,-1.535) (1.332,-1.500) (1.336,-1.465) (1.339,-1.430) (1.342,-1.395) (1.346,-1.360) (1.350,-1.326) (1.353,-1.291) (1.357,-1.256) (1.361,-1.221) (1.365,-1.186) (1.370,-1.151) (1.374,-1.116) (1.378,-1.082) (1.383,-1.047) (1.387,-1.012) (1.392,-0.977) (1.396,-0.943) (1.401,-0.908) (1.405,-0.873) (1.410,-0.838) (1.414,-0.803) (1.419,-0.769) (1.423,-0.734) (1.428,-0.699) (1.432,-0.664) (1.436,-0.629) (1.440,-0.595) (1.444,-0.560) (1.448,-0.525) (1.452,-0.490) (1.455,-0.455) (1.459,-0.420) (1.462,-0.385) (1.464,-0.350) (1.467,-0.315) (1.469,-0.280) (1.472,-0.245) (1.473,-0.210) (1.475,-0.175) (1.476,-0.140) (1.477,-0.105) (1.478,-0.070) (1.479,-0.035) (1.479,0.000) (1.479,0.035) (1.478,0.070) (1.477,0.105) (1.476,0.140) (1.475,0.175) (1.473,0.210) (1.472,0.245) (1.469,0.280) (1.467,0.315) (1.464,0.350) (1.462,0.385) (1.459,0.420) (1.455,0.455) (1.452,0.490) (1.448,0.525) (1.444,0.560) (1.440,0.595) (1.436,0.629) (1.432,0.664) (1.428,0.699) (1.423,0.734) (1.419,0.769) (1.414,0.803) (1.410,0.838) (1.405,0.873) (1.401,0.908) (1.396,0.943) (1.392,0.977) (1.387,1.012) (1.383,1.047) (1.378,1.082) (1.374,1.116) (1.370,1.151) (1.365,1.186) (1.361,1.221) (1.357,1.256) (1.353,1.291) (1.350,1.326) (1.346,1.360) (1.342,1.395) (1.339,1.430) (1.336,1.465) (1.332,1.500) (1.329,1.535) (1.326,1.570) (1.324,1.605) (1.321,1.640) (1.318,1.675) (1.316,1.710) (1.313,1.745) (1.311,1.780) (1.308,1.815) (1.306,1.850) (1.304,1.885) (1.302,1.920) (1.300,1.955) (1.298,1.990) (1.297,2.025) (1.295,2.060) (1.293,2.095) (1.292,2.130) (1.290,2.165) (1.289,2.200)};
\addplot[equipot] coordinates {(0.781,-2.200) (0.782,-2.165) (0.784,-2.130) (0.785,-2.094) (0.787,-2.059) (0.788,-2.024) (0.790,-1.989) (0.792,-1.954) (0.794,-1.918) (0.797,-1.883) (0.799,-1.848) (0.801,-1.813) (0.804,-1.778) (0.807,-1.743) (0.810,-1.708) (0.813,-1.673) (0.817,-1.638) (0.820,-1.602) (0.824,-1.567) (0.828,-1.532) (0.833,-1.498) (0.837,-1.463) (0.842,-1.428) (0.848,-1.393) (0.853,-1.358) (0.859,-1.323) (0.866,-1.289) (0.872,-1.254) (0.879,-1.220) (0.887,-1.185) (0.894,-1.151) (0.903,-1.117) (0.911,-1.082) (0.920,-1.048) (0.929,-1.014) (0.938,-0.980) (0.947,-0.946) (0.957,-0.912) (0.967,-0.878) (0.977,-0.845) (0.987,-0.811) (0.997,-0.777) (1.007,-0.743) (1.017,-0.710) (1.027,-0.676) (1.036,-0.642) (1.046,-0.608) (1.055,-0.574) (1.063,-0.540) (1.072,-0.505) (1.080,-0.471) (1.088,-0.437) (1.095,-0.402) (1.102,-0.368) (1.108,-0.333) (1.113,-0.298) (1.119,-0.263) (1.123,-0.228) (1.127,-0.193) (1.130,-0.158) (1.133,-0.123) (1.135,-0.088) (1.136,-0.053) (1.137,-0.018) (1.137,0.018) (1.136,0.053) (1.135,0.088) (1.133,0.123) (1.130,0.158) (1.127,0.193) (1.123,0.228) (1.119,0.263) (1.113,0.298) (1.108,0.333) (1.102,0.368) (1.095,0.402) (1.088,0.437) (1.080,0.471) (1.072,0.505) (1.063,0.540) (1.055,0.574) (1.046,0.608) (1.036,0.642) (1.027,0.676) (1.017,0.710) (1.007,0.743) (0.997,0.777) (0.987,0.811) (0.977,0.845) (0.967,0.878) (0.957,0.912) (0.947,0.946) (0.938,0.980) (0.929,1.014) (0.920,1.048) (0.911,1.082) (0.903,1.117) (0.894,1.151) (0.887,1.185) (0.879,1.220) (0.872,1.254) (0.866,1.289) (0.859,1.323) (0.853,1.358) (0.848,1.393) (0.842,1.428) (0.837,1.463) (0.833,1.498) (0.828,1.532) (0.824,1.567) (0.820,1.602) (0.817,1.638) (0.813,1.673) (0.810,1.708) (0.807,1.743) (0.804,1.778) (0.801,1.813) (0.799,1.848) (0.797,1.883) (0.794,1.918) (0.792,1.954) (0.790,1.989) (0.788,2.024) (0.787,2.059) (0.785,2.094) (0.784,2.130) (0.782,2.165) (0.781,2.200)};
\addplot[equipot] coordinates {(0.262,-2.200) (0.263,-2.165) (0.263,-2.130) (0.264,-2.095) (0.265,-2.060) (0.266,-2.025) (0.266,-1.990) (0.267,-1.955) (0.268,-1.919) (0.270,-1.884) (0.271,-1.849) (0.272,-1.814) (0.273,-1.779) (0.275,-1.744) (0.277,-1.709) (0.278,-1.674) (0.280,-1.639) (0.283,-1.604) (0.285,-1.569) (0.288,-1.534) (0.291,-1.499) (0.294,-1.464) (0.298,-1.429) (0.302,-1.395) (0.307,-1.360) (0.312,-1.325) (0.318,-1.291) (0.324,-1.256) (0.332,-1.222) (0.340,-1.188) (0.350,-1.154) (0.360,-1.121) (0.373,-1.088) (0.386,-1.055) (0.401,-1.024) (0.417,-0.993) (0.435,-0.962) (0.453,-0.933) (0.473,-0.904) (0.493,-0.875) (0.500,-0.842) (0.500,-0.807) (0.500,-0.772) (0.500,-0.737) (0.500,-0.701) (0.500,-0.666) (0.500,-0.631) (0.500,-0.596) (0.500,-0.561) (0.500,-0.526) (0.500,-0.491) (0.500,-0.456) (0.500,-0.421) (0.500,-0.386) (0.500,-0.351) (0.500,-0.316) (0.500,-0.281) (0.500,-0.246) (0.500,-0.210) (0.500,-0.175) (0.500,-0.140) (0.500,-0.105) (0.500,-0.070) (0.500,-0.035) (0.500,0.000) (0.500,0.035) (0.500,0.070) (0.500,0.105) (0.500,0.140) (0.500,0.175) (0.500,0.210) (0.500,0.246) (0.500,0.281) (0.500,0.316) (0.500,0.351) (0.500,0.386) (0.500,0.421) (0.500,0.456) (0.500,0.491) (0.500,0.526) (0.500,0.561) (0.500,0.596) (0.500,0.631) (0.500,0.666) (0.500,0.701) (0.500,0.737) (0.500,0.772) (0.500,0.807) (0.500,0.842) (0.493,0.875) (0.473,0.904) (0.453,0.933) (0.435,0.962) (0.417,0.993) (0.401,1.024) (0.386,1.055) (0.373,1.088) (0.360,1.121) (0.350,1.154) (0.340,1.188) (0.332,1.222) (0.324,1.256) (0.318,1.291) (0.312,1.325) (0.307,1.360) (0.302,1.395) (0.298,1.429) (0.294,1.464) (0.291,1.499) (0.288,1.534) (0.285,1.569) (0.283,1.604) (0.280,1.639) (0.278,1.674) (0.277,1.709) (0.275,1.744) (0.273,1.779) (0.272,1.814) (0.271,1.849) (0.270,1.884) (0.268,1.919) (0.267,1.955) (0.266,1.990) (0.266,2.025) (0.265,2.060) (0.264,2.095) (0.263,2.130) (0.263,2.165) (0.262,2.200)};
\addplot[equipot] coordinates {(-0.262,-2.200) (-0.263,-2.165) (-0.263,-2.130) (-0.264,-2.095) (-0.265,-2.060) (-0.266,-2.025) (-0.266,-1.990) (-0.267,-1.955) (-0.268,-1.919) (-0.270,-1.884) (-0.271,-1.849) (-0.272,-1.814) (-0.273,-1.779) (-0.275,-1.744) (-0.277,-1.709) (-0.278,-1.674) (-0.280,-1.639) (-0.283,-1.604) (-0.285,-1.569) (-0.288,-1.534) (-0.291,-1.499) (-0.294,-1.464) (-0.298,-1.429) (-0.302,-1.395) (-0.307,-1.360) (-0.312,-1.325) (-0.318,-1.291) (-0.324,-1.256) (-0.332,-1.222) (-0.340,-1.188) (-0.350,-1.154) (-0.360,-1.121) (-0.373,-1.088) (-0.386,-1.055) (-0.401,-1.024) (-0.417,-0.993) (-0.435,-0.962) (-0.453,-0.933) (-0.473,-0.904) (-0.493,-0.875) (-0.500,-0.842) (-0.500,-0.807) (-0.500,-0.772) (-0.500,-0.737) (-0.500,-0.701) (-0.500,-0.666) (-0.500,-0.631) (-0.500,-0.596) (-0.500,-0.561) (-0.500,-0.526) (-0.500,-0.491) (-0.500,-0.456) (-0.500,-0.421) (-0.500,-0.386) (-0.500,-0.351) (-0.500,-0.316) (-0.500,-0.281) (-0.500,-0.246) (-0.500,-0.210) (-0.500,-0.175) (-0.500,-0.140) (-0.500,-0.105) (-0.500,-0.070) (-0.500,-0.035) (-0.500,-0.000) (-0.500,0.035) (-0.500,0.070) (-0.500,0.105) (-0.500,0.140) (-0.500,0.175) (-0.500,0.210) (-0.500,0.246) (-0.500,0.281) (-0.500,0.316) (-0.500,0.351) (-0.500,0.386) (-0.500,0.421) (-0.500,0.456) (-0.500,0.491) (-0.500,0.526) (-0.500,0.561) (-0.500,0.596) (-0.500,0.631) (-0.500,0.666) (-0.500,0.701) (-0.500,0.737) (-0.500,0.772) (-0.500,0.807) (-0.500,0.842) (-0.493,0.875) (-0.473,0.904) (-0.453,0.933) (-0.435,0.962) (-0.417,0.993) (-0.401,1.024) (-0.386,1.055) (-0.373,1.088) (-0.360,1.121) (-0.350,1.154) (-0.340,1.188) (-0.332,1.222) (-0.324,1.256) (-0.318,1.291) (-0.312,1.325) (-0.307,1.360) (-0.302,1.395) (-0.298,1.429) (-0.294,1.464) (-0.291,1.499) (-0.288,1.534) (-0.285,1.569) (-0.283,1.604) (-0.280,1.639) (-0.278,1.674) (-0.277,1.709) (-0.275,1.744) (-0.273,1.779) (-0.272,1.814) (-0.271,1.849) (-0.270,1.884) (-0.268,1.919) (-0.267,1.955) (-0.266,1.990) (-0.266,2.025) (-0.265,2.060) (-0.264,2.095) (-0.263,2.130) (-0.263,2.165) (-0.262,2.200)};
\addplot[equipot] coordinates {(-0.781,-2.200) (-0.782,-2.165) (-0.784,-2.130) (-0.785,-2.094) (-0.787,-2.059) (-0.788,-2.024) (-0.790,-1.989) (-0.792,-1.954) (-0.794,-1.918) (-0.797,-1.883) (-0.799,-1.848) (-0.801,-1.813) (-0.804,-1.778) (-0.807,-1.743) (-0.810,-1.708) (-0.813,-1.673) (-0.817,-1.638) (-0.820,-1.602) (-0.824,-1.567) (-0.828,-1.532) (-0.833,-1.498) (-0.837,-1.463) (-0.842,-1.428) (-0.848,-1.393) (-0.853,-1.358) (-0.859,-1.323) (-0.866,-1.289) (-0.872,-1.254) (-0.879,-1.220) (-0.887,-1.185) (-0.894,-1.151) (-0.903,-1.117) (-0.911,-1.082) (-0.920,-1.048) (-0.929,-1.014) (-0.938,-0.980) (-0.947,-0.946) (-0.957,-0.912) (-0.967,-0.878) (-0.977,-0.845) (-0.987,-0.811) (-0.997,-0.777) (-1.007,-0.743) (-1.017,-0.710) (-1.027,-0.676) (-1.036,-0.642) (-1.046,-0.608) (-1.055,-0.574) (-1.063,-0.540) (-1.072,-0.505) (-1.080,-0.471) (-1.088,-0.437) (-1.095,-0.402) (-1.102,-0.368) (-1.108,-0.333) (-1.113,-0.298) (-1.119,-0.263) (-1.123,-0.228) (-1.127,-0.193) (-1.130,-0.158) (-1.133,-0.123) (-1.135,-0.088) (-1.136,-0.053) (-1.137,-0.018) (-1.137,0.018) (-1.136,0.053) (-1.135,0.088) (-1.133,0.123) (-1.130,0.158) (-1.127,0.193) (-1.123,0.228) (-1.119,0.263) (-1.113,0.298) (-1.108,0.333) (-1.102,0.368) (-1.095,0.402) (-1.088,0.437) (-1.080,0.471) (-1.072,0.505) (-1.063,0.540) (-1.055,0.574) (-1.046,0.608) (-1.036,0.642) (-1.027,0.676) (-1.017,0.710) (-1.007,0.743) (-0.997,0.777) (-0.987,0.811) (-0.977,0.845) (-0.967,0.878) (-0.957,0.912) (-0.947,0.946) (-0.938,0.980) (-0.929,1.014) (-0.920,1.048) (-0.911,1.082) (-0.903,1.117) (-0.894,1.151) (-0.887,1.185) (-0.879,1.220) (-0.872,1.254) (-0.866,1.289) (-0.859,1.323) (-0.853,1.358) (-0.848,1.393) (-0.842,1.428) (-0.837,1.463) (-0.833,1.498) (-0.828,1.532) (-0.824,1.567) (-0.820,1.602) (-0.817,1.638) (-0.813,1.673) (-0.810,1.708) (-0.807,1.743) (-0.804,1.778) (-0.801,1.813) (-0.799,1.848) (-0.797,1.883) (-0.794,1.918) (-0.792,1.954) (-0.790,1.989) (-0.788,2.024) (-0.787,2.059) (-0.785,2.094) (-0.784,2.130) (-0.782,2.165) (-0.781,2.200)};
\addplot[equipot] coordinates {(-1.289,-2.200) (-1.290,-2.165) (-1.292,-2.130) (-1.293,-2.095) (-1.295,-2.060) (-1.297,-2.025) (-1.298,-1.990) (-1.300,-1.955) (-1.302,-1.920) (-1.304,-1.885) (-1.306,-1.850) (-1.308,-1.815) (-1.311,-1.780) (-1.313,-1.745) (-1.316,-1.710) (-1.318,-1.675) (-1.321,-1.640) (-1.324,-1.605) (-1.326,-1.570) (-1.329,-1.535) (-1.332,-1.500) (-1.336,-1.465) (-1.339,-1.430) (-1.342,-1.395) (-1.346,-1.360) (-1.350,-1.326) (-1.353,-1.291) (-1.357,-1.256) (-1.361,-1.221) (-1.365,-1.186) (-1.370,-1.151) (-1.374,-1.116) (-1.378,-1.082) (-1.383,-1.047) (-1.387,-1.012) (-1.392,-0.977) (-1.396,-0.943) (-1.401,-0.908) (-1.405,-0.873) (-1.410,-0.838) (-1.414,-0.803) (-1.419,-0.769) (-1.423,-0.734) (-1.428,-0.699) (-1.432,-0.664) (-1.436,-0.629) (-1.440,-0.595) (-1.444,-0.560) (-1.448,-0.525) (-1.452,-0.490) (-1.455,-0.455) (-1.459,-0.420) (-1.462,-0.385) (-1.464,-0.350) (-1.467,-0.315) (-1.469,-0.280) (-1.472,-0.245) (-1.473,-0.210) (-1.475,-0.175) (-1.476,-0.140) (-1.477,-0.105) (-1.478,-0.070) (-1.479,-0.035) (-1.479,0.000) (-1.479,0.035) (-1.478,0.070) (-1.477,0.105) (-1.476,0.140) (-1.475,0.175) (-1.473,0.210) (-1.472,0.245) (-1.469,0.280) (-1.467,0.315) (-1.464,0.350) (-1.462,0.385) (-1.459,0.420) (-1.455,0.455) (-1.452,0.490) (-1.448,0.525) (-1.444,0.560) (-1.440,0.595) (-1.436,0.629) (-1.432,0.664) (-1.428,0.699) (-1.423,0.734) (-1.419,0.769) (-1.414,0.803) (-1.410,0.838) (-1.405,0.873) (-1.401,0.908) (-1.396,0.943) (-1.392,0.977) (-1.387,1.012) (-1.383,1.047) (-1.378,1.082) (-1.374,1.116) (-1.370,1.151) (-1.365,1.186) (-1.361,1.221) (-1.357,1.256) (-1.353,1.291) (-1.350,1.326) (-1.346,1.360) (-1.342,1.395) (-1.339,1.430) (-1.336,1.465) (-1.332,1.500) (-1.329,1.535) (-1.326,1.570) (-1.324,1.605) (-1.321,1.640) (-1.318,1.675) (-1.316,1.710) (-1.313,1.745) (-1.311,1.780) (-1.308,1.815) (-1.306,1.850) (-1.304,1.885) (-1.302,1.920) (-1.300,1.955) (-1.298,1.990) (-1.297,2.025) (-1.295,2.060) (-1.293,2.095) (-1.292,2.130) (-1.290,2.165) (-1.289,2.200)};
\addplot[equipot] coordinates {(-1.789,-2.200) (-1.790,-2.165) (-1.792,-2.130) (-1.793,-2.094) (-1.794,-2.059) (-1.795,-2.024) (-1.797,-1.989) (-1.798,-1.953) (-1.799,-1.918) (-1.801,-1.883) (-1.802,-1.848) (-1.804,-1.813) (-1.806,-1.777) (-1.807,-1.742) (-1.809,-1.707) (-1.810,-1.672) (-1.812,-1.637) (-1.814,-1.601) (-1.816,-1.566) (-1.818,-1.531) (-1.820,-1.496) (-1.822,-1.461) (-1.824,-1.425) (-1.826,-1.390) (-1.828,-1.355) (-1.830,-1.320) (-1.832,-1.285) (-1.834,-1.250) (-1.836,-1.214) (-1.838,-1.179) (-1.840,-1.144) (-1.843,-1.109) (-1.845,-1.074) (-1.847,-1.039) (-1.849,-1.003) (-1.851,-0.968) (-1.854,-0.933) (-1.856,-0.898) (-1.858,-0.863) (-1.860,-0.827) (-1.862,-0.792) (-1.864,-0.757) (-1.866,-0.722) (-1.868,-0.687) (-1.870,-0.652) (-1.872,-0.616) (-1.874,-0.581) (-1.876,-0.546) (-1.877,-0.511) (-1.879,-0.476) (-1.881,-0.440) (-1.882,-0.405) (-1.883,-0.370) (-1.884,-0.335) (-1.885,-0.300) (-1.886,-0.264) (-1.887,-0.229) (-1.888,-0.194) (-1.889,-0.159) (-1.889,-0.123) (-1.890,-0.088) (-1.890,-0.053) (-1.890,-0.018) (-1.890,0.018) (-1.890,0.053) (-1.890,0.088) (-1.889,0.123) (-1.889,0.159) (-1.888,0.194) (-1.887,0.229) (-1.886,0.264) (-1.885,0.300) (-1.884,0.335) (-1.883,0.370) (-1.882,0.405) (-1.881,0.440) (-1.879,0.476) (-1.877,0.511) (-1.876,0.546) (-1.874,0.581) (-1.872,0.616) (-1.870,0.652) (-1.868,0.687) (-1.866,0.722) (-1.864,0.757) (-1.862,0.792) (-1.860,0.827) (-1.858,0.863) (-1.856,0.898) (-1.854,0.933) (-1.851,0.968) (-1.849,1.003) (-1.847,1.039) (-1.845,1.074) (-1.843,1.109) (-1.840,1.144) (-1.838,1.179) (-1.836,1.214) (-1.834,1.250) (-1.832,1.285) (-1.830,1.320) (-1.828,1.355) (-1.826,1.390) (-1.824,1.425) (-1.822,1.461) (-1.820,1.496) (-1.818,1.531) (-1.816,1.566) (-1.814,1.601) (-1.812,1.637) (-1.810,1.672) (-1.809,1.707) (-1.807,1.742) (-1.806,1.777) (-1.804,1.813) (-1.802,1.848) (-1.801,1.883) (-1.799,1.918) (-1.798,1.953) (-1.797,1.989) (-1.795,2.024) (-1.794,2.059) (-1.793,2.094) (-1.792,2.130) (-1.790,2.165) (-1.789,2.200)};
\addplot[equipot] coordinates {(-2.286,-2.200) (-2.287,-2.165) (-2.287,-2.130) (-2.288,-2.094) (-2.289,-2.059) (-2.290,-2.024) (-2.291,-1.989) (-2.292,-1.954) (-2.293,-1.918) (-2.294,-1.883) (-2.295,-1.848) (-2.296,-1.813) (-2.297,-1.778) (-2.298,-1.742) (-2.299,-1.707) (-2.300,-1.672) (-2.301,-1.637) (-2.302,-1.602) (-2.303,-1.566) (-2.304,-1.531) (-2.306,-1.496) (-2.307,-1.461) (-2.308,-1.426) (-2.309,-1.390) (-2.310,-1.355) (-2.311,-1.320) (-2.312,-1.285) (-2.314,-1.250) (-2.315,-1.214) (-2.316,-1.179) (-2.317,-1.144) (-2.318,-1.109) (-2.319,-1.074) (-2.321,-1.039) (-2.322,-1.003) (-2.323,-0.968) (-2.324,-0.933) (-2.325,-0.898) (-2.326,-0.863) (-2.327,-0.827) (-2.328,-0.792) (-2.329,-0.757) (-2.330,-0.722) (-2.331,-0.687) (-2.332,-0.651) (-2.333,-0.616) (-2.334,-0.581) (-2.335,-0.546) (-2.335,-0.511) (-2.336,-0.475) (-2.337,-0.440) (-2.338,-0.405) (-2.338,-0.370) (-2.339,-0.335) (-2.339,-0.299) (-2.340,-0.264) (-2.340,-0.229) (-2.340,-0.194) (-2.341,-0.158) (-2.341,-0.123) (-2.341,-0.088) (-2.341,-0.053) (-2.341,-0.018) (-2.341,0.018) (-2.341,0.053) (-2.341,0.088) (-2.341,0.123) (-2.341,0.158) (-2.340,0.194) (-2.340,0.229) (-2.340,0.264) (-2.339,0.299) (-2.339,0.335) (-2.338,0.370) (-2.338,0.405) (-2.337,0.440) (-2.336,0.475) (-2.335,0.511) (-2.335,0.546) (-2.334,0.581) (-2.333,0.616) (-2.332,0.651) (-2.331,0.687) (-2.330,0.722) (-2.329,0.757) (-2.328,0.792) (-2.327,0.827) (-2.326,0.863) (-2.325,0.898) (-2.324,0.933) (-2.323,0.968) (-2.322,1.003) (-2.321,1.039) (-2.319,1.074) (-2.318,1.109) (-2.317,1.144) (-2.316,1.179) (-2.315,1.214) (-2.314,1.250) (-2.312,1.285) (-2.311,1.320) (-2.310,1.355) (-2.309,1.390) (-2.308,1.426) (-2.307,1.461) (-2.306,1.496) (-2.304,1.531) (-2.303,1.566) (-2.302,1.602) (-2.301,1.637) (-2.300,1.672) (-2.299,1.707) (-2.298,1.742) (-2.297,1.778) (-2.296,1.813) (-2.295,1.848) (-2.294,1.883) (-2.293,1.918) (-2.292,1.954) (-2.291,1.989) (-2.290,2.024) (-2.289,2.059) (-2.288,2.094) (-2.287,2.130) (-2.287,2.165) (-2.286,2.200)};
\addplot[equipot] coordinates {(-2.781,-2.200) (-2.782,-2.165) (-2.782,-2.130) (-2.783,-2.094) (-2.784,-2.059) (-2.784,-2.024) (-2.785,-1.989) (-2.785,-1.954) (-2.786,-1.918) (-2.787,-1.883) (-2.787,-1.848) (-2.788,-1.813) (-2.789,-1.778) (-2.789,-1.742) (-2.790,-1.707) (-2.791,-1.672) (-2.791,-1.637) (-2.792,-1.602) (-2.793,-1.566) (-2.793,-1.531) (-2.794,-1.496) (-2.795,-1.461) (-2.795,-1.426) (-2.796,-1.390) (-2.797,-1.355) (-2.797,-1.320) (-2.798,-1.285) (-2.799,-1.250) (-2.799,-1.214) (-2.800,-1.179) (-2.801,-1.144) (-2.801,-1.109) (-2.802,-1.074) (-2.802,-1.038) (-2.803,-1.003) (-2.804,-0.968) (-2.804,-0.933) (-2.805,-0.898) (-2.805,-0.862) (-2.806,-0.827) (-2.807,-0.792) (-2.807,-0.757) (-2.808,-0.722) (-2.808,-0.686) (-2.809,-0.651) (-2.809,-0.616) (-2.809,-0.581) (-2.810,-0.546) (-2.810,-0.510) (-2.811,-0.475) (-2.811,-0.440) (-2.811,-0.405) (-2.812,-0.370) (-2.812,-0.334) (-2.812,-0.299) (-2.812,-0.264) (-2.813,-0.229) (-2.813,-0.194) (-2.813,-0.158) (-2.813,-0.123) (-2.813,-0.088) (-2.813,-0.053) (-2.813,-0.018) (-2.813,0.018) (-2.813,0.053) (-2.813,0.088) (-2.813,0.123) (-2.813,0.158) (-2.813,0.194) (-2.813,0.229) (-2.812,0.264) (-2.812,0.299) (-2.812,0.334) (-2.812,0.370) (-2.811,0.405) (-2.811,0.440) (-2.811,0.475) (-2.810,0.510) (-2.810,0.546) (-2.809,0.581) (-2.809,0.616) (-2.809,0.651) (-2.808,0.686) (-2.808,0.722) (-2.807,0.757) (-2.807,0.792) (-2.806,0.827) (-2.805,0.862) (-2.805,0.898) (-2.804,0.933) (-2.804,0.968) (-2.803,1.003) (-2.802,1.038) (-2.802,1.074) (-2.801,1.109) (-2.801,1.144) (-2.800,1.179) (-2.799,1.214) (-2.799,1.250) (-2.798,1.285) (-2.797,1.320) (-2.797,1.355) (-2.796,1.390) (-2.795,1.426) (-2.795,1.461) (-2.794,1.496) (-2.793,1.531) (-2.793,1.566) (-2.792,1.602) (-2.791,1.637) (-2.791,1.672) (-2.790,1.707) (-2.789,1.742) (-2.789,1.778) (-2.788,1.813) (-2.787,1.848) (-2.787,1.883) (-2.786,1.918) (-2.785,1.954) (-2.785,1.989) (-2.784,2.024) (-2.784,2.059) (-2.783,2.094) (-2.782,2.130) (-2.782,2.165) (-2.781,2.200)};
\addplot[lineaD, flechaen=0.142, flechaen=0.858] coordinates {(-3.000,0.196) (-2.965,0.196) (-2.930,0.196) (-2.895,0.196) (-2.860,0.196) (-2.824,0.196) (-2.789,0.196) (-2.754,0.195) (-2.719,0.195) (-2.684,0.195) (-2.649,0.195) (-2.614,0.195) (-2.579,0.194) (-2.544,0.194) (-2.509,0.194) (-2.473,0.194) (-2.438,0.193) (-2.403,0.193) (-2.368,0.193) (-2.333,0.193) (-2.298,0.192) (-2.263,0.192) (-2.228,0.192) (-2.193,0.191) (-2.158,0.191) (-2.122,0.190) (-2.087,0.190) (-2.052,0.189) (-2.017,0.189) (-1.982,0.188) (-1.947,0.188) (-1.912,0.187) (-1.877,0.187) (-1.842,0.186) (-1.806,0.185) (-1.771,0.184) (-1.736,0.183) (-1.701,0.183) (-1.666,0.182) (-1.631,0.181) (-1.596,0.179) (-1.561,0.178) (-1.526,0.177) (-1.491,0.176) (-1.456,0.174) (-1.421,0.173) (-1.385,0.171) (-1.350,0.169) (-1.315,0.167) (-1.280,0.165) (-1.245,0.163) (-1.210,0.161) (-1.175,0.158) (-1.140,0.155) (-1.105,0.152) (-1.070,0.149) (-1.035,0.146) (-1.000,0.143) (-0.965,0.141) (-0.930,0.141) (-0.895,0.141) (-0.860,0.141) (-0.825,0.141) (-0.790,0.141) (-0.755,0.141) (-0.720,0.141) (-0.685,0.141) (-0.649,0.141) (-0.614,0.141) (-0.579,0.141) (-0.544,0.141) (-0.509,0.141) (-0.474,0.141) (-0.439,0.141) (-0.404,0.141) (-0.369,0.141) (-0.333,0.141) (-0.298,0.141) (-0.263,0.141) (-0.228,0.141) (-0.193,0.141) (-0.158,0.141) (-0.123,0.141) (-0.088,0.141) (-0.053,0.141) (-0.018,0.141) (0.018,0.141) (0.053,0.141) (0.088,0.141) (0.123,0.141) (0.158,0.141) (0.193,0.141) (0.228,0.141) (0.263,0.141) (0.298,0.141) (0.333,0.141) (0.369,0.141) (0.404,0.141) (0.439,0.141) (0.474,0.141) (0.509,0.141) (0.544,0.141) (0.579,0.141) (0.614,0.141) (0.649,0.141) (0.685,0.141) (0.720,0.141) (0.755,0.141) (0.790,0.141) (0.825,0.141) (0.860,0.141) (0.895,0.141) (0.930,0.141) (0.965,0.141) (1.000,0.143) (1.035,0.146) (1.070,0.149) (1.105,0.152) (1.140,0.155) (1.175,0.158) (1.210,0.161) (1.245,0.163) (1.280,0.165) (1.315,0.167) (1.350,0.169) (1.385,0.171) (1.421,0.173) (1.456,0.174) (1.491,0.176) (1.526,0.177) (1.561,0.178) (1.596,0.179) (1.631,0.181) (1.666,0.182) (1.701,0.183) (1.736,0.183) (1.771,0.184) (1.806,0.185) (1.842,0.186) (1.877,0.187) (1.912,0.187) (1.947,0.188) (1.982,0.188) (2.017,0.189) (2.052,0.189) (2.087,0.190) (2.122,0.190) (2.158,0.191) (2.193,0.191) (2.228,0.192) (2.263,0.192) (2.298,0.192) (2.333,0.193) (2.368,0.193) (2.403,0.193) (2.438,0.193) (2.473,0.194) (2.509,0.194) (2.544,0.194) (2.579,0.194) (2.614,0.195) (2.649,0.195) (2.684,0.195) (2.719,0.195) (2.754,0.195) (2.789,0.196) (2.824,0.196) (2.860,0.196) (2.895,0.196) (2.930,0.196) (2.965,0.196) (3.000,0.196)};
\addplot[lineaD, flechaen=0.141, flechaen=0.859] coordinates {(-3.000,0.393) (-2.965,0.393) (-2.930,0.392) (-2.895,0.392) (-2.859,0.392) (-2.824,0.392) (-2.789,0.391) (-2.754,0.391) (-2.719,0.391) (-2.684,0.390) (-2.648,0.390) (-2.613,0.390) (-2.578,0.389) (-2.543,0.389) (-2.508,0.388) (-2.473,0.388) (-2.438,0.387) (-2.402,0.387) (-2.367,0.386) (-2.332,0.386) (-2.297,0.385) (-2.262,0.384) (-2.227,0.384) (-2.191,0.383) (-2.156,0.382) (-2.121,0.381) (-2.086,0.381) (-2.051,0.380) (-2.016,0.379) (-1.981,0.378) (-1.945,0.377) (-1.910,0.375) (-1.875,0.374) (-1.840,0.373) (-1.805,0.371) (-1.770,0.370) (-1.735,0.368) (-1.700,0.367) (-1.664,0.365) (-1.629,0.363) (-1.594,0.361) (-1.559,0.359) (-1.524,0.356) (-1.489,0.354) (-1.454,0.351) (-1.419,0.348) (-1.384,0.345) (-1.349,0.342) (-1.314,0.338) (-1.279,0.334) (-1.244,0.330) (-1.209,0.326) (-1.174,0.321) (-1.140,0.316) (-1.105,0.310) (-1.070,0.304) (-1.036,0.298) (-1.001,0.292) (-0.967,0.284) (-0.932,0.283) (-0.896,0.283) (-0.861,0.283) (-0.826,0.283) (-0.791,0.283) (-0.756,0.283) (-0.721,0.283) (-0.686,0.283) (-0.650,0.283) (-0.615,0.283) (-0.580,0.283) (-0.545,0.283) (-0.510,0.283) (-0.475,0.283) (-0.439,0.283) (-0.404,0.283) (-0.369,0.283) (-0.334,0.283) (-0.299,0.283) (-0.264,0.283) (-0.229,0.283) (-0.193,0.283) (-0.158,0.283) (-0.123,0.283) (-0.088,0.283) (-0.053,0.283) (-0.018,0.283) (0.018,0.283) (0.053,0.283) (0.088,0.283) (0.123,0.283) (0.158,0.283) (0.193,0.283) (0.229,0.283) (0.264,0.283) (0.299,0.283) (0.334,0.283) (0.369,0.283) (0.404,0.283) (0.439,0.283) (0.475,0.283) (0.510,0.283) (0.545,0.283) (0.580,0.283) (0.615,0.283) (0.650,0.283) (0.686,0.283) (0.721,0.283) (0.756,0.283) (0.791,0.283) (0.826,0.283) (0.861,0.283) (0.896,0.283) (0.932,0.283) (0.967,0.284) (1.001,0.292) (1.036,0.298) (1.070,0.304) (1.105,0.310) (1.140,0.316) (1.174,0.321) (1.209,0.326) (1.244,0.330) (1.279,0.334) (1.314,0.338) (1.349,0.342) (1.384,0.345) (1.419,0.348) (1.454,0.351) (1.489,0.354) (1.524,0.356) (1.559,0.359) (1.594,0.361) (1.629,0.363) (1.664,0.365) (1.700,0.367) (1.735,0.368) (1.770,0.370) (1.805,0.371) (1.840,0.373) (1.875,0.374) (1.910,0.375) (1.945,0.377) (1.981,0.378) (2.016,0.379) (2.051,0.380) (2.086,0.381) (2.121,0.381) (2.156,0.382) (2.191,0.383) (2.227,0.384) (2.262,0.384) (2.297,0.385) (2.332,0.386) (2.367,0.386) (2.402,0.387) (2.438,0.387) (2.473,0.388) (2.508,0.388) (2.543,0.389) (2.578,0.389) (2.613,0.390) (2.648,0.390) (2.684,0.390) (2.719,0.391) (2.754,0.391) (2.789,0.391) (2.824,0.392) (2.859,0.392) (2.895,0.392) (2.930,0.392) (2.965,0.393) (3.000,0.393)};
\addplot[lineaD, flechaen=0.141, flechaen=0.859] coordinates {(-3.000,0.590) (-2.965,0.589) (-2.930,0.589) (-2.895,0.589) (-2.860,0.588) (-2.825,0.588) (-2.790,0.587) (-2.755,0.587) (-2.720,0.587) (-2.685,0.586) (-2.650,0.586) (-2.615,0.585) (-2.580,0.584) (-2.545,0.584) (-2.510,0.583) (-2.474,0.583) (-2.439,0.582) (-2.404,0.581) (-2.369,0.580) (-2.334,0.580) (-2.299,0.579) (-2.264,0.578) (-2.229,0.577) (-2.194,0.576) (-2.159,0.575) (-2.124,0.574) (-2.089,0.572) (-2.054,0.571) (-2.019,0.570) (-1.984,0.569) (-1.949,0.567) (-1.914,0.565) (-1.879,0.564) (-1.844,0.562) (-1.809,0.560) (-1.774,0.558) (-1.739,0.556) (-1.704,0.554) (-1.669,0.551) (-1.634,0.549) (-1.599,0.546) (-1.565,0.543) (-1.530,0.540) (-1.495,0.537) (-1.460,0.533) (-1.425,0.529) (-1.390,0.525) (-1.356,0.521) (-1.321,0.516) (-1.286,0.511) (-1.251,0.506) (-1.217,0.500) (-1.182,0.494) (-1.148,0.487) (-1.114,0.480) (-1.080,0.473) (-1.045,0.465) (-1.011,0.456) (-0.978,0.447) (-0.944,0.437) (-0.911,0.426) (-0.876,0.424) (-0.841,0.424) (-0.806,0.424) (-0.771,0.424) (-0.736,0.424) (-0.701,0.424) (-0.666,0.424) (-0.631,0.424) (-0.596,0.424) (-0.561,0.424) (-0.526,0.424) (-0.491,0.424) (-0.455,0.424) (-0.420,0.424) (-0.385,0.424) (-0.350,0.424) (-0.315,0.424) (-0.280,0.424) (-0.245,0.424) (-0.210,0.424) (-0.175,0.424) (-0.140,0.424) (-0.105,0.424) (-0.070,0.424) (-0.035,0.424) (0.000,0.424) (0.035,0.424) (0.070,0.424) (0.105,0.424) (0.140,0.424) (0.175,0.424) (0.210,0.424) (0.245,0.424) (0.280,0.424) (0.315,0.424) (0.350,0.424) (0.385,0.424) (0.420,0.424) (0.455,0.424) (0.491,0.424) (0.526,0.424) (0.561,0.424) (0.596,0.424) (0.631,0.424) (0.666,0.424) (0.701,0.424) (0.736,0.424) (0.771,0.424) (0.806,0.424) (0.841,0.424) (0.876,0.424) (0.911,0.426) (0.944,0.437) (0.978,0.447) (1.011,0.456) (1.045,0.465) (1.080,0.473) (1.114,0.480) (1.148,0.487) (1.182,0.494) (1.217,0.500) (1.251,0.506) (1.286,0.511) (1.321,0.516) (1.356,0.521) (1.390,0.525) (1.425,0.529) (1.460,0.533) (1.495,0.537) (1.530,0.540) (1.565,0.543) (1.599,0.546) (1.634,0.549) (1.669,0.551) (1.704,0.554) (1.739,0.556) (1.774,0.558) (1.809,0.560) (1.844,0.562) (1.879,0.564) (1.914,0.565) (1.949,0.567) (1.984,0.569) (2.019,0.570) (2.054,0.571) (2.089,0.572) (2.124,0.574) (2.159,0.575) (2.194,0.576) (2.229,0.577) (2.264,0.578) (2.299,0.579) (2.334,0.580) (2.369,0.580) (2.404,0.581) (2.439,0.582) (2.474,0.583) (2.510,0.583) (2.545,0.584) (2.580,0.584) (2.615,0.585) (2.650,0.586) (2.685,0.586) (2.720,0.587) (2.755,0.587) (2.790,0.587) (2.825,0.588) (2.860,0.588) (2.895,0.589) (2.930,0.589) (2.965,0.589) (3.000,0.590)};
\addplot[lineaD, flechaen=0.141, flechaen=0.859] coordinates {(-3.000,0.787) (-2.965,0.786) (-2.930,0.786) (-2.895,0.786) (-2.859,0.785) (-2.824,0.785) (-2.789,0.784) (-2.754,0.783) (-2.719,0.783) (-2.684,0.782) (-2.649,0.782) (-2.613,0.781) (-2.578,0.780) (-2.543,0.780) (-2.508,0.779) (-2.473,0.778) (-2.438,0.777) (-2.402,0.776) (-2.367,0.775) (-2.332,0.774) (-2.297,0.773) (-2.262,0.772) (-2.227,0.771) (-2.192,0.770) (-2.157,0.769) (-2.121,0.767) (-2.086,0.766) (-2.051,0.764) (-2.016,0.763) (-1.981,0.761) (-1.946,0.759) (-1.911,0.758) (-1.876,0.756) (-1.841,0.754) (-1.805,0.751) (-1.770,0.749) (-1.735,0.747) (-1.700,0.744) (-1.665,0.741) (-1.630,0.738) (-1.595,0.735) (-1.560,0.732) (-1.525,0.728) (-1.490,0.724) (-1.455,0.720) (-1.420,0.716) (-1.386,0.711) (-1.351,0.706) (-1.316,0.701) (-1.281,0.696) (-1.247,0.690) (-1.212,0.683) (-1.178,0.676) (-1.143,0.669) (-1.109,0.661) (-1.075,0.653) (-1.041,0.643) (-1.007,0.634) (-0.974,0.623) (-0.940,0.612) (-0.907,0.600) (-0.874,0.587) (-0.842,0.574) (-0.809,0.566) (-0.773,0.566) (-0.738,0.566) (-0.703,0.566) (-0.668,0.566) (-0.633,0.566) (-0.598,0.566) (-0.562,0.566) (-0.527,0.566) (-0.492,0.566) (-0.457,0.566) (-0.422,0.566) (-0.387,0.566) (-0.352,0.566) (-0.316,0.566) (-0.281,0.566) (-0.246,0.566) (-0.211,0.566) (-0.176,0.566) (-0.141,0.566) (-0.105,0.566) (-0.070,0.566) (-0.035,0.566) (0.000,0.566) (0.035,0.566) (0.070,0.566) (0.105,0.566) (0.141,0.566) (0.176,0.566) (0.211,0.566) (0.246,0.566) (0.281,0.566) (0.316,0.566) (0.352,0.566) (0.387,0.566) (0.422,0.566) (0.457,0.566) (0.492,0.566) (0.527,0.566) (0.562,0.566) (0.598,0.566) (0.633,0.566) (0.668,0.566) (0.703,0.566) (0.738,0.566) (0.773,0.566) (0.809,0.566) (0.842,0.574) (0.874,0.587) (0.907,0.600) (0.940,0.612) (0.974,0.623) (1.007,0.634) (1.041,0.643) (1.075,0.653) (1.109,0.661) (1.143,0.669) (1.178,0.676) (1.212,0.683) (1.247,0.690) (1.281,0.696) (1.316,0.701) (1.351,0.706) (1.386,0.711) (1.420,0.716) (1.455,0.720) (1.490,0.724) (1.525,0.728) (1.560,0.732) (1.595,0.735) (1.630,0.738) (1.665,0.741) (1.700,0.744) (1.735,0.747) (1.770,0.749) (1.805,0.751) (1.841,0.754) (1.876,0.756) (1.911,0.758) (1.946,0.759) (1.981,0.761) (2.016,0.763) (2.051,0.764) (2.086,0.766) (2.121,0.767) (2.157,0.769) (2.192,0.770) (2.227,0.771) (2.262,0.772) (2.297,0.773) (2.332,0.774) (2.367,0.775) (2.402,0.776) (2.438,0.777) (2.473,0.778) (2.508,0.779) (2.543,0.780) (2.578,0.780) (2.613,0.781) (2.649,0.782) (2.684,0.782) (2.719,0.783) (2.754,0.783) (2.789,0.784) (2.824,0.785) (2.859,0.785) (2.895,0.786) (2.930,0.786) (2.965,0.786) (3.000,0.787)};
\addplot[lineaD, flechaen=0.140, flechaen=0.860] coordinates {(-3.000,0.984) (-2.965,0.984) (-2.930,0.983) (-2.895,0.983) (-2.860,0.982) (-2.825,0.982) (-2.789,0.981) (-2.754,0.981) (-2.719,0.980) (-2.684,0.979) (-2.649,0.979) (-2.614,0.978) (-2.579,0.977) (-2.544,0.976) (-2.509,0.975) (-2.474,0.974) (-2.439,0.973) (-2.404,0.972) (-2.368,0.971) (-2.333,0.970) (-2.298,0.969) (-2.263,0.968) (-2.228,0.967) (-2.193,0.966) (-2.158,0.964) (-2.123,0.963) (-2.088,0.961) (-2.053,0.960) (-2.018,0.958) (-1.983,0.956) (-1.948,0.954) (-1.913,0.952) (-1.878,0.950) (-1.843,0.948) (-1.808,0.946) (-1.773,0.943) (-1.738,0.941) (-1.703,0.938) (-1.668,0.935) (-1.633,0.932) (-1.598,0.929) (-1.563,0.926) (-1.528,0.922) (-1.493,0.918) (-1.458,0.914) (-1.423,0.910) (-1.389,0.906) (-1.354,0.901) (-1.319,0.896) (-1.284,0.890) (-1.250,0.884) (-1.215,0.878) (-1.181,0.872) (-1.146,0.865) (-1.112,0.857) (-1.078,0.849) (-1.044,0.841) (-1.010,0.831) (-0.976,0.822) (-0.943,0.811) (-0.910,0.800) (-0.877,0.788) (-0.844,0.775) (-0.812,0.761) (-0.780,0.746) (-0.749,0.731) (-0.718,0.713) (-0.684,0.707) (-0.649,0.707) (-0.614,0.707) (-0.579,0.707) (-0.544,0.707) (-0.509,0.707) (-0.474,0.707) (-0.439,0.707) (-0.404,0.707) (-0.368,0.707) (-0.333,0.707) (-0.298,0.707) (-0.263,0.707) (-0.228,0.707) (-0.193,0.707) (-0.158,0.707) (-0.123,0.707) (-0.088,0.707) (-0.053,0.707) (-0.018,0.707) (0.018,0.707) (0.053,0.707) (0.088,0.707) (0.123,0.707) (0.158,0.707) (0.193,0.707) (0.228,0.707) (0.263,0.707) (0.298,0.707) (0.333,0.707) (0.368,0.707) (0.404,0.707) (0.439,0.707) (0.474,0.707) (0.509,0.707) (0.544,0.707) (0.579,0.707) (0.614,0.707) (0.649,0.707) (0.684,0.707) (0.718,0.713) (0.749,0.731) (0.780,0.746) (0.812,0.761) (0.844,0.775) (0.877,0.788) (0.910,0.800) (0.943,0.811) (0.976,0.822) (1.010,0.831) (1.044,0.841) (1.078,0.849) (1.112,0.857) (1.146,0.865) (1.181,0.872) (1.215,0.878) (1.250,0.884) (1.284,0.890) (1.319,0.896) (1.354,0.901) (1.389,0.906) (1.423,0.910) (1.458,0.914) (1.493,0.918) (1.528,0.922) (1.563,0.926) (1.598,0.929) (1.633,0.932) (1.668,0.935) (1.703,0.938) (1.738,0.941) (1.773,0.943) (1.808,0.946) (1.843,0.948) (1.878,0.950) (1.913,0.952) (1.948,0.954) (1.983,0.956) (2.018,0.958) (2.053,0.960) (2.088,0.961) (2.123,0.963) (2.158,0.964) (2.193,0.966) (2.228,0.967) (2.263,0.968) (2.298,0.969) (2.333,0.970) (2.368,0.971) (2.404,0.972) (2.439,0.973) (2.474,0.974) (2.509,0.975) (2.544,0.976) (2.579,0.977) (2.614,0.978) (2.649,0.979) (2.684,0.979) (2.719,0.980) (2.754,0.981) (2.789,0.981) (2.825,0.982) (2.860,0.982) (2.895,0.983) (2.930,0.983) (2.965,0.984) (3.000,0.984)};
\addplot[lineaD, flechaen=0.139, flechaen=0.861] coordinates {(-3.000,1.182) (-2.965,1.182) (-2.930,1.181) (-2.895,1.181) (-2.860,1.180) (-2.825,1.180) (-2.790,1.179) (-2.755,1.178) (-2.720,1.178) (-2.685,1.177) (-2.650,1.176) (-2.615,1.175) (-2.580,1.174) (-2.545,1.174) (-2.510,1.173) (-2.475,1.172) (-2.440,1.171) (-2.404,1.170) (-2.369,1.169) (-2.334,1.168) (-2.299,1.166) (-2.264,1.165) (-2.229,1.164) (-2.194,1.163) (-2.159,1.161) (-2.124,1.160) (-2.089,1.158) (-2.054,1.157) (-2.019,1.155) (-1.984,1.153) (-1.949,1.151) (-1.914,1.149) (-1.879,1.147) (-1.844,1.145) (-1.809,1.143) (-1.774,1.141) (-1.740,1.138) (-1.705,1.136) (-1.670,1.133) (-1.635,1.130) (-1.600,1.127) (-1.565,1.124) (-1.530,1.121) (-1.495,1.117) (-1.460,1.113) (-1.426,1.110) (-1.391,1.105) (-1.356,1.101) (-1.321,1.097) (-1.287,1.092) (-1.252,1.087) (-1.217,1.081) (-1.183,1.076) (-1.148,1.070) (-1.114,1.063) (-1.079,1.056) (-1.045,1.049) (-1.011,1.042) (-0.977,1.034) (-0.943,1.025) (-0.909,1.016) (-0.875,1.006) (-0.842,0.996) (-0.808,0.985) (-0.775,0.974) (-0.742,0.961) (-0.710,0.948) (-0.678,0.934) (-0.646,0.919) (-0.615,0.903) (-0.585,0.885) (-0.555,0.867) (-0.526,0.849) (-0.491,0.849) (-0.456,0.849) (-0.420,0.849) (-0.385,0.849) (-0.350,0.849) (-0.315,0.849) (-0.280,0.849) (-0.245,0.849) (-0.210,0.849) (-0.175,0.849) (-0.140,0.849) (-0.105,0.849) (-0.070,0.849) (-0.035,0.849) (0.000,0.849) (0.035,0.849) (0.070,0.849) (0.105,0.849) (0.140,0.849) (0.175,0.849) (0.210,0.849) (0.245,0.849) (0.280,0.849) (0.315,0.849) (0.350,0.849) (0.385,0.849) (0.420,0.849) (0.456,0.849) (0.491,0.849) (0.526,0.849) (0.555,0.867) (0.585,0.885) (0.615,0.903) (0.646,0.919) (0.678,0.934) (0.710,0.948) (0.742,0.961) (0.775,0.974) (0.808,0.985) (0.842,0.996) (0.875,1.006) (0.909,1.016) (0.943,1.025) (0.977,1.034) (1.011,1.042) (1.045,1.049) (1.079,1.056) (1.114,1.063) (1.148,1.070) (1.183,1.076) (1.217,1.081) (1.252,1.087) (1.287,1.092) (1.321,1.097) (1.356,1.101) (1.391,1.105) (1.426,1.110) (1.460,1.113) (1.495,1.117) (1.530,1.121) (1.565,1.124) (1.600,1.127) (1.635,1.130) (1.670,1.133) (1.705,1.136) (1.740,1.138) (1.774,1.141) (1.809,1.143) (1.844,1.145) (1.879,1.147) (1.914,1.149) (1.949,1.151) (1.984,1.153) (2.019,1.155) (2.054,1.157) (2.089,1.158) (2.124,1.160) (2.159,1.161) (2.194,1.163) (2.229,1.164) (2.264,1.165) (2.299,1.166) (2.334,1.168) (2.369,1.169) (2.404,1.170) (2.440,1.171) (2.475,1.172) (2.510,1.173) (2.545,1.174) (2.580,1.174) (2.615,1.175) (2.650,1.176) (2.685,1.177) (2.720,1.178) (2.755,1.178) (2.790,1.179) (2.825,1.180) (2.860,1.180) (2.895,1.181) (2.930,1.181) (2.965,1.182) (3.000,1.182)};
\addplot[lineaD, flechaen=0.139, flechaen=0.861] coordinates {(-3.000,1.381) (-2.965,1.380) (-2.930,1.380) (-2.895,1.379) (-2.860,1.379) (-2.825,1.378) (-2.790,1.377) (-2.755,1.377) (-2.720,1.376) (-2.685,1.375) (-2.650,1.374) (-2.615,1.374) (-2.580,1.373) (-2.545,1.372) (-2.510,1.371) (-2.475,1.370) (-2.439,1.369) (-2.404,1.368) (-2.369,1.367) (-2.334,1.366) (-2.299,1.365) (-2.264,1.364) (-2.229,1.362) (-2.194,1.361) (-2.159,1.360) (-2.124,1.358) (-2.089,1.357) (-2.054,1.355) (-2.019,1.354) (-1.984,1.352) (-1.949,1.350) (-1.914,1.348) (-1.879,1.347) (-1.844,1.345) (-1.809,1.343) (-1.774,1.340) (-1.739,1.338) (-1.704,1.336) (-1.669,1.333) (-1.635,1.331) (-1.600,1.328) (-1.565,1.325) (-1.530,1.322) (-1.495,1.319) (-1.460,1.316) (-1.425,1.313) (-1.390,1.309) (-1.355,1.305) (-1.321,1.302) (-1.286,1.298) (-1.251,1.293) (-1.216,1.289) (-1.181,1.284) (-1.147,1.280) (-1.112,1.274) (-1.077,1.269) (-1.043,1.264) (-1.008,1.258) (-0.974,1.252) (-0.939,1.246) (-0.905,1.239) (-0.871,1.232) (-0.836,1.225) (-0.802,1.217) (-0.768,1.209) (-0.734,1.201) (-0.700,1.193) (-0.666,1.184) (-0.632,1.174) (-0.599,1.164) (-0.565,1.154) (-0.532,1.144) (-0.498,1.133) (-0.465,1.121) (-0.432,1.109) (-0.399,1.097) (-0.367,1.085) (-0.334,1.071) (-0.302,1.058) (-0.270,1.044) (-0.237,1.031) (-0.205,1.017) (-0.173,1.003) (-0.140,0.991) (-0.105,0.990) (-0.070,0.990) (-0.035,0.990) (-0.000,0.990) (0.035,0.990) (0.070,0.990) (0.105,0.990) (0.140,0.991) (0.173,1.003) (0.205,1.017) (0.237,1.031) (0.270,1.044) (0.302,1.058) (0.334,1.071) (0.367,1.085) (0.399,1.097) (0.432,1.109) (0.465,1.121) (0.498,1.133) (0.532,1.144) (0.565,1.154) (0.599,1.164) (0.632,1.174) (0.666,1.184) (0.700,1.193) (0.734,1.201) (0.768,1.209) (0.802,1.217) (0.836,1.225) (0.871,1.232) (0.905,1.239) (0.939,1.246) (0.974,1.252) (1.008,1.258) (1.043,1.264) (1.077,1.269) (1.112,1.274) (1.147,1.280) (1.181,1.284) (1.216,1.289) (1.251,1.293) (1.286,1.298) (1.321,1.302) (1.355,1.305) (1.390,1.309) (1.425,1.313) (1.460,1.316) (1.495,1.319) (1.530,1.322) (1.565,1.325) (1.600,1.328) (1.635,1.331) (1.669,1.333) (1.704,1.336) (1.739,1.338) (1.774,1.340) (1.809,1.343) (1.844,1.345) (1.879,1.347) (1.914,1.348) (1.949,1.350) (1.984,1.352) (2.019,1.354) (2.054,1.355) (2.089,1.357) (2.124,1.358) (2.159,1.360) (2.194,1.361) (2.229,1.362) (2.264,1.364) (2.299,1.365) (2.334,1.366) (2.369,1.367) (2.404,1.368) (2.439,1.369) (2.475,1.370) (2.510,1.371) (2.545,1.372) (2.580,1.373) (2.615,1.374) (2.650,1.374) (2.685,1.375) (2.720,1.376) (2.755,1.377) (2.790,1.377) (2.825,1.378) (2.860,1.379) (2.895,1.379) (2.930,1.380) (2.965,1.380) (3.000,1.381)};
\addplot[lineaD, flechaen=0.141, flechaen=0.859] coordinates {(-3.000,1.580) (-2.965,1.579) (-2.930,1.579) (-2.895,1.578) (-2.860,1.578) (-2.825,1.577) (-2.790,1.576) (-2.755,1.576) (-2.720,1.575) (-2.685,1.574) (-2.650,1.573) (-2.615,1.573) (-2.580,1.572) (-2.545,1.571) (-2.510,1.570) (-2.475,1.569) (-2.440,1.568) (-2.405,1.567) (-2.370,1.566) (-2.335,1.565) (-2.300,1.564) (-2.265,1.563) (-2.230,1.562) (-2.195,1.561) (-2.160,1.559) (-2.125,1.558) (-2.090,1.557) (-2.055,1.555) (-2.020,1.554) (-1.985,1.552) (-1.950,1.551) (-1.915,1.549) (-1.880,1.547) (-1.845,1.546) (-1.810,1.544) (-1.775,1.542) (-1.740,1.540) (-1.705,1.538) (-1.670,1.536) (-1.635,1.533) (-1.600,1.531) (-1.565,1.529) (-1.530,1.526) (-1.495,1.524) (-1.460,1.521) (-1.425,1.518) (-1.390,1.515) (-1.355,1.512) (-1.320,1.509) (-1.286,1.506) (-1.251,1.503) (-1.216,1.499) (-1.181,1.496) (-1.146,1.492) (-1.111,1.488) (-1.077,1.484) (-1.042,1.480) (-1.007,1.476) (-0.972,1.471) (-0.937,1.467) (-0.903,1.462) (-0.868,1.458) (-0.833,1.453) (-0.799,1.448) (-0.764,1.443) (-0.729,1.438) (-0.695,1.432) (-0.660,1.427) (-0.626,1.421) (-0.591,1.416) (-0.556,1.410) (-0.522,1.405) (-0.487,1.399) (-0.453,1.394) (-0.418,1.388) (-0.383,1.383) (-0.349,1.378) (-0.314,1.373) (-0.279,1.368) (-0.245,1.364) (-0.210,1.360) (-0.175,1.357) (-0.140,1.354) (-0.105,1.352) (-0.070,1.350) (-0.035,1.349) (0.000,1.348) (0.035,1.349) (0.070,1.350) (0.105,1.352) (0.140,1.354) (0.175,1.357) (0.210,1.360) (0.245,1.364) (0.279,1.368) (0.314,1.373) (0.349,1.378) (0.383,1.383) (0.418,1.388) (0.453,1.394) (0.487,1.399) (0.522,1.405) (0.556,1.410) (0.591,1.416) (0.626,1.421) (0.660,1.427) (0.695,1.432) (0.729,1.438) (0.764,1.443) (0.799,1.448) (0.833,1.453) (0.868,1.458) (0.903,1.462) (0.937,1.467) (0.972,1.471) (1.007,1.476) (1.042,1.480) (1.077,1.484) (1.111,1.488) (1.146,1.492) (1.181,1.496) (1.216,1.499) (1.251,1.503) (1.286,1.506) (1.320,1.509) (1.355,1.512) (1.390,1.515) (1.425,1.518) (1.460,1.521) (1.495,1.524) (1.530,1.526) (1.565,1.529) (1.600,1.531) (1.635,1.533) (1.670,1.536) (1.705,1.538) (1.740,1.540) (1.775,1.542) (1.810,1.544) (1.845,1.546) (1.880,1.547) (1.915,1.549) (1.950,1.551) (1.985,1.552) (2.020,1.554) (2.055,1.555) (2.090,1.557) (2.125,1.558) (2.160,1.559) (2.195,1.561) (2.230,1.562) (2.265,1.563) (2.300,1.564) (2.335,1.565) (2.370,1.566) (2.405,1.567) (2.440,1.568) (2.475,1.569) (2.510,1.570) (2.545,1.571) (2.580,1.572) (2.615,1.573) (2.650,1.573) (2.685,1.574) (2.720,1.575) (2.755,1.576) (2.790,1.576) (2.825,1.577) (2.860,1.578) (2.895,1.578) (2.930,1.579) (2.965,1.579) (3.000,1.580)};
\addplot[lineaD, flechaen=0.141, flechaen=0.859] coordinates {(-3.000,1.779) (-2.965,1.779) (-2.930,1.778) (-2.895,1.777) (-2.859,1.777) (-2.824,1.776) (-2.789,1.776) (-2.754,1.775) (-2.719,1.774) (-2.684,1.774) (-2.649,1.773) (-2.613,1.772) (-2.578,1.771) (-2.543,1.770) (-2.508,1.770) (-2.473,1.769) (-2.438,1.768) (-2.403,1.767) (-2.367,1.766) (-2.332,1.765) (-2.297,1.764) (-2.262,1.763) (-2.227,1.762) (-2.192,1.761) (-2.157,1.760) (-2.121,1.758) (-2.086,1.757) (-2.051,1.756) (-2.016,1.755) (-1.981,1.753) (-1.946,1.752) (-1.911,1.750) (-1.876,1.749) (-1.841,1.747) (-1.805,1.746) (-1.770,1.744) (-1.735,1.742) (-1.700,1.741) (-1.665,1.739) (-1.630,1.737) (-1.595,1.735) (-1.560,1.733) (-1.525,1.731) (-1.489,1.729) (-1.454,1.727) (-1.419,1.724) (-1.384,1.722) (-1.349,1.720) (-1.314,1.717) (-1.279,1.715) (-1.244,1.712) (-1.209,1.709) (-1.174,1.707) (-1.139,1.704) (-1.104,1.701) (-1.069,1.698) (-1.034,1.695) (-0.999,1.692) (-0.964,1.689) (-0.929,1.686) (-0.894,1.682) (-0.859,1.679) (-0.824,1.676) (-0.789,1.673) (-0.754,1.669) (-0.719,1.666) (-0.684,1.663) (-0.649,1.659) (-0.614,1.656) (-0.579,1.653) (-0.544,1.649) (-0.509,1.646) (-0.474,1.643) (-0.439,1.640) (-0.404,1.637) (-0.369,1.635) (-0.334,1.632) (-0.299,1.630) (-0.263,1.628) (-0.228,1.626) (-0.193,1.624) (-0.158,1.622) (-0.123,1.621) (-0.088,1.620) (-0.053,1.620) (-0.018,1.619) (0.018,1.619) (0.053,1.620) (0.088,1.620) (0.123,1.621) (0.158,1.622) (0.193,1.624) (0.228,1.626) (0.263,1.628) (0.299,1.630) (0.334,1.632) (0.369,1.635) (0.404,1.637) (0.439,1.640) (0.474,1.643) (0.509,1.646) (0.544,1.649) (0.579,1.653) (0.614,1.656) (0.649,1.659) (0.684,1.663) (0.719,1.666) (0.754,1.669) (0.789,1.673) (0.824,1.676) (0.859,1.679) (0.894,1.682) (0.929,1.686) (0.964,1.689) (0.999,1.692) (1.034,1.695) (1.069,1.698) (1.104,1.701) (1.139,1.704) (1.174,1.707) (1.209,1.709) (1.244,1.712) (1.279,1.715) (1.314,1.717) (1.349,1.720) (1.384,1.722) (1.419,1.724) (1.454,1.727) (1.489,1.729) (1.525,1.731) (1.560,1.733) (1.595,1.735) (1.630,1.737) (1.665,1.739) (1.700,1.741) (1.735,1.742) (1.770,1.744) (1.805,1.746) (1.841,1.747) (1.876,1.749) (1.911,1.750) (1.946,1.752) (1.981,1.753) (2.016,1.755) (2.051,1.756) (2.086,1.757) (2.121,1.758) (2.157,1.760) (2.192,1.761) (2.227,1.762) (2.262,1.763) (2.297,1.764) (2.332,1.765) (2.367,1.766) (2.403,1.767) (2.438,1.768) (2.473,1.769) (2.508,1.770) (2.543,1.770) (2.578,1.771) (2.613,1.772) (2.649,1.773) (2.684,1.774) (2.719,1.774) (2.754,1.775) (2.789,1.776) (2.824,1.776) (2.859,1.777) (2.895,1.777) (2.930,1.778) (2.965,1.779) (3.000,1.779)};
\addplot[lineaD, flechaen=0.142, flechaen=0.858] coordinates {(-3.000,1.979) (-2.965,1.978) (-2.930,1.978) (-2.895,1.977) (-2.860,1.977) (-2.824,1.976) (-2.789,1.975) (-2.754,1.975) (-2.719,1.974) (-2.684,1.973) (-2.649,1.973) (-2.614,1.972) (-2.579,1.971) (-2.544,1.971) (-2.508,1.970) (-2.473,1.969) (-2.438,1.968) (-2.403,1.967) (-2.368,1.967) (-2.333,1.966) (-2.298,1.965) (-2.263,1.964) (-2.228,1.963) (-2.192,1.962) (-2.157,1.961) (-2.122,1.960) (-2.087,1.959) (-2.052,1.958) (-2.017,1.956) (-1.982,1.955) (-1.947,1.954) (-1.912,1.953) (-1.877,1.951) (-1.841,1.950) (-1.806,1.949) (-1.771,1.947) (-1.736,1.946) (-1.701,1.944) (-1.666,1.943) (-1.631,1.941) (-1.596,1.940) (-1.561,1.938) (-1.526,1.936) (-1.491,1.935) (-1.455,1.933) (-1.420,1.931) (-1.385,1.929) (-1.350,1.927) (-1.315,1.926) (-1.280,1.924) (-1.245,1.922) (-1.210,1.920) (-1.175,1.917) (-1.140,1.915) (-1.105,1.913) (-1.070,1.911) (-1.035,1.909) (-1.000,1.907) (-0.965,1.905) (-0.930,1.902) (-0.895,1.900) (-0.859,1.898) (-0.824,1.896) (-0.789,1.893) (-0.754,1.891) (-0.719,1.889) (-0.684,1.887) (-0.649,1.885) (-0.614,1.882) (-0.579,1.880) (-0.544,1.878) (-0.509,1.876) (-0.474,1.875) (-0.439,1.873) (-0.404,1.871) (-0.369,1.869) (-0.334,1.868) (-0.298,1.867) (-0.263,1.865) (-0.228,1.864) (-0.193,1.863) (-0.158,1.862) (-0.123,1.862) (-0.088,1.861) (-0.053,1.861) (-0.018,1.861) (0.018,1.861) (0.053,1.861) (0.088,1.861) (0.123,1.862) (0.158,1.862) (0.193,1.863) (0.228,1.864) (0.263,1.865) (0.298,1.867) (0.334,1.868) (0.369,1.869) (0.404,1.871) (0.439,1.873) (0.474,1.875) (0.509,1.876) (0.544,1.878) (0.579,1.880) (0.614,1.882) (0.649,1.885) (0.684,1.887) (0.719,1.889) (0.754,1.891) (0.789,1.893) (0.824,1.896) (0.859,1.898) (0.895,1.900) (0.930,1.902) (0.965,1.905) (1.000,1.907) (1.035,1.909) (1.070,1.911) (1.105,1.913) (1.140,1.915) (1.175,1.917) (1.210,1.920) (1.245,1.922) (1.280,1.924) (1.315,1.926) (1.350,1.927) (1.385,1.929) (1.420,1.931) (1.455,1.933) (1.491,1.935) (1.526,1.936) (1.561,1.938) (1.596,1.940) (1.631,1.941) (1.666,1.943) (1.701,1.944) (1.736,1.946) (1.771,1.947) (1.806,1.949) (1.841,1.950) (1.877,1.951) (1.912,1.953) (1.947,1.954) (1.982,1.955) (2.017,1.956) (2.052,1.958) (2.087,1.959) (2.122,1.960) (2.157,1.961) (2.192,1.962) (2.228,1.963) (2.263,1.964) (2.298,1.965) (2.333,1.966) (2.368,1.967) (2.403,1.967) (2.438,1.968) (2.473,1.969) (2.508,1.970) (2.544,1.971) (2.579,1.971) (2.614,1.972) (2.649,1.973) (2.684,1.973) (2.719,1.974) (2.754,1.975) (2.789,1.975) (2.824,1.976) (2.860,1.977) (2.895,1.977) (2.930,1.978) (2.965,1.978) (3.000,1.979)};
\addplot[lineaD, flechaen=0.142, flechaen=0.858] coordinates {(-3.000,2.179) (-2.965,2.178) (-2.930,2.178) (-2.895,2.177) (-2.860,2.177) (-2.824,2.176) (-2.789,2.176) (-2.754,2.175) (-2.719,2.174) (-2.684,2.174) (-2.649,2.173) (-2.614,2.173) (-2.579,2.172) (-2.544,2.171) (-2.509,2.170) (-2.473,2.170) (-2.438,2.169) (-2.403,2.168) (-2.368,2.167) (-2.333,2.167) (-2.298,2.166) (-2.263,2.165) (-2.228,2.164) (-2.193,2.163) (-2.158,2.162) (-2.123,2.161) (-2.087,2.161) (-2.052,2.160) (-2.017,2.159) (-1.982,2.158) (-1.947,2.156) (-1.912,2.155) (-1.877,2.154) (-1.842,2.153) (-1.807,2.152) (-1.772,2.151) (-1.737,2.150) (-1.701,2.148) (-1.666,2.147) (-1.631,2.146) (-1.596,2.145) (-1.561,2.143) (-1.526,2.142) (-1.491,2.141) (-1.456,2.139) (-1.421,2.138) (-1.386,2.136) (-1.351,2.135) (-1.316,2.133) (-1.281,2.132) (-1.245,2.130) (-1.210,2.129) (-1.175,2.127) (-1.140,2.126) (-1.105,2.124) (-1.070,2.122) (-1.035,2.121) (-1.000,2.119) (-0.965,2.118) (-0.930,2.116) (-0.895,2.114) (-0.860,2.113) (-0.825,2.111) (-0.790,2.110) (-0.754,2.108) (-0.719,2.107) (-0.684,2.105) (-0.649,2.104) (-0.614,2.102) (-0.579,2.101) (-0.544,2.100) (-0.509,2.098) (-0.474,2.097) (-0.439,2.096) (-0.404,2.095) (-0.369,2.094) (-0.333,2.093) (-0.298,2.092) (-0.263,2.091) (-0.228,2.090) (-0.193,2.090) (-0.158,2.089) (-0.123,2.089) (-0.088,2.089) (-0.053,2.088) (-0.018,2.088) (0.018,2.088) (0.053,2.088) (0.088,2.089) (0.123,2.089) (0.158,2.089) (0.193,2.090) (0.228,2.090) (0.263,2.091) (0.298,2.092) (0.333,2.093) (0.369,2.094) (0.404,2.095) (0.439,2.096) (0.474,2.097) (0.509,2.098) (0.544,2.100) (0.579,2.101) (0.614,2.102) (0.649,2.104) (0.684,2.105) (0.719,2.107) (0.754,2.108) (0.790,2.110) (0.825,2.111) (0.860,2.113) (0.895,2.114) (0.930,2.116) (0.965,2.118) (1.000,2.119) (1.035,2.121) (1.070,2.122) (1.105,2.124) (1.140,2.126) (1.175,2.127) (1.210,2.129) (1.245,2.130) (1.281,2.132) (1.316,2.133) (1.351,2.135) (1.386,2.136) (1.421,2.138) (1.456,2.139) (1.491,2.141) (1.526,2.142) (1.561,2.143) (1.596,2.145) (1.631,2.146) (1.666,2.147) (1.701,2.148) (1.737,2.150) (1.772,2.151) (1.807,2.152) (1.842,2.153) (1.877,2.154) (1.912,2.155) (1.947,2.156) (1.982,2.158) (2.017,2.159) (2.052,2.160) (2.087,2.161) (2.123,2.161) (2.158,2.162) (2.193,2.163) (2.228,2.164) (2.263,2.165) (2.298,2.166) (2.333,2.167) (2.368,2.167) (2.403,2.168) (2.438,2.169) (2.473,2.170) (2.509,2.170) (2.544,2.171) (2.579,2.172) (2.614,2.173) (2.649,2.173) (2.684,2.174) (2.719,2.174) (2.754,2.175) (2.789,2.176) (2.824,2.176) (2.860,2.177) (2.895,2.177) (2.930,2.178) (2.965,2.178) (3.000,2.179)};
\addplot[lineaD, flechaen=0.142, flechaen=0.858] coordinates {(-3.000,-0.196) (-2.965,-0.196) (-2.930,-0.196) (-2.895,-0.196) (-2.860,-0.196) (-2.824,-0.196) (-2.789,-0.196) (-2.754,-0.195) (-2.719,-0.195) (-2.684,-0.195) (-2.649,-0.195) (-2.614,-0.195) (-2.579,-0.194) (-2.544,-0.194) (-2.509,-0.194) (-2.473,-0.194) (-2.438,-0.193) (-2.403,-0.193) (-2.368,-0.193) (-2.333,-0.193) (-2.298,-0.192) (-2.263,-0.192) (-2.228,-0.192) (-2.193,-0.191) (-2.158,-0.191) (-2.122,-0.190) (-2.087,-0.190) (-2.052,-0.189) (-2.017,-0.189) (-1.982,-0.188) (-1.947,-0.188) (-1.912,-0.187) (-1.877,-0.187) (-1.842,-0.186) (-1.806,-0.185) (-1.771,-0.184) (-1.736,-0.183) (-1.701,-0.183) (-1.666,-0.182) (-1.631,-0.181) (-1.596,-0.179) (-1.561,-0.178) (-1.526,-0.177) (-1.491,-0.176) (-1.456,-0.174) (-1.421,-0.173) (-1.385,-0.171) (-1.350,-0.169) (-1.315,-0.167) (-1.280,-0.165) (-1.245,-0.163) (-1.210,-0.161) (-1.175,-0.158) (-1.140,-0.155) (-1.105,-0.152) (-1.070,-0.149) (-1.035,-0.146) (-1.000,-0.143) (-0.965,-0.141) (-0.930,-0.141) (-0.895,-0.141) (-0.860,-0.141) (-0.825,-0.141) (-0.790,-0.141) (-0.755,-0.141) (-0.720,-0.141) (-0.685,-0.141) (-0.649,-0.141) (-0.614,-0.141) (-0.579,-0.141) (-0.544,-0.141) (-0.509,-0.141) (-0.474,-0.141) (-0.439,-0.141) (-0.404,-0.141) (-0.369,-0.141) (-0.333,-0.141) (-0.298,-0.141) (-0.263,-0.141) (-0.228,-0.141) (-0.193,-0.141) (-0.158,-0.141) (-0.123,-0.141) (-0.088,-0.141) (-0.053,-0.141) (-0.018,-0.141) (0.018,-0.141) (0.053,-0.141) (0.088,-0.141) (0.123,-0.141) (0.158,-0.141) (0.193,-0.141) (0.228,-0.141) (0.263,-0.141) (0.298,-0.141) (0.333,-0.141) (0.369,-0.141) (0.404,-0.141) (0.439,-0.141) (0.474,-0.141) (0.509,-0.141) (0.544,-0.141) (0.579,-0.141) (0.614,-0.141) (0.649,-0.141) (0.685,-0.141) (0.720,-0.141) (0.755,-0.141) (0.790,-0.141) (0.825,-0.141) (0.860,-0.141) (0.895,-0.141) (0.930,-0.141) (0.965,-0.141) (1.000,-0.143) (1.035,-0.146) (1.070,-0.149) (1.105,-0.152) (1.140,-0.155) (1.175,-0.158) (1.210,-0.161) (1.245,-0.163) (1.280,-0.165) (1.315,-0.167) (1.350,-0.169) (1.385,-0.171) (1.421,-0.173) (1.456,-0.174) (1.491,-0.176) (1.526,-0.177) (1.561,-0.178) (1.596,-0.179) (1.631,-0.181) (1.666,-0.182) (1.701,-0.183) (1.736,-0.183) (1.771,-0.184) (1.806,-0.185) (1.842,-0.186) (1.877,-0.187) (1.912,-0.187) (1.947,-0.188) (1.982,-0.188) (2.017,-0.189) (2.052,-0.189) (2.087,-0.190) (2.122,-0.190) (2.158,-0.191) (2.193,-0.191) (2.228,-0.192) (2.263,-0.192) (2.298,-0.192) (2.333,-0.193) (2.368,-0.193) (2.403,-0.193) (2.438,-0.193) (2.473,-0.194) (2.509,-0.194) (2.544,-0.194) (2.579,-0.194) (2.614,-0.195) (2.649,-0.195) (2.684,-0.195) (2.719,-0.195) (2.754,-0.195) (2.789,-0.196) (2.824,-0.196) (2.860,-0.196) (2.895,-0.196) (2.930,-0.196) (2.965,-0.196) (3.000,-0.196)};
\addplot[lineaD, flechaen=0.141, flechaen=0.859] coordinates {(-3.000,-0.393) (-2.965,-0.393) (-2.930,-0.392) (-2.895,-0.392) (-2.859,-0.392) (-2.824,-0.392) (-2.789,-0.391) (-2.754,-0.391) (-2.719,-0.391) (-2.684,-0.390) (-2.648,-0.390) (-2.613,-0.390) (-2.578,-0.389) (-2.543,-0.389) (-2.508,-0.388) (-2.473,-0.388) (-2.438,-0.387) (-2.402,-0.387) (-2.367,-0.386) (-2.332,-0.386) (-2.297,-0.385) (-2.262,-0.384) (-2.227,-0.384) (-2.191,-0.383) (-2.156,-0.382) (-2.121,-0.381) (-2.086,-0.381) (-2.051,-0.380) (-2.016,-0.379) (-1.981,-0.378) (-1.945,-0.377) (-1.910,-0.375) (-1.875,-0.374) (-1.840,-0.373) (-1.805,-0.371) (-1.770,-0.370) (-1.735,-0.368) (-1.700,-0.367) (-1.664,-0.365) (-1.629,-0.363) (-1.594,-0.361) (-1.559,-0.359) (-1.524,-0.356) (-1.489,-0.354) (-1.454,-0.351) (-1.419,-0.348) (-1.384,-0.345) (-1.349,-0.342) (-1.314,-0.338) (-1.279,-0.334) (-1.244,-0.330) (-1.209,-0.326) (-1.174,-0.321) (-1.140,-0.316) (-1.105,-0.310) (-1.070,-0.304) (-1.036,-0.298) (-1.001,-0.292) (-0.967,-0.284) (-0.932,-0.283) (-0.896,-0.283) (-0.861,-0.283) (-0.826,-0.283) (-0.791,-0.283) (-0.756,-0.283) (-0.721,-0.283) (-0.686,-0.283) (-0.650,-0.283) (-0.615,-0.283) (-0.580,-0.283) (-0.545,-0.283) (-0.510,-0.283) (-0.475,-0.283) (-0.439,-0.283) (-0.404,-0.283) (-0.369,-0.283) (-0.334,-0.283) (-0.299,-0.283) (-0.264,-0.283) (-0.229,-0.283) (-0.193,-0.283) (-0.158,-0.283) (-0.123,-0.283) (-0.088,-0.283) (-0.053,-0.283) (-0.018,-0.283) (0.018,-0.283) (0.053,-0.283) (0.088,-0.283) (0.123,-0.283) (0.158,-0.283) (0.193,-0.283) (0.229,-0.283) (0.264,-0.283) (0.299,-0.283) (0.334,-0.283) (0.369,-0.283) (0.404,-0.283) (0.439,-0.283) (0.475,-0.283) (0.510,-0.283) (0.545,-0.283) (0.580,-0.283) (0.615,-0.283) (0.650,-0.283) (0.686,-0.283) (0.721,-0.283) (0.756,-0.283) (0.791,-0.283) (0.826,-0.283) (0.861,-0.283) (0.896,-0.283) (0.932,-0.283) (0.967,-0.284) (1.001,-0.292) (1.036,-0.298) (1.070,-0.304) (1.105,-0.310) (1.140,-0.316) (1.174,-0.321) (1.209,-0.326) (1.244,-0.330) (1.279,-0.334) (1.314,-0.338) (1.349,-0.342) (1.384,-0.345) (1.419,-0.348) (1.454,-0.351) (1.489,-0.354) (1.524,-0.356) (1.559,-0.359) (1.594,-0.361) (1.629,-0.363) (1.664,-0.365) (1.700,-0.367) (1.735,-0.368) (1.770,-0.370) (1.805,-0.371) (1.840,-0.373) (1.875,-0.374) (1.910,-0.375) (1.945,-0.377) (1.981,-0.378) (2.016,-0.379) (2.051,-0.380) (2.086,-0.381) (2.121,-0.381) (2.156,-0.382) (2.191,-0.383) (2.227,-0.384) (2.262,-0.384) (2.297,-0.385) (2.332,-0.386) (2.367,-0.386) (2.402,-0.387) (2.438,-0.387) (2.473,-0.388) (2.508,-0.388) (2.543,-0.389) (2.578,-0.389) (2.613,-0.390) (2.648,-0.390) (2.684,-0.390) (2.719,-0.391) (2.754,-0.391) (2.789,-0.391) (2.824,-0.392) (2.859,-0.392) (2.895,-0.392) (2.930,-0.392) (2.965,-0.393) (3.000,-0.393)};
\addplot[lineaD, flechaen=0.141, flechaen=0.859] coordinates {(-3.000,-0.590) (-2.965,-0.589) (-2.930,-0.589) (-2.895,-0.589) (-2.860,-0.588) (-2.825,-0.588) (-2.790,-0.587) (-2.755,-0.587) (-2.720,-0.587) (-2.685,-0.586) (-2.650,-0.586) (-2.615,-0.585) (-2.580,-0.584) (-2.545,-0.584) (-2.510,-0.583) (-2.474,-0.583) (-2.439,-0.582) (-2.404,-0.581) (-2.369,-0.580) (-2.334,-0.580) (-2.299,-0.579) (-2.264,-0.578) (-2.229,-0.577) (-2.194,-0.576) (-2.159,-0.575) (-2.124,-0.574) (-2.089,-0.572) (-2.054,-0.571) (-2.019,-0.570) (-1.984,-0.569) (-1.949,-0.567) (-1.914,-0.565) (-1.879,-0.564) (-1.844,-0.562) (-1.809,-0.560) (-1.774,-0.558) (-1.739,-0.556) (-1.704,-0.554) (-1.669,-0.551) (-1.634,-0.549) (-1.599,-0.546) (-1.565,-0.543) (-1.530,-0.540) (-1.495,-0.537) (-1.460,-0.533) (-1.425,-0.529) (-1.390,-0.525) (-1.356,-0.521) (-1.321,-0.516) (-1.286,-0.511) (-1.251,-0.506) (-1.217,-0.500) (-1.182,-0.494) (-1.148,-0.487) (-1.114,-0.480) (-1.080,-0.473) (-1.045,-0.465) (-1.011,-0.456) (-0.978,-0.447) (-0.944,-0.437) (-0.911,-0.426) (-0.876,-0.424) (-0.841,-0.424) (-0.806,-0.424) (-0.771,-0.424) (-0.736,-0.424) (-0.701,-0.424) (-0.666,-0.424) (-0.631,-0.424) (-0.596,-0.424) (-0.561,-0.424) (-0.526,-0.424) (-0.491,-0.424) (-0.455,-0.424) (-0.420,-0.424) (-0.385,-0.424) (-0.350,-0.424) (-0.315,-0.424) (-0.280,-0.424) (-0.245,-0.424) (-0.210,-0.424) (-0.175,-0.424) (-0.140,-0.424) (-0.105,-0.424) (-0.070,-0.424) (-0.035,-0.424) (0.000,-0.424) (0.035,-0.424) (0.070,-0.424) (0.105,-0.424) (0.140,-0.424) (0.175,-0.424) (0.210,-0.424) (0.245,-0.424) (0.280,-0.424) (0.315,-0.424) (0.350,-0.424) (0.385,-0.424) (0.420,-0.424) (0.455,-0.424) (0.491,-0.424) (0.526,-0.424) (0.561,-0.424) (0.596,-0.424) (0.631,-0.424) (0.666,-0.424) (0.701,-0.424) (0.736,-0.424) (0.771,-0.424) (0.806,-0.424) (0.841,-0.424) (0.876,-0.424) (0.911,-0.426) (0.944,-0.437) (0.978,-0.447) (1.011,-0.456) (1.045,-0.465) (1.080,-0.473) (1.114,-0.480) (1.148,-0.487) (1.182,-0.494) (1.217,-0.500) (1.251,-0.506) (1.286,-0.511) (1.321,-0.516) (1.356,-0.521) (1.390,-0.525) (1.425,-0.529) (1.460,-0.533) (1.495,-0.537) (1.530,-0.540) (1.565,-0.543) (1.599,-0.546) (1.634,-0.549) (1.669,-0.551) (1.704,-0.554) (1.739,-0.556) (1.774,-0.558) (1.809,-0.560) (1.844,-0.562) (1.879,-0.564) (1.914,-0.565) (1.949,-0.567) (1.984,-0.569) (2.019,-0.570) (2.054,-0.571) (2.089,-0.572) (2.124,-0.574) (2.159,-0.575) (2.194,-0.576) (2.229,-0.577) (2.264,-0.578) (2.299,-0.579) (2.334,-0.580) (2.369,-0.580) (2.404,-0.581) (2.439,-0.582) (2.474,-0.583) (2.510,-0.583) (2.545,-0.584) (2.580,-0.584) (2.615,-0.585) (2.650,-0.586) (2.685,-0.586) (2.720,-0.587) (2.755,-0.587) (2.790,-0.587) (2.825,-0.588) (2.860,-0.588) (2.895,-0.589) (2.930,-0.589) (2.965,-0.589) (3.000,-0.590)};
\addplot[lineaD, flechaen=0.141, flechaen=0.859] coordinates {(-3.000,-0.787) (-2.965,-0.786) (-2.930,-0.786) (-2.895,-0.786) (-2.859,-0.785) (-2.824,-0.785) (-2.789,-0.784) (-2.754,-0.783) (-2.719,-0.783) (-2.684,-0.782) (-2.649,-0.782) (-2.613,-0.781) (-2.578,-0.780) (-2.543,-0.780) (-2.508,-0.779) (-2.473,-0.778) (-2.438,-0.777) (-2.402,-0.776) (-2.367,-0.775) (-2.332,-0.774) (-2.297,-0.773) (-2.262,-0.772) (-2.227,-0.771) (-2.192,-0.770) (-2.157,-0.769) (-2.121,-0.767) (-2.086,-0.766) (-2.051,-0.764) (-2.016,-0.763) (-1.981,-0.761) (-1.946,-0.759) (-1.911,-0.758) (-1.876,-0.756) (-1.841,-0.754) (-1.805,-0.751) (-1.770,-0.749) (-1.735,-0.747) (-1.700,-0.744) (-1.665,-0.741) (-1.630,-0.738) (-1.595,-0.735) (-1.560,-0.732) (-1.525,-0.728) (-1.490,-0.724) (-1.455,-0.720) (-1.420,-0.716) (-1.386,-0.711) (-1.351,-0.706) (-1.316,-0.701) (-1.281,-0.696) (-1.247,-0.690) (-1.212,-0.683) (-1.178,-0.676) (-1.143,-0.669) (-1.109,-0.661) (-1.075,-0.653) (-1.041,-0.643) (-1.007,-0.634) (-0.974,-0.623) (-0.940,-0.612) (-0.907,-0.600) (-0.874,-0.587) (-0.842,-0.574) (-0.809,-0.566) (-0.773,-0.566) (-0.738,-0.566) (-0.703,-0.566) (-0.668,-0.566) (-0.633,-0.566) (-0.598,-0.566) (-0.562,-0.566) (-0.527,-0.566) (-0.492,-0.566) (-0.457,-0.566) (-0.422,-0.566) (-0.387,-0.566) (-0.352,-0.566) (-0.316,-0.566) (-0.281,-0.566) (-0.246,-0.566) (-0.211,-0.566) (-0.176,-0.566) (-0.141,-0.566) (-0.105,-0.566) (-0.070,-0.566) (-0.035,-0.566) (0.000,-0.566) (0.035,-0.566) (0.070,-0.566) (0.105,-0.566) (0.141,-0.566) (0.176,-0.566) (0.211,-0.566) (0.246,-0.566) (0.281,-0.566) (0.316,-0.566) (0.352,-0.566) (0.387,-0.566) (0.422,-0.566) (0.457,-0.566) (0.492,-0.566) (0.527,-0.566) (0.562,-0.566) (0.598,-0.566) (0.633,-0.566) (0.668,-0.566) (0.703,-0.566) (0.738,-0.566) (0.773,-0.566) (0.809,-0.566) (0.842,-0.574) (0.874,-0.587) (0.907,-0.600) (0.940,-0.612) (0.974,-0.623) (1.007,-0.634) (1.041,-0.643) (1.075,-0.653) (1.109,-0.661) (1.143,-0.669) (1.178,-0.676) (1.212,-0.683) (1.247,-0.690) (1.281,-0.696) (1.316,-0.701) (1.351,-0.706) (1.386,-0.711) (1.420,-0.716) (1.455,-0.720) (1.490,-0.724) (1.525,-0.728) (1.560,-0.732) (1.595,-0.735) (1.630,-0.738) (1.665,-0.741) (1.700,-0.744) (1.735,-0.747) (1.770,-0.749) (1.805,-0.751) (1.841,-0.754) (1.876,-0.756) (1.911,-0.758) (1.946,-0.759) (1.981,-0.761) (2.016,-0.763) (2.051,-0.764) (2.086,-0.766) (2.121,-0.767) (2.157,-0.769) (2.192,-0.770) (2.227,-0.771) (2.262,-0.772) (2.297,-0.773) (2.332,-0.774) (2.367,-0.775) (2.402,-0.776) (2.438,-0.777) (2.473,-0.778) (2.508,-0.779) (2.543,-0.780) (2.578,-0.780) (2.613,-0.781) (2.649,-0.782) (2.684,-0.782) (2.719,-0.783) (2.754,-0.783) (2.789,-0.784) (2.824,-0.785) (2.859,-0.785) (2.895,-0.786) (2.930,-0.786) (2.965,-0.786) (3.000,-0.787)};
\addplot[lineaD, flechaen=0.140, flechaen=0.860] coordinates {(-3.000,-0.984) (-2.965,-0.984) (-2.930,-0.983) (-2.895,-0.983) (-2.860,-0.982) (-2.825,-0.982) (-2.789,-0.981) (-2.754,-0.981) (-2.719,-0.980) (-2.684,-0.979) (-2.649,-0.979) (-2.614,-0.978) (-2.579,-0.977) (-2.544,-0.976) (-2.509,-0.975) (-2.474,-0.974) (-2.439,-0.973) (-2.404,-0.972) (-2.368,-0.971) (-2.333,-0.970) (-2.298,-0.969) (-2.263,-0.968) (-2.228,-0.967) (-2.193,-0.966) (-2.158,-0.964) (-2.123,-0.963) (-2.088,-0.961) (-2.053,-0.960) (-2.018,-0.958) (-1.983,-0.956) (-1.948,-0.954) (-1.913,-0.952) (-1.878,-0.950) (-1.843,-0.948) (-1.808,-0.946) (-1.773,-0.943) (-1.738,-0.941) (-1.703,-0.938) (-1.668,-0.935) (-1.633,-0.932) (-1.598,-0.929) (-1.563,-0.926) (-1.528,-0.922) (-1.493,-0.918) (-1.458,-0.914) (-1.423,-0.910) (-1.389,-0.906) (-1.354,-0.901) (-1.319,-0.896) (-1.284,-0.890) (-1.250,-0.884) (-1.215,-0.878) (-1.181,-0.872) (-1.146,-0.865) (-1.112,-0.857) (-1.078,-0.849) (-1.044,-0.841) (-1.010,-0.831) (-0.976,-0.822) (-0.943,-0.811) (-0.910,-0.800) (-0.877,-0.788) (-0.844,-0.775) (-0.812,-0.761) (-0.780,-0.746) (-0.749,-0.731) (-0.718,-0.713) (-0.684,-0.707) (-0.649,-0.707) (-0.614,-0.707) (-0.579,-0.707) (-0.544,-0.707) (-0.509,-0.707) (-0.474,-0.707) (-0.439,-0.707) (-0.404,-0.707) (-0.368,-0.707) (-0.333,-0.707) (-0.298,-0.707) (-0.263,-0.707) (-0.228,-0.707) (-0.193,-0.707) (-0.158,-0.707) (-0.123,-0.707) (-0.088,-0.707) (-0.053,-0.707) (-0.018,-0.707) (0.018,-0.707) (0.053,-0.707) (0.088,-0.707) (0.123,-0.707) (0.158,-0.707) (0.193,-0.707) (0.228,-0.707) (0.263,-0.707) (0.298,-0.707) (0.333,-0.707) (0.368,-0.707) (0.404,-0.707) (0.439,-0.707) (0.474,-0.707) (0.509,-0.707) (0.544,-0.707) (0.579,-0.707) (0.614,-0.707) (0.649,-0.707) (0.684,-0.707) (0.718,-0.713) (0.749,-0.731) (0.780,-0.746) (0.812,-0.761) (0.844,-0.775) (0.877,-0.788) (0.910,-0.800) (0.943,-0.811) (0.976,-0.822) (1.010,-0.831) (1.044,-0.841) (1.078,-0.849) (1.112,-0.857) (1.146,-0.865) (1.181,-0.872) (1.215,-0.878) (1.250,-0.884) (1.284,-0.890) (1.319,-0.896) (1.354,-0.901) (1.389,-0.906) (1.423,-0.910) (1.458,-0.914) (1.493,-0.918) (1.528,-0.922) (1.563,-0.926) (1.598,-0.929) (1.633,-0.932) (1.668,-0.935) (1.703,-0.938) (1.738,-0.941) (1.773,-0.943) (1.808,-0.946) (1.843,-0.948) (1.878,-0.950) (1.913,-0.952) (1.948,-0.954) (1.983,-0.956) (2.018,-0.958) (2.053,-0.960) (2.088,-0.961) (2.123,-0.963) (2.158,-0.964) (2.193,-0.966) (2.228,-0.967) (2.263,-0.968) (2.298,-0.969) (2.333,-0.970) (2.368,-0.971) (2.404,-0.972) (2.439,-0.973) (2.474,-0.974) (2.509,-0.975) (2.544,-0.976) (2.579,-0.977) (2.614,-0.978) (2.649,-0.979) (2.684,-0.979) (2.719,-0.980) (2.754,-0.981) (2.789,-0.981) (2.825,-0.982) (2.860,-0.982) (2.895,-0.983) (2.930,-0.983) (2.965,-0.984) (3.000,-0.984)};
\addplot[lineaD, flechaen=0.139, flechaen=0.861] coordinates {(-3.000,-1.182) (-2.965,-1.182) (-2.930,-1.181) (-2.895,-1.181) (-2.860,-1.180) (-2.825,-1.180) (-2.790,-1.179) (-2.755,-1.178) (-2.720,-1.178) (-2.685,-1.177) (-2.650,-1.176) (-2.615,-1.175) (-2.580,-1.174) (-2.545,-1.174) (-2.510,-1.173) (-2.475,-1.172) (-2.440,-1.171) (-2.404,-1.170) (-2.369,-1.169) (-2.334,-1.168) (-2.299,-1.166) (-2.264,-1.165) (-2.229,-1.164) (-2.194,-1.163) (-2.159,-1.161) (-2.124,-1.160) (-2.089,-1.158) (-2.054,-1.157) (-2.019,-1.155) (-1.984,-1.153) (-1.949,-1.151) (-1.914,-1.149) (-1.879,-1.147) (-1.844,-1.145) (-1.809,-1.143) (-1.774,-1.141) (-1.740,-1.138) (-1.705,-1.136) (-1.670,-1.133) (-1.635,-1.130) (-1.600,-1.127) (-1.565,-1.124) (-1.530,-1.121) (-1.495,-1.117) (-1.460,-1.113) (-1.426,-1.110) (-1.391,-1.105) (-1.356,-1.101) (-1.321,-1.097) (-1.287,-1.092) (-1.252,-1.087) (-1.217,-1.081) (-1.183,-1.076) (-1.148,-1.070) (-1.114,-1.063) (-1.079,-1.056) (-1.045,-1.049) (-1.011,-1.042) (-0.977,-1.034) (-0.943,-1.025) (-0.909,-1.016) (-0.875,-1.006) (-0.842,-0.996) (-0.808,-0.985) (-0.775,-0.974) (-0.742,-0.961) (-0.710,-0.948) (-0.678,-0.934) (-0.646,-0.919) (-0.615,-0.903) (-0.585,-0.885) (-0.555,-0.867) (-0.526,-0.849) (-0.491,-0.849) (-0.456,-0.849) (-0.420,-0.849) (-0.385,-0.849) (-0.350,-0.849) (-0.315,-0.849) (-0.280,-0.849) (-0.245,-0.849) (-0.210,-0.849) (-0.175,-0.849) (-0.140,-0.849) (-0.105,-0.849) (-0.070,-0.849) (-0.035,-0.849) (0.000,-0.849) (0.035,-0.849) (0.070,-0.849) (0.105,-0.849) (0.140,-0.849) (0.175,-0.849) (0.210,-0.849) (0.245,-0.849) (0.280,-0.849) (0.315,-0.849) (0.350,-0.849) (0.385,-0.849) (0.420,-0.849) (0.456,-0.849) (0.491,-0.849) (0.526,-0.849) (0.555,-0.867) (0.585,-0.885) (0.615,-0.903) (0.646,-0.919) (0.678,-0.934) (0.710,-0.948) (0.742,-0.961) (0.775,-0.974) (0.808,-0.985) (0.842,-0.996) (0.875,-1.006) (0.909,-1.016) (0.943,-1.025) (0.977,-1.034) (1.011,-1.042) (1.045,-1.049) (1.079,-1.056) (1.114,-1.063) (1.148,-1.070) (1.183,-1.076) (1.217,-1.081) (1.252,-1.087) (1.287,-1.092) (1.321,-1.097) (1.356,-1.101) (1.391,-1.105) (1.426,-1.110) (1.460,-1.113) (1.495,-1.117) (1.530,-1.121) (1.565,-1.124) (1.600,-1.127) (1.635,-1.130) (1.670,-1.133) (1.705,-1.136) (1.740,-1.138) (1.774,-1.141) (1.809,-1.143) (1.844,-1.145) (1.879,-1.147) (1.914,-1.149) (1.949,-1.151) (1.984,-1.153) (2.019,-1.155) (2.054,-1.157) (2.089,-1.158) (2.124,-1.160) (2.159,-1.161) (2.194,-1.163) (2.229,-1.164) (2.264,-1.165) (2.299,-1.166) (2.334,-1.168) (2.369,-1.169) (2.404,-1.170) (2.440,-1.171) (2.475,-1.172) (2.510,-1.173) (2.545,-1.174) (2.580,-1.174) (2.615,-1.175) (2.650,-1.176) (2.685,-1.177) (2.720,-1.178) (2.755,-1.178) (2.790,-1.179) (2.825,-1.180) (2.860,-1.180) (2.895,-1.181) (2.930,-1.181) (2.965,-1.182) (3.000,-1.182)};
\addplot[lineaD, flechaen=0.139, flechaen=0.861] coordinates {(-3.000,-1.381) (-2.965,-1.380) (-2.930,-1.380) (-2.895,-1.379) (-2.860,-1.379) (-2.825,-1.378) (-2.790,-1.377) (-2.755,-1.377) (-2.720,-1.376) (-2.685,-1.375) (-2.650,-1.374) (-2.615,-1.374) (-2.580,-1.373) (-2.545,-1.372) (-2.510,-1.371) (-2.475,-1.370) (-2.439,-1.369) (-2.404,-1.368) (-2.369,-1.367) (-2.334,-1.366) (-2.299,-1.365) (-2.264,-1.364) (-2.229,-1.362) (-2.194,-1.361) (-2.159,-1.360) (-2.124,-1.358) (-2.089,-1.357) (-2.054,-1.355) (-2.019,-1.354) (-1.984,-1.352) (-1.949,-1.350) (-1.914,-1.348) (-1.879,-1.347) (-1.844,-1.345) (-1.809,-1.343) (-1.774,-1.340) (-1.739,-1.338) (-1.704,-1.336) (-1.669,-1.333) (-1.635,-1.331) (-1.600,-1.328) (-1.565,-1.325) (-1.530,-1.322) (-1.495,-1.319) (-1.460,-1.316) (-1.425,-1.313) (-1.390,-1.309) (-1.355,-1.305) (-1.321,-1.302) (-1.286,-1.298) (-1.251,-1.293) (-1.216,-1.289) (-1.181,-1.284) (-1.147,-1.280) (-1.112,-1.274) (-1.077,-1.269) (-1.043,-1.264) (-1.008,-1.258) (-0.974,-1.252) (-0.939,-1.246) (-0.905,-1.239) (-0.871,-1.232) (-0.836,-1.225) (-0.802,-1.217) (-0.768,-1.209) (-0.734,-1.201) (-0.700,-1.193) (-0.666,-1.184) (-0.632,-1.174) (-0.599,-1.164) (-0.565,-1.154) (-0.532,-1.144) (-0.498,-1.133) (-0.465,-1.121) (-0.432,-1.109) (-0.399,-1.097) (-0.367,-1.085) (-0.334,-1.071) (-0.302,-1.058) (-0.270,-1.044) (-0.237,-1.031) (-0.205,-1.017) (-0.173,-1.003) (-0.140,-0.991) (-0.105,-0.990) (-0.070,-0.990) (-0.035,-0.990) (-0.000,-0.990) (0.035,-0.990) (0.070,-0.990) (0.105,-0.990) (0.140,-0.991) (0.173,-1.003) (0.205,-1.017) (0.237,-1.031) (0.270,-1.044) (0.302,-1.058) (0.334,-1.071) (0.367,-1.085) (0.399,-1.097) (0.432,-1.109) (0.465,-1.121) (0.498,-1.133) (0.532,-1.144) (0.565,-1.154) (0.599,-1.164) (0.632,-1.174) (0.666,-1.184) (0.700,-1.193) (0.734,-1.201) (0.768,-1.209) (0.802,-1.217) (0.836,-1.225) (0.871,-1.232) (0.905,-1.239) (0.939,-1.246) (0.974,-1.252) (1.008,-1.258) (1.043,-1.264) (1.077,-1.269) (1.112,-1.274) (1.147,-1.280) (1.181,-1.284) (1.216,-1.289) (1.251,-1.293) (1.286,-1.298) (1.321,-1.302) (1.355,-1.305) (1.390,-1.309) (1.425,-1.313) (1.460,-1.316) (1.495,-1.319) (1.530,-1.322) (1.565,-1.325) (1.600,-1.328) (1.635,-1.331) (1.669,-1.333) (1.704,-1.336) (1.739,-1.338) (1.774,-1.340) (1.809,-1.343) (1.844,-1.345) (1.879,-1.347) (1.914,-1.348) (1.949,-1.350) (1.984,-1.352) (2.019,-1.354) (2.054,-1.355) (2.089,-1.357) (2.124,-1.358) (2.159,-1.360) (2.194,-1.361) (2.229,-1.362) (2.264,-1.364) (2.299,-1.365) (2.334,-1.366) (2.369,-1.367) (2.404,-1.368) (2.439,-1.369) (2.475,-1.370) (2.510,-1.371) (2.545,-1.372) (2.580,-1.373) (2.615,-1.374) (2.650,-1.374) (2.685,-1.375) (2.720,-1.376) (2.755,-1.377) (2.790,-1.377) (2.825,-1.378) (2.860,-1.379) (2.895,-1.379) (2.930,-1.380) (2.965,-1.380) (3.000,-1.381)};
\addplot[lineaD, flechaen=0.141, flechaen=0.859] coordinates {(-3.000,-1.580) (-2.965,-1.579) (-2.930,-1.579) (-2.895,-1.578) (-2.860,-1.578) (-2.825,-1.577) (-2.790,-1.576) (-2.755,-1.576) (-2.720,-1.575) (-2.685,-1.574) (-2.650,-1.573) (-2.615,-1.573) (-2.580,-1.572) (-2.545,-1.571) (-2.510,-1.570) (-2.475,-1.569) (-2.440,-1.568) (-2.405,-1.567) (-2.370,-1.566) (-2.335,-1.565) (-2.300,-1.564) (-2.265,-1.563) (-2.230,-1.562) (-2.195,-1.561) (-2.160,-1.559) (-2.125,-1.558) (-2.090,-1.557) (-2.055,-1.555) (-2.020,-1.554) (-1.985,-1.552) (-1.950,-1.551) (-1.915,-1.549) (-1.880,-1.547) (-1.845,-1.546) (-1.810,-1.544) (-1.775,-1.542) (-1.740,-1.540) (-1.705,-1.538) (-1.670,-1.536) (-1.635,-1.533) (-1.600,-1.531) (-1.565,-1.529) (-1.530,-1.526) (-1.495,-1.524) (-1.460,-1.521) (-1.425,-1.518) (-1.390,-1.515) (-1.355,-1.512) (-1.320,-1.509) (-1.286,-1.506) (-1.251,-1.503) (-1.216,-1.499) (-1.181,-1.496) (-1.146,-1.492) (-1.111,-1.488) (-1.077,-1.484) (-1.042,-1.480) (-1.007,-1.476) (-0.972,-1.471) (-0.937,-1.467) (-0.903,-1.462) (-0.868,-1.458) (-0.833,-1.453) (-0.799,-1.448) (-0.764,-1.443) (-0.729,-1.438) (-0.695,-1.432) (-0.660,-1.427) (-0.626,-1.421) (-0.591,-1.416) (-0.556,-1.410) (-0.522,-1.405) (-0.487,-1.399) (-0.453,-1.394) (-0.418,-1.388) (-0.383,-1.383) (-0.349,-1.378) (-0.314,-1.373) (-0.279,-1.368) (-0.245,-1.364) (-0.210,-1.360) (-0.175,-1.357) (-0.140,-1.354) (-0.105,-1.352) (-0.070,-1.350) (-0.035,-1.349) (0.000,-1.348) (0.035,-1.349) (0.070,-1.350) (0.105,-1.352) (0.140,-1.354) (0.175,-1.357) (0.210,-1.360) (0.245,-1.364) (0.279,-1.368) (0.314,-1.373) (0.349,-1.378) (0.383,-1.383) (0.418,-1.388) (0.453,-1.394) (0.487,-1.399) (0.522,-1.405) (0.556,-1.410) (0.591,-1.416) (0.626,-1.421) (0.660,-1.427) (0.695,-1.432) (0.729,-1.438) (0.764,-1.443) (0.799,-1.448) (0.833,-1.453) (0.868,-1.458) (0.903,-1.462) (0.937,-1.467) (0.972,-1.471) (1.007,-1.476) (1.042,-1.480) (1.077,-1.484) (1.111,-1.488) (1.146,-1.492) (1.181,-1.496) (1.216,-1.499) (1.251,-1.503) (1.286,-1.506) (1.320,-1.509) (1.355,-1.512) (1.390,-1.515) (1.425,-1.518) (1.460,-1.521) (1.495,-1.524) (1.530,-1.526) (1.565,-1.529) (1.600,-1.531) (1.635,-1.533) (1.670,-1.536) (1.705,-1.538) (1.740,-1.540) (1.775,-1.542) (1.810,-1.544) (1.845,-1.546) (1.880,-1.547) (1.915,-1.549) (1.950,-1.551) (1.985,-1.552) (2.020,-1.554) (2.055,-1.555) (2.090,-1.557) (2.125,-1.558) (2.160,-1.559) (2.195,-1.561) (2.230,-1.562) (2.265,-1.563) (2.300,-1.564) (2.335,-1.565) (2.370,-1.566) (2.405,-1.567) (2.440,-1.568) (2.475,-1.569) (2.510,-1.570) (2.545,-1.571) (2.580,-1.572) (2.615,-1.573) (2.650,-1.573) (2.685,-1.574) (2.720,-1.575) (2.755,-1.576) (2.790,-1.576) (2.825,-1.577) (2.860,-1.578) (2.895,-1.578) (2.930,-1.579) (2.965,-1.579) (3.000,-1.580)};
\addplot[lineaD, flechaen=0.141, flechaen=0.859] coordinates {(-3.000,-1.779) (-2.965,-1.779) (-2.930,-1.778) (-2.895,-1.777) (-2.859,-1.777) (-2.824,-1.776) (-2.789,-1.776) (-2.754,-1.775) (-2.719,-1.774) (-2.684,-1.774) (-2.649,-1.773) (-2.613,-1.772) (-2.578,-1.771) (-2.543,-1.770) (-2.508,-1.770) (-2.473,-1.769) (-2.438,-1.768) (-2.403,-1.767) (-2.367,-1.766) (-2.332,-1.765) (-2.297,-1.764) (-2.262,-1.763) (-2.227,-1.762) (-2.192,-1.761) (-2.157,-1.760) (-2.121,-1.758) (-2.086,-1.757) (-2.051,-1.756) (-2.016,-1.755) (-1.981,-1.753) (-1.946,-1.752) (-1.911,-1.750) (-1.876,-1.749) (-1.841,-1.747) (-1.805,-1.746) (-1.770,-1.744) (-1.735,-1.742) (-1.700,-1.741) (-1.665,-1.739) (-1.630,-1.737) (-1.595,-1.735) (-1.560,-1.733) (-1.525,-1.731) (-1.489,-1.729) (-1.454,-1.727) (-1.419,-1.724) (-1.384,-1.722) (-1.349,-1.720) (-1.314,-1.717) (-1.279,-1.715) (-1.244,-1.712) (-1.209,-1.709) (-1.174,-1.707) (-1.139,-1.704) (-1.104,-1.701) (-1.069,-1.698) (-1.034,-1.695) (-0.999,-1.692) (-0.964,-1.689) (-0.929,-1.686) (-0.894,-1.682) (-0.859,-1.679) (-0.824,-1.676) (-0.789,-1.673) (-0.754,-1.669) (-0.719,-1.666) (-0.684,-1.663) (-0.649,-1.659) (-0.614,-1.656) (-0.579,-1.653) (-0.544,-1.649) (-0.509,-1.646) (-0.474,-1.643) (-0.439,-1.640) (-0.404,-1.637) (-0.369,-1.635) (-0.334,-1.632) (-0.299,-1.630) (-0.263,-1.628) (-0.228,-1.626) (-0.193,-1.624) (-0.158,-1.622) (-0.123,-1.621) (-0.088,-1.620) (-0.053,-1.620) (-0.018,-1.619) (0.018,-1.619) (0.053,-1.620) (0.088,-1.620) (0.123,-1.621) (0.158,-1.622) (0.193,-1.624) (0.228,-1.626) (0.263,-1.628) (0.299,-1.630) (0.334,-1.632) (0.369,-1.635) (0.404,-1.637) (0.439,-1.640) (0.474,-1.643) (0.509,-1.646) (0.544,-1.649) (0.579,-1.653) (0.614,-1.656) (0.649,-1.659) (0.684,-1.663) (0.719,-1.666) (0.754,-1.669) (0.789,-1.673) (0.824,-1.676) (0.859,-1.679) (0.894,-1.682) (0.929,-1.686) (0.964,-1.689) (0.999,-1.692) (1.034,-1.695) (1.069,-1.698) (1.104,-1.701) (1.139,-1.704) (1.174,-1.707) (1.209,-1.709) (1.244,-1.712) (1.279,-1.715) (1.314,-1.717) (1.349,-1.720) (1.384,-1.722) (1.419,-1.724) (1.454,-1.727) (1.489,-1.729) (1.525,-1.731) (1.560,-1.733) (1.595,-1.735) (1.630,-1.737) (1.665,-1.739) (1.700,-1.741) (1.735,-1.742) (1.770,-1.744) (1.805,-1.746) (1.841,-1.747) (1.876,-1.749) (1.911,-1.750) (1.946,-1.752) (1.981,-1.753) (2.016,-1.755) (2.051,-1.756) (2.086,-1.757) (2.121,-1.758) (2.157,-1.760) (2.192,-1.761) (2.227,-1.762) (2.262,-1.763) (2.297,-1.764) (2.332,-1.765) (2.367,-1.766) (2.403,-1.767) (2.438,-1.768) (2.473,-1.769) (2.508,-1.770) (2.543,-1.770) (2.578,-1.771) (2.613,-1.772) (2.649,-1.773) (2.684,-1.774) (2.719,-1.774) (2.754,-1.775) (2.789,-1.776) (2.824,-1.776) (2.859,-1.777) (2.895,-1.777) (2.930,-1.778) (2.965,-1.779) (3.000,-1.779)};
\addplot[lineaD, flechaen=0.142, flechaen=0.858] coordinates {(-3.000,-1.979) (-2.965,-1.978) (-2.930,-1.978) (-2.895,-1.977) (-2.860,-1.977) (-2.824,-1.976) (-2.789,-1.975) (-2.754,-1.975) (-2.719,-1.974) (-2.684,-1.973) (-2.649,-1.973) (-2.614,-1.972) (-2.579,-1.971) (-2.544,-1.971) (-2.508,-1.970) (-2.473,-1.969) (-2.438,-1.968) (-2.403,-1.967) (-2.368,-1.967) (-2.333,-1.966) (-2.298,-1.965) (-2.263,-1.964) (-2.228,-1.963) (-2.192,-1.962) (-2.157,-1.961) (-2.122,-1.960) (-2.087,-1.959) (-2.052,-1.958) (-2.017,-1.956) (-1.982,-1.955) (-1.947,-1.954) (-1.912,-1.953) (-1.877,-1.951) (-1.841,-1.950) (-1.806,-1.949) (-1.771,-1.947) (-1.736,-1.946) (-1.701,-1.944) (-1.666,-1.943) (-1.631,-1.941) (-1.596,-1.940) (-1.561,-1.938) (-1.526,-1.936) (-1.491,-1.935) (-1.455,-1.933) (-1.420,-1.931) (-1.385,-1.929) (-1.350,-1.927) (-1.315,-1.926) (-1.280,-1.924) (-1.245,-1.922) (-1.210,-1.920) (-1.175,-1.917) (-1.140,-1.915) (-1.105,-1.913) (-1.070,-1.911) (-1.035,-1.909) (-1.000,-1.907) (-0.965,-1.905) (-0.930,-1.902) (-0.895,-1.900) (-0.859,-1.898) (-0.824,-1.896) (-0.789,-1.893) (-0.754,-1.891) (-0.719,-1.889) (-0.684,-1.887) (-0.649,-1.885) (-0.614,-1.882) (-0.579,-1.880) (-0.544,-1.878) (-0.509,-1.876) (-0.474,-1.875) (-0.439,-1.873) (-0.404,-1.871) (-0.369,-1.869) (-0.334,-1.868) (-0.298,-1.867) (-0.263,-1.865) (-0.228,-1.864) (-0.193,-1.863) (-0.158,-1.862) (-0.123,-1.862) (-0.088,-1.861) (-0.053,-1.861) (-0.018,-1.861) (0.018,-1.861) (0.053,-1.861) (0.088,-1.861) (0.123,-1.862) (0.158,-1.862) (0.193,-1.863) (0.228,-1.864) (0.263,-1.865) (0.298,-1.867) (0.334,-1.868) (0.369,-1.869) (0.404,-1.871) (0.439,-1.873) (0.474,-1.875) (0.509,-1.876) (0.544,-1.878) (0.579,-1.880) (0.614,-1.882) (0.649,-1.885) (0.684,-1.887) (0.719,-1.889) (0.754,-1.891) (0.789,-1.893) (0.824,-1.896) (0.859,-1.898) (0.895,-1.900) (0.930,-1.902) (0.965,-1.905) (1.000,-1.907) (1.035,-1.909) (1.070,-1.911) (1.105,-1.913) (1.140,-1.915) (1.175,-1.917) (1.210,-1.920) (1.245,-1.922) (1.280,-1.924) (1.315,-1.926) (1.350,-1.927) (1.385,-1.929) (1.420,-1.931) (1.455,-1.933) (1.491,-1.935) (1.526,-1.936) (1.561,-1.938) (1.596,-1.940) (1.631,-1.941) (1.666,-1.943) (1.701,-1.944) (1.736,-1.946) (1.771,-1.947) (1.806,-1.949) (1.841,-1.950) (1.877,-1.951) (1.912,-1.953) (1.947,-1.954) (1.982,-1.955) (2.017,-1.956) (2.052,-1.958) (2.087,-1.959) (2.122,-1.960) (2.157,-1.961) (2.192,-1.962) (2.228,-1.963) (2.263,-1.964) (2.298,-1.965) (2.333,-1.966) (2.368,-1.967) (2.403,-1.967) (2.438,-1.968) (2.473,-1.969) (2.508,-1.970) (2.544,-1.971) (2.579,-1.971) (2.614,-1.972) (2.649,-1.973) (2.684,-1.973) (2.719,-1.974) (2.754,-1.975) (2.789,-1.975) (2.824,-1.976) (2.860,-1.977) (2.895,-1.977) (2.930,-1.978) (2.965,-1.978) (3.000,-1.979)};
\addplot[lineaD, flechaen=0.142, flechaen=0.858] coordinates {(-3.000,-2.179) (-2.965,-2.178) (-2.930,-2.178) (-2.895,-2.177) (-2.860,-2.177) (-2.824,-2.176) (-2.789,-2.176) (-2.754,-2.175) (-2.719,-2.174) (-2.684,-2.174) (-2.649,-2.173) (-2.614,-2.173) (-2.579,-2.172) (-2.544,-2.171) (-2.509,-2.170) (-2.473,-2.170) (-2.438,-2.169) (-2.403,-2.168) (-2.368,-2.167) (-2.333,-2.167) (-2.298,-2.166) (-2.263,-2.165) (-2.228,-2.164) (-2.193,-2.163) (-2.158,-2.162) (-2.123,-2.161) (-2.087,-2.161) (-2.052,-2.160) (-2.017,-2.159) (-1.982,-2.158) (-1.947,-2.156) (-1.912,-2.155) (-1.877,-2.154) (-1.842,-2.153) (-1.807,-2.152) (-1.772,-2.151) (-1.737,-2.150) (-1.701,-2.148) (-1.666,-2.147) (-1.631,-2.146) (-1.596,-2.145) (-1.561,-2.143) (-1.526,-2.142) (-1.491,-2.141) (-1.456,-2.139) (-1.421,-2.138) (-1.386,-2.136) (-1.351,-2.135) (-1.316,-2.133) (-1.281,-2.132) (-1.245,-2.130) (-1.210,-2.129) (-1.175,-2.127) (-1.140,-2.126) (-1.105,-2.124) (-1.070,-2.122) (-1.035,-2.121) (-1.000,-2.119) (-0.965,-2.118) (-0.930,-2.116) (-0.895,-2.114) (-0.860,-2.113) (-0.825,-2.111) (-0.790,-2.110) (-0.754,-2.108) (-0.719,-2.107) (-0.684,-2.105) (-0.649,-2.104) (-0.614,-2.102) (-0.579,-2.101) (-0.544,-2.100) (-0.509,-2.098) (-0.474,-2.097) (-0.439,-2.096) (-0.404,-2.095) (-0.369,-2.094) (-0.333,-2.093) (-0.298,-2.092) (-0.263,-2.091) (-0.228,-2.090) (-0.193,-2.090) (-0.158,-2.089) (-0.123,-2.089) (-0.088,-2.089) (-0.053,-2.088) (-0.018,-2.088) (0.018,-2.088) (0.053,-2.088) (0.088,-2.089) (0.123,-2.089) (0.158,-2.089) (0.193,-2.090) (0.228,-2.090) (0.263,-2.091) (0.298,-2.092) (0.333,-2.093) (0.369,-2.094) (0.404,-2.095) (0.439,-2.096) (0.474,-2.097) (0.509,-2.098) (0.544,-2.100) (0.579,-2.101) (0.614,-2.102) (0.649,-2.104) (0.684,-2.105) (0.719,-2.107) (0.754,-2.108) (0.790,-2.110) (0.825,-2.111) (0.860,-2.113) (0.895,-2.114) (0.930,-2.116) (0.965,-2.118) (1.000,-2.119) (1.035,-2.121) (1.070,-2.122) (1.105,-2.124) (1.140,-2.126) (1.175,-2.127) (1.210,-2.129) (1.245,-2.130) (1.281,-2.132) (1.316,-2.133) (1.351,-2.135) (1.386,-2.136) (1.421,-2.138) (1.456,-2.139) (1.491,-2.141) (1.526,-2.142) (1.561,-2.143) (1.596,-2.145) (1.631,-2.146) (1.666,-2.147) (1.701,-2.148) (1.737,-2.150) (1.772,-2.151) (1.807,-2.152) (1.842,-2.153) (1.877,-2.154) (1.912,-2.155) (1.947,-2.156) (1.982,-2.158) (2.017,-2.159) (2.052,-2.160) (2.087,-2.161) (2.123,-2.161) (2.158,-2.162) (2.193,-2.163) (2.228,-2.164) (2.263,-2.165) (2.298,-2.166) (2.333,-2.167) (2.368,-2.167) (2.403,-2.168) (2.438,-2.169) (2.473,-2.170) (2.509,-2.170) (2.544,-2.171) (2.579,-2.172) (2.614,-2.173) (2.649,-2.173) (2.684,-2.174) (2.719,-2.174) (2.754,-2.175) (2.789,-2.176) (2.824,-2.176) (2.860,-2.177) (2.895,-2.177) (2.930,-2.178) (2.965,-2.178) (3.000,-2.179)};
\addplot[lineaD, flechaen=0.142, flechaen=0.858] coordinates {(-3.000,0.000) (3.000,0.000)};

  \draw[texto, line width=0.6pt] (axis cs:0,0) circle[radius=1];
\end{axis}
\end{tikzpicture}
\caption{Líneas de $\vD$}
\end{subfigure}
\caption{Esfera dieléctrica ($\er=4$, sombreada) en un campo uniforme $\vE_0$ dirigido de izquierda a
derecha; trazos: equipotenciales. Solución exacta de la \cref{ej:esfera-dielectrica}. (a) El campo
interior es uniforme y más débil, $\vE_{\mathrm{int}}=3\vE_0/(\er+2)$: parte de las líneas de $\vE$ termina en la carga
ligada negativa del hemisferio izquierdo y renace en la positiva del derecho. (b) Las líneas de $\vD$, en
cambio, son continuas (no hay carga libre) y se concentran dentro de la esfera, donde
$D_{\mathrm{int}}=3\er\eo E_0/(\er+2)>\eo E_0$. Ambas familias se refractan en la superficie según la \cref{fig:refraccion}.}
\label{fig:esfera-dielectrica}
\end{figure}

